% Leçon 34 — Expression écrite : économie du Cameroun (为什么喀麦隆经济比泰发展？) — Chinois 1ère A — Cours signé OnBuch+
% Programme officiel MINESEC. Compiler avec : tectonic ZHA34-ecrit-economie-du-cameroun.tex
\newcommand{\DOCMATIERE}{Chinois}
\newcommand{\DOCNIVEAU}{1ère A}
\newcommand{\DOCMODULE}{Module 4 : Activités économiques et monde du travail}
\newcommand{\DOCLECON}{ZHA34}
\newcommand{\DOCTITRE}{Leçon 34 — Expression écrite : économie du Cameroun (为什么喀麦隆经济比泰发展？)}
% OnBuch+ — préambule « cours signés » ENRICHI (compilé avec tectonic / XeLaTeX)
% Système visuel « Pop » d'OnBuch+ : fond crème (jamais blanc), bordures encre
% épaisses, ombres « dures » décalées, titres Archivo Black en capitales.
% Version MATHÉMATIQUES : graphes (pgfplots/tikz), boîtes pédagogiques
% (définition, propriété, méthode, attention, le savais-tu, exercice résolu,
% exercice, TP, bilan), page de garde et signature OnBuch+ ancrée en bas.
%
% Macros attendues dans main.tex AVANT \input{preamble} :
%   \DOCMATIERE  \DOCNIVEAU  \DOCMODULE  \DOCLECON (numéro)  \DOCTITRE
\documentclass[11pt]{article}

\usepackage{fontspec}
\usepackage[french]{babel} % français = langue principale ; le chinois via \zh{..} / textebox (xeCJK)
\usepackage{xeCJK}
\usepackage{pifont} % \ding{51} \ding{55} utilisés par les modèles
\usepackage[a4paper,margin=2cm,top=3.4cm,bottom=2.8cm,headheight=1.4cm,footskip=1.3cm]{geometry}
\usepackage{amsmath,amssymb,amsfonts}
\usepackage{mathtools} % \xmapsto, \coloneqq... souvent utilisés par les modèles
\usepackage{cancel} % simplifications barrées (bilans de forces)
% \abs{z} (module, valeur absolue) et \norm{u} (norme), utilisés par les modèles
\providecommand{\abs}{}\let\abs\relax\DeclarePairedDelimiter{\abs}{\lvert}{\rvert}
\providecommand{\norm}{}\let\norm\relax\DeclarePairedDelimiter{\norm}{\lVert}{\rVert}
\usepackage[table]{xcolor} % [table] : fournit aussi \rowcolors (alternance de lignes)
\usepackage{pagecolor}
\usepackage{tikz}
\usetikzlibrary{arrows.meta,positioning,calc,shapes.geometric,shapes.misc,shapes.arrows,decorations.pathmorphing,decorations.pathreplacing,decorations.markings,patterns,shadows,mindmap,trees,fit,backgrounds,matrix,intersections,angles,quotes}
\usepackage{pgfplots}
\usepackage[version=4]{mhchem} % équations nucléaires \ce{^{235}_{92}U}
\usepackage[european,siunitx]{circuitikz} % schémas électriques
\pgfplotsset{compat=1.18}
\usepgfplotslibrary{fillbetween,groupplots} % aires entre courbes (calcul intégral)
\usepackage{enumitem}
\usepackage{fancyhdr}
% [most] charge toutes les bibliothèques tcolorbox (listings, minted, raster,
% documentation...) et gonfle inutilement la mémoire TeX sur un document long
% (~300 boîtes) : on ne charge que ce qui est réellement utilisé.
\usepackage[skins,breakable]{tcolorbox}
\usepackage{graphicx}
\usepackage{colortbl}
\usepackage{booktabs}
\usepackage{tabularx}
\usepackage{diagbox} % case d'en-tête diagonale des tableaux à double entrée
% Colonnes extensibles centrée / gauche / droite (C, L, R), souvent utilisées par les modèles
\newcolumntype{C}{>{\centering\arraybackslash}X}
\newcolumntype{L}{>{\raggedright\arraybackslash}X}
\newcolumntype{R}{>{\raggedleft\arraybackslash}X}
\usepackage{array}
\usepackage{multicol}
\usepackage{multirow}
\usepackage{siunitx}
% parse-numbers=false : accepte aussi \SI{2,5 \times 10^{-3}}{...}
\sisetup{locale=FR,output-decimal-marker={,},group-minimum-digits=4,parse-numbers=false}
\DeclareSIUnit\ohm{\ensuremath{\symup{\Omega}}} % U+2126 absent de la police
\DeclareSIUnit\division{div} % réglages d'oscilloscope (ms/div, V/div)
\DeclareSIUnit\tr{tr} % tours (vitesse de rotation, ex. \SI{1500}{\tr\per\minute})
\DeclareSIUnit\molar{M} % concentration molaire (ex. \SI{3}{\molar}, \SI{1}{\micro\molar})
\DeclareSIUnit\an{an} % année (le modèle confond parfois avec \year, un compteur TeX) — ex. \SI{5}{\kilo\meter\per\an}
\DeclareSIUnit\FCFA{FCFA} % franc CFA (ex. \SI{500}{\FCFA\per\kilo\gram})
\DeclareSIUnit\degreeBrix{\textdegree Brix} % teneur en sucre d'un sirop/jus (réfractométrie)
\DeclareSIUnit\month{mois} % le modèle utilise parfois \month, un compteur TeX, comme unité de durée
\usepackage{qrcode}
% \degree : commande fréquemment utilisée par les modèles pour un angle ou une
% température en dehors de \SI{}{} (non fournie par les paquets chargés).
\providecommand{\degree}{^{\circ}}
% \qed (fin de démonstration), utilisé par les modèles sans amsthm
\providecommand{\qed}{\hfill$\square$}
% \barre : double barre des valeurs interdites dans un tableau de variations (tkz-tab)
\providecommand{\barre}{\ensuremath{\|}}
% environnement proof (amsthm non chargé) utilisé par les modèles
\ifdefined\proof\else\newenvironment{proof}[1][Démonstration]{\par\noindent\textit{#1.}\ }{\qed\par}\fi
\usepackage{lastpage}
\usepackage{hyperref}
\hypersetup{hidelinks}

% ============================================================
% Palette « Pop » OnBuch+ — mêmes hex que la classe P (plus_ui.dart)
% ============================================================
\definecolor{poporange}{HTML}{F59321}
\definecolor{popdark}{HTML}{E07A0C}
\definecolor{popcream}{HTML}{FFF6E8}
\definecolor{popink}{HTML}{1C1714}
\definecolor{poppurple}{HTML}{7A5AE0}
\definecolor{popgreen}{HTML}{2E9E6A}
\definecolor{poppink}{HTML}{E0457B}
\definecolor{popblue}{HTML}{2F7DE1}
\definecolor{popgold}{HTML}{C98A12}
\definecolor{popmuted}{HTML}{8A7F76}
\definecolor{popline}{HTML}{EADFD2}
\definecolor{popgray}{HTML}{8A7F76}
\colorlet{popgrey}{popgray}
% Alias tolérants (le modèle nomme parfois les couleurs différemment)
\colorlet{popwhite}{white}
\colorlet{popblack}{popink}
\colorlet{popred}{poppink}
\colorlet{popyellow}{popgold}
% Teintes claires (fonds de tableaux/encadrés) — le modèle nomme parfois
% les couleurs avec un suffixe L (ex. popblueL) sans les avoir définies.
\colorlet{poporangeL}{poporange!20}
\colorlet{popdarkL}{popdark!20}
\colorlet{poppurpleL}{poppurple!20}
\colorlet{popgreenL}{popgreen!20}
\colorlet{poppinkL}{poppink!20}
\colorlet{popblueL}{popblue!20}
\colorlet{popgoldL}{popgold!20}
\colorlet{popredL}{popred!20}
\colorlet{popgrayL}{popgray!20}
% Teintes claires pour fonds de tableaux / zones
\colorlet{popblueL}{popblue!10!white}
\colorlet{poporangeL}{poporange!14!white}
\colorlet{popgreenL}{popgreen!12!white}
\colorlet{poppurpleL}{poppurple!12!white}
\colorlet{poppinkL}{poppink!12!white}
\colorlet{popgoldL}{popgold!14!white}

\pagecolor{popcream}

% ============================================================
% Typographie
% ============================================================
\defaultfontfeatures{Ligatures=TeX}
\setmainfont{PlusJakartaSans-Medium.ttf}[Path=../../../fonts/,BoldFont=PlusJakartaSans-Bold.ttf]
% Symboles, API et emojis absents de la police principale : repli sur DejaVu Sans (caractères actifs)
\newfontface\symfont{DejaVuSans.ttf}[Path=../../../fonts/]
\makeatletter
\newcommand\symactive[1]{\catcode#1=13 \begingroup\lccode`\~=#1 \lowercase{\endgroup\def~}{{\symfont\char#1}}}
\newcommand\symblank[1]{\catcode#1=13 \begingroup\lccode`\~=#1 \lowercase{\endgroup\def~}{}}
\newcommand\symhyph[1]{\catcode#1=13 \begingroup\lccode`\~=#1 \lowercase{\endgroup\def~}{\mbox{-}}}
\newcount\symc
\newcommand\symrange[2]{\symc=#1 \@whilenum\symc<#2 \do{\expandafter\symactive\expandafter{\the\symc}\advance\symc by 1 }}
\newcommand\symblankrange[2]{\symc=#1 \@whilenum\symc<#2 \do{\expandafter\symblank\expandafter{\the\symc}\advance\symc by 1 }}
\makeatother
\symrange{"0250}{"02FF}\symactive{"2713}\symactive{"2714}\symactive{"2717}\symactive{"2718}\symactive{"2610}\symactive{"2611}\symactive{"2612}
\symrange{"2600}{"27BF}
\catcode"2705=13 \begingroup\lccode`\~="2705 \lowercase{\endgroup\def~}{{\symfont\char"2713}}\symblank{"26BD}\symblank{"26A1}
\symrange{"2460}{"257F}\symactive{"223C}\symactive{"2248}\symactive{"2260}
\symactive{"01F8}\symactive{"01F9}\symactive{"1E3E}\symactive{"1E3F}
\symhyph{"2011}\symhyph{"2E1A}\symblankrange{"1F300}{"1FAFF}\symblank{"FE0F}
% Caractères chinois : WenQuanYi Zen Hei (sous-ensemble GB2312 + pinyin), gras/italique simulés
\setCJKmainfont{WenQuanYiZenHei-subset.ttf}[Path=../../../fonts/,AutoFakeBold=3,AutoFakeSlant=0.2]
\xeCJKsetup{CJKspace=false}
% Pinyin : voyelles à caron (ǎ ǐ ǒ ǔ ǖ ǘ ǚ ǜ) absentes de la police principale -> caractères actifs
\newfontface\pyfont{WenQuanYiZenHei-subset.ttf}[Path=../../../fonts/]
\newcommand\pyactive[1]{\catcode#1=13 \begingroup\lccode`\~=#1 \lowercase{\endgroup\def~}{{\pyfont\char#1}}}
\pyactive{"01CD}\pyactive{"01CE}\pyactive{"01CF}\pyactive{"01D0}\pyactive{"01D1}\pyactive{"01D2}\pyactive{"01D3}\pyactive{"01D4}\pyactive{"01D5}\pyactive{"01D6}\pyactive{"01D7}\pyactive{"01D8}\pyactive{"01D9}\pyactive{"01DA}\pyactive{"01DB}\pyactive{"01DC}
% Maths en sans-serif (Fira Math) avec chiffres et lettres droites de Plus
% Jakarta, pour que les formules se fondent dans le texte.
\usepackage[math-style=ISO,bold-style=ISO]{unicode-math}
\setmathfont{FiraMath-Regular.otf}[Path=../../../fonts/]
\setmathfont{PlusJakartaSans-Medium.ttf}[Path=../../../fonts/,range={up/num,up/{Latin,latin},\mathup}]
\usepackage{icomma} % virgule décimale sans espace en mode math
% Caractères mathématiques Unicode absents des polices de texte (Plus Jakarta,
% Archivo Black) : rendus par leur équivalent mathématique, en texte comme en
% formule — sinon ils s'affichent en carré vide (ex. « Arithmétique dans ℤ »).
\usepackage{newunicodechar}
\newunicodechar{ℝ}{\ensuremath{\mathbb{R}}}
\newunicodechar{ℤ}{\ensuremath{\mathbb{Z}}}
\newunicodechar{ℂ}{\ensuremath{\mathbb{C}}}
\newunicodechar{ℕ}{\ensuremath{\mathbb{N}}}
\newunicodechar{ℚ}{\ensuremath{\mathbb{Q}}}
\newunicodechar{𝔻}{\ensuremath{\mathbb{D}}}
\newunicodechar{α}{\ensuremath{\alpha}}
\newunicodechar{β}{\ensuremath{\beta}}
\newunicodechar{γ}{\ensuremath{\gamma}}
\newunicodechar{δ}{\ensuremath{\delta}}
\newunicodechar{ε}{\ensuremath{\varepsilon}}
\newunicodechar{θ}{\ensuremath{\theta}}
\newunicodechar{λ}{\ensuremath{\lambda}}
\newunicodechar{μ}{\ensuremath{\mu}}
\newunicodechar{σ}{\ensuremath{\sigma}}
\newunicodechar{τ}{\ensuremath{\tau}}
\newunicodechar{φ}{\ensuremath{\varphi}}
\newunicodechar{ψ}{\ensuremath{\psi}}
\newunicodechar{ω}{\ensuremath{\omega}}
\newunicodechar{Σ}{\ensuremath{\Sigma}}
% majuscules produites par \MakeUppercase dans les titres (α-aminés -> Α)
\newunicodechar{Α}{\ensuremath{\alpha}}
\newunicodechar{Β}{\ensuremath{\beta}}
\newunicodechar{Γ}{\ensuremath{\Gamma}}
\newunicodechar{Δ}{\ensuremath{\Delta}}
\newunicodechar{Ε}{\ensuremath{\varepsilon}}
\newunicodechar{Θ}{\ensuremath{\Theta}}
\newunicodechar{Λ}{\ensuremath{\Lambda}}
\newunicodechar{Μ}{\ensuremath{\mu}}
\newunicodechar{Τ}{\ensuremath{\tau}}
\newunicodechar{Φ}{\ensuremath{\Phi}}
\newunicodechar{Ψ}{\ensuremath{\Psi}}
\newunicodechar{Ω}{\ensuremath{\Omega}}
\newunicodechar{↦}{\ensuremath{\mapsto}}
\newunicodechar{∈}{\ensuremath{\in}}
\newunicodechar{∉}{\ensuremath{\notin}}
\newunicodechar{∧}{\ensuremath{\wedge}}
\newunicodechar{⇔}{\ensuremath{\Leftrightarrow}}
\newunicodechar{⟺}{\ensuremath{\Longleftrightarrow}}
\newunicodechar{⇒}{\ensuremath{\Rightarrow}}
\newunicodechar{⟹}{\ensuremath{\Longrightarrow}}
\newunicodechar{∘}{\ensuremath{\circ}}
\newunicodechar{∪}{\ensuremath{\cup}}
\newunicodechar{∩}{\ensuremath{\cap}}
\newunicodechar{≡}{\ensuremath{\equiv}}
\newunicodechar{⊂}{\ensuremath{\subset}}
\newunicodechar{⊥}{\ensuremath{\perp}}
\newunicodechar{∃}{\ensuremath{\exists}}
\newunicodechar{∀}{\ensuremath{\forall}}
\newunicodechar{∛}{\ensuremath{\sqrt[3]{\,}}}
\newunicodechar{∜}{\ensuremath{\sqrt[4]{\,}}}
\newunicodechar{⃗}{\ensuremath{{}^{\to}}}
\newunicodechar{ⁿ}{\ifmmode^{n}\else\textsuperscript{\ensuremath{n}}\fi}
\newunicodechar{ˣ}{\ifmmode^{x}\else\textsuperscript{\ensuremath{x}}\fi}
\newunicodechar{ᵇ}{\ifmmode^{b}\else\textsuperscript{\ensuremath{b}}\fi}
\newunicodechar{ʳ}{\ifmmode^{r}\else\textsuperscript{\ensuremath{r}}\fi}
\newunicodechar{ᵘ}{\ifmmode^{u}\else\textsuperscript{\ensuremath{u}}\fi}
\newunicodechar{ʸ}{\ifmmode^{y}\else\textsuperscript{\ensuremath{y}}\fi}
\newunicodechar{ᵃ}{\ifmmode^{a}\else\textsuperscript{\ensuremath{a}}\fi}
\newunicodechar{ᵉ}{\ifmmode^{e}\else\textsuperscript{\ensuremath{e}}\fi}
\newunicodechar{ᵏ}{\ifmmode^{k}\else\textsuperscript{\ensuremath{k}}\fi}
\newunicodechar{ᵖ}{\ifmmode^{p}\else\textsuperscript{\ensuremath{p}}\fi}
\newunicodechar{ᴺ}{\ifmmode^{N}\else\textsuperscript{\ensuremath{N}}\fi}
\newunicodechar{ᶜ}{\ifmmode^{c}\else\textsuperscript{\ensuremath{c}}\fi}
\newunicodechar{ʰ}{\ifmmode^{h}\else\textsuperscript{\ensuremath{h}}\fi}
\newunicodechar{ᵠ}{\ifmmode^{q}\else\textsuperscript{\ensuremath{q}}\fi}
\newunicodechar{ₙ}{\ifmmode_{n}\else\textsubscript{\ensuremath{n}}\fi}
\newunicodechar{ₐ}{\ifmmode_{a}\else\textsubscript{\ensuremath{a}}\fi}
\newunicodechar{ᵢ}{\ifmmode_{i}\else\textsubscript{\ensuremath{i}}\fi}
\newunicodechar{ᵣ}{\ifmmode_{r}\else\textsubscript{\ensuremath{r}}\fi}
\newunicodechar{ₖ}{\ifmmode_{k}\else\textsubscript{\ensuremath{k}}\fi}
\newunicodechar{ᵦ}{\ifmmode_{\beta}\else\textsubscript{\ensuremath{\beta}}\fi}
\newunicodechar{ₓ}{\ifmmode_{x}\else\textsubscript{\ensuremath{x}}\fi}
\newfontfamily\popdisplay{ArchivoBlack-Regular.ttf}[Path=../../../fonts/]
\newfontfamily\poplogo{SpaceGrotesk-Bold.ttf}[Path=../../../fonts/]
\newfontfamily\popsemi{PlusJakartaSans-SemiBold.ttf}[Path=../../../fonts/]
\newfontfamily\popxbold{PlusJakartaSans-ExtraBold.ttf}[Path=../../../fonts/]
\newfontfamily\popxboldls{PlusJakartaSans-ExtraBold.ttf}[Path=../../../fonts/,LetterSpace=3.0]

\setlist[enumerate]{leftmargin=1.7em,itemsep=5pt,topsep=4pt}
\setlist[itemize]{leftmargin=1.7em,itemsep=4pt,topsep=4pt,label={\color{poporange}\textbullet}}
\setlength{\parindent}{0pt}
\setlength{\parskip}{5pt}
\renewcommand{\arraystretch}{1.25}

% pgfplots : style Pop par défaut
\pgfplotsset{
  every axis/.append style={axis line style={popink,line width=1pt},tick style={popink},
    grid style={popline,line width=0.5pt},label style={font=\small},tick label style={font=\footnotesize}},
  popcurve/.style={poporange,line width=1.8pt,no markers,smooth},
}
% Même style disponible dans un \draw TikZ simple (hors pgfplots)
\tikzset{popcurve/.style={poporange,line width=1.8pt}}
\tikzset{popgrid/.style={popline,very thin}}
\colorlet{popcurve}{poporange} % le nom du style est parfois utilisé comme couleur

% ============================================================
% Pill / squiggle
% ============================================================
\newcommand{\poppill}[3]{%
  \tikz[baseline=-0.55ex]\node[draw=popink,line width=1.2pt,rounded corners=2.6pt,%
    fill=#2,inner xsep=6.5pt,inner ysep=3pt]{%
    \popxboldls\fontsize{7}{7}\selectfont\color{#3}\MakeUppercase{#1}};%
}
\newcommand{\popsquiggle}[1]{%
  \tikz[baseline=-0.4ex]{
    \draw[#1,line width=2.6pt,line cap=round]
      plot [smooth] coordinates {(0,0.09) (0.34,0.01) (0.68,0.21) (1.02,0.01) (1.36,0.10)};
  }%
}
% Mot-clé surligné (marqueur orange sous le texte)
\newcommand{\cle}[1]{{\popxbold\color{popdark}#1}}

% ============================================================
% En-tête / pied
% ============================================================
\pagestyle{fancy}
\fancyhf{}
\renewcommand{\headrulewidth}{0pt}
\renewcommand{\footrulewidth}{0pt}
\lhead{\raisebox{-0.15ex}{\poplogo\fontsize{13}{13}\selectfont\color{popink}OnBuch\color{poporange}+}}
\rhead{\poppill{\DOCMATIERE\ \textbullet\ \DOCNIVEAU\ \textbullet\ Leçon \DOCLECON}{white}{popink}}
\lfoot{\popxbold\fontsize{7.6}{7.6}\selectfont\color{popmuted}\MakeUppercase{Cours signé OnBuch+}\ \textbullet\ {\popsemi\DOCTITRE}}
\rfoot{\tikz[baseline=-0.6ex]\node[rounded corners=8.5pt,fill=popink,minimum height=17pt,minimum width=17pt,inner xsep=5pt,inner sep=0]{\popxbold\fontsize{8}{8}\selectfont\color{popcream}\thepage\,/\,\pageref*{LastPage}};}

% ============================================================
% Titres
% ============================================================
\newcommand{\chaptitre}[2]{%
  \begin{tcolorbox}[enhanced,colback=poporange,colframe=popink,boxrule=2.6pt,arc=11pt,
      drop shadow=popink,left=15pt,right=15pt,top=12pt,bottom=11pt,before skip=3pt,after skip=18pt]
    \poppill{#2}{popink}{popcream}\par\vspace{9pt}
    {\raggedright\hyphenpenalty=10000\exhyphenpenalty=10000\popdisplay\fontsize{22}{24}\selectfont\color{popcream}\MakeUppercase{#1}\par}
  \end{tcolorbox}
}
% Section numérotée : pastille orange avec le numéro + titre Archivo
\newcounter{popsec}
\newcommand{\coursec}[1]{%
  \stepcounter{popsec}\setcounter{popsub}{0}%
  \par\vspace{16pt}\noindent
  \begin{minipage}{\linewidth}
  \tikz[baseline=-0.7ex]\node[draw=popink,line width=1.6pt,fill=poporange,rounded corners=4pt,minimum size=22pt,inner sep=0]{\popdisplay\fontsize{12}{12}\selectfont\color{popcream}\thepopsec};%
  \hspace{8pt}{\popdisplay\fontsize{15}{17}\selectfont\color{popink}\MakeUppercase{#1}}\par
  \vspace{4pt}\noindent{\color{poporange}\rule{36pt}{3.6pt}}
  \end{minipage}\par\nopagebreak\vspace{8pt}
}
\newcounter{popsub}
\newcommand{\courssub}[1]{%
  \stepcounter{popsub}%
  \par\vspace{9pt}\noindent{\popxbold\fontsize{11.5}{13}\selectfont\color{popink}{\color{poporange}\thepopsec.\thepopsub}\hspace{6pt}#1}\par\nopagebreak\vspace{4pt}
}

% ============================================================
% Boîtes pédagogiques — même recette : fond blanc, bordure épaisse colorée,
% ombre dure encre, badge plein. Titre optionnel [#1].
% ============================================================
\tcbset{popbase/.style={breakable,enhanced,colback=white,boxrule=2.3pt,arc=9pt,
  shadow={3pt}{-3pt}{0pt}{popink},left=14pt,right=14pt,top=11pt,bottom=11pt,before skip=11pt,after skip=15pt,
  fontupper=\popsemi\fontsize{10.6}{14.5}\selectfont\color{popink},
  fontlower=\popsemi\fontsize{10.6}{14.5}\selectfont\color{popink}}}
% Coche dessinée (le glyphe ✓ n'existe pas dans les polices du document)
\newcommand{\popcheck}[1]{\tikz[baseline=-0.5ex]\draw[#1,line width=1.6pt,line cap=round,line join=round](-0.13,0.0)--(-0.04,-0.09)--(0.13,0.1);}
\providecommand{\coche}{\popcheck{popgreen}}
\providecommand{\resultat}[1]{\par\smallskip\noindent\fcolorbox{popgreen}{popgreenL}{\parbox{\dimexpr\linewidth-2\fboxsep-2\fboxrule\relax}{\textbf{Résultat.} #1}}\par\smallskip}
\newcommand{\popboxhead}[3]{\poppill{#1}{#2}{white}\ifx\relax#3\relax\else\hspace{6pt}{\popxbold\fontsize{10.6}{12}\selectfont\color{popink}#3}\fi\par\vspace{7pt}}

\newtcolorbox{aretenir}[1][]{popbase,colframe=popgreen,before upper={\popboxhead{À retenir}{popgreen}{#1}}}
% Texte en chinois (document de lecture, dialogue, extrait) : cadre
% d'étude. Un titre optionnel (source, auteur) s'affiche dans l'en-tête.
\newtcolorbox{textebox}[1][]{popbase,colframe=popblue,before upper={\popboxhead{文本 · Texte}{popblue}{#1}},after upper={}}
% Marques ✓ / ✗ (forme correcte / fautive) souvent utilisées par les modèles
\providecommand{\cross}{\textcolor{poppink}{\ensuremath{\times}}}
\providecommand{\xmark}{\cross}\providecommand{\crossmark}{\cross}\providecommand{\cmark}{\ensuremath{\checkmark}}
% Ligne de réponse / justification à compléter (inventée par les modèles)
\providecommand{\justification}[1]{\par\noindent\textit{Justification~:} \dotfill\par}
% \textipa (package tipa non chargé) : la police ne couvre pas l'API -> simple passage
\providecommand{\textipa}[1]{#1}
% Soulignement (ulem non chargé) : \uline, \ul
\providecommand{\uline}[1]{\underline{#1}}\providecommand{\ul}[1]{\underline{#1}}
% \glossaire{\item ...} inventée par les modèles : liste de glossaire
\providecommand{\glossaire}[1]{\begin{itemize}#1\end{itemize}}
% \alert (beamer) inventée par les modèles : texte mis en évidence
\providecommand{\alert}[1]{\textbf{\textcolor{poppink}{#1}}}
% variantes ASCII d'ordinaux écrites en macro
\providecommand{\iere}{\textsuperscript{re}}\providecommand{\ieme}{\textsuperscript{e}}\providecommand{\ere}{\textsuperscript{re}}\providecommand{\eme}{\textsuperscript{e}}\providecommand{\er}{\textsuperscript{er}}\providecommand{\ie}{\textsuperscript{e}}
% Ordinaux français écrits en macro par les modèles : 1\ière, 3\ième
\newcommand{\ière}{\textsuperscript{re}}\newcommand{\ième}{\textsuperscript{e}}
% \exemple (inventée par les modèles) : étiquette « Exemple : »
\providecommand{\exemple}{\textbf{Exemple~:}\ }
% \square utilisable aussi hors mode math (cases à cocher)
\AtBeginDocument{\let\squareMath\square\renewcommand{\square}{\ensuremath{\squareMath}}}
% Mot ou phrase en chinois à l'intérieur d'un paragraphe français
\newcommand{\zh}[1]{\ifmmode\text{#1}\else{#1}\fi}
\let\eng\zh \let\esp\zh \let\all\zh % alias tolérés
% flèches d'intonation inventées par les modèles
\providecommand{\textsearrow}{}\renewcommand{\textsearrow}{\ensuremath{\searrow}}\providecommand{\textnearrow}{}\renewcommand{\textnearrow}{\ensuremath{\nearrow}}
% \fr{..} : mot français (inventé par les modèles)
\providecommand{\fr}[1]{\textit{#1}}
% zhbox : encadré de texte chinois inventé par les modèles
\newenvironment{zhbox}{\begin{quote}}{\end{quote}}
% \courssubsub : titre de niveau 3 inventé par les modèles
\providecommand{\courssubsub}[1]{\par\medskip\noindent\textbf{#1}\par\smallskip}
% \zhbf : texte chinois en gras inventé par les modèles
\providecommand{\zhbf}[1]{\textbf{#1}}
% \vs : « versus » inventé par les modèles
\providecommand{\vs}{\textit{vs}\ }
% \blank : case à compléter inventée par les modèles
\providecommand{\blank}[1][]{\underline{\hspace{1.6em}}}
\newtcolorbox{exemplebox}[1][]{popbase,colframe=poporange,before upper={\popboxhead{Exemple}{poporange}{#1}}}
\newtcolorbox{definition}[1][]{popbase,colframe=popblue,before upper={\popboxhead{Définition}{popblue}{#1}}}
\newtcolorbox{propriete}[1][]{popbase,colframe=poppurple,before upper={\popboxhead{Propriété}{poppurple}{#1}}}
\newtcolorbox{methode}[1][]{popbase,colframe=popgold,before upper={\popboxhead{Méthode}{popgold}{#1}}}
\newtcolorbox{attention}[1][]{popbase,colframe=poppink,before upper={\popboxhead{Attention}{poppink}{#1}}}
\newtcolorbox{experience}[1][]{popbase,colframe=popblue,colback=popblueL,before upper={\popboxhead{Expérience}{popblue}{#1}}}
% « Le savais-tu ? » : carte violette pleine (contexte, culture, Cameroun)
\newtcolorbox{savaistu}[1][]{popbase,colback=poppurple,colframe=popink,
  fontupper=\popsemi\fontsize{10.4}{14.2}\selectfont\color{white},
  before upper={\poppill{Le savais-tu\,?}{white}{poppurple}\ifx\relax#1\relax\else\hspace{6pt}{\popxbold\color{white}#1}\fi\par\vspace{7pt}}}
% Exercice résolu : énoncé en haut, \tcblower puis la solution sur fond crème
\newcounter{popexo}\newcounter{popexr}
\newtcolorbox{exoresolu}[1][]{popbase,colframe=poporange,colbacklower=popcream,
  segmentation style={popink,line width=1.2pt,dashed},
  before upper={\stepcounter{popexr}\popboxhead{Exercice résolu \thepopexr}{poporange}{#1}},
  before lower={\poppill{Solution}{popink}{popcream}\par\vspace{6pt}}}
% Exercice (non résolu en place) — la correction est dans la partie Corrigés
\newtcolorbox{exercice}[1][]{popbase,colframe=popink,
  before upper={\stepcounter{popexo}\popboxhead{Exercice \thepopexo}{popink}{#1}}}
% Corrigé
\newtcolorbox{corrige}[1][]{popbase,colframe=popgreen,colback=popgreenL,
  before upper={\popboxhead{Corrigé}{popgreen}{#1}}}
% Objectifs / prérequis / bilan
\newtcolorbox{objectifs}{popbase,colframe=popink,colback=poporangeL,
  before upper={\popboxhead{Objectifs de la leçon}{popink}{}}}
\newtcolorbox{prerequis}{popbase,colframe=popink,colback=popblueL,
  before upper={\popboxhead{Prérequis}{popblue}{}}}
\newtcolorbox{bilan}{popbase,colframe=popink,colback=white,boxrule=2.8pt,
  before upper={\popboxhead{Fiche bilan}{poporange}{L'essentiel à connaître pour l'examen}}}
% Figure encadrée avec légende
\newtcolorbox{popfigure}[1][]{enhanced,breakable=false,colback=white,colframe=popline,boxrule=1.4pt,arc=8pt,
  left=10pt,right=10pt,top=10pt,bottom=8pt,before skip=10pt,after skip=12pt,halign=center,
  lower separated=false,halign lower=center,
  fontlower=\popsemi\fontsize{9}{11.5}\selectfont\color{popmuted},
  #1}
\newcommand{\legende}[1]{\par\vspace{6pt}{\popsemi\fontsize{9}{11.5}\selectfont\color{popmuted}#1}}

% ============================================================
% Signature OnBuch+
% ============================================================
\newcommand{\poplogoinline}[1]{{\poplogo\fontsize{#1}{#1}\selectfont\color{popink}OnBuch\color{poporange}+}}
\newcommand{\signatureonbuch}{%
  \begin{tcolorbox}[enhanced,colback=popink,colframe=popink,boxrule=2.2pt,arc=10pt,
      drop shadow=poporange,left=14pt,right=14pt,top=10pt,bottom=10pt,before skip=0pt,after skip=0pt]
    \tikz[baseline=-0.7ex]\node[circle,fill=popgreen,draw=popcream,line width=1.3pt,minimum size=22pt,inner sep=0]{\popcheck{white}};%
    \hspace{9pt}\begin{minipage}[c]{0.62\linewidth}
      {\popxbold\fontsize{12}{13}\selectfont\color{popcream}Signé\ \textbullet\ {\poplogo OnBuch\color{poporange}+}}\par\vspace{2pt}
      {\raggedright\popsemi\fontsize{8}{10.5}\selectfont\color{popline}\MakeUppercase{Équipe pédagogique OnBuch+}\\\MakeUppercase{Conforme au programme officiel MINESEC}\par}
    \end{minipage}\hfill
    {\poplogo\fontsize{22}{22}\selectfont\color{popcream}OnBuch\color{poporange}+}
  \end{tcolorbox}%
}

% ============================================================
% Page de garde
% #1 titre · #2 matière · #3 niveau · #4 module · #5 n° leçon · #6 durée ·
% #7 description / objectifs (texte ou itemize)
% ============================================================
\newcommand{\pagegarde}[7]{%
  \thispagestyle{empty}
  \newgeometry{a4paper,margin=1.7cm,top=1.6cm,bottom=1.4cm}%
  \begin{tikzpicture}[remember picture,overlay]
    \fill[poporange!18] ([xshift=-3.2cm,yshift=-3.6cm]current page.north east) circle (4.6cm);
    \fill[poppurple!14] ([xshift=2.4cm,yshift=7.6cm]current page.south west) circle (3.8cm);
  \end{tikzpicture}%
  {\poplogo\fontsize{30}{30}\selectfont\color{popink}OnBuch\color{poporange}+}\hfill
  \raisebox{0.6ex}{\poppill{Cours signé \textbullet\ Premium}{popink}{popcream}}\par
  \vspace{1pt}\popsquiggle{poppurple}\par
  \vspace{8pt}
  \begin{tcolorbox}[enhanced,colback=poporange,colframe=popink,boxrule=2.8pt,arc=13pt,
      drop shadow=popink,left=18pt,right=18pt,top=12pt,bottom=13pt,after skip=10pt]
    \poppill{#2 \textbullet\ #3}{popink}{popcream}\hspace{5pt}\poppill{#4}{white}{popink}\par\vspace{8pt}
    {\popdisplay\fontsize{14}{14}\selectfont\color{popink}LEÇON #5}\par\vspace{4pt}
    {\raggedright\hyphenpenalty=10000\exhyphenpenalty=10000\popdisplay\fontsize{25}{27}\selectfont\color{popcream}\MakeUppercase{#1}\par}
    \vspace{8pt}
    {\popxbold\fontsize{9.5}{9.5}\selectfont\color{popink}\MakeUppercase{Durée conseillée : #6}}
  \end{tcolorbox}
  \begin{tcolorbox}[enhanced,colback=white,colframe=popink,boxrule=2.2pt,arc=10pt,
      drop shadow=popink,left=16pt,right=16pt,top=10pt,bottom=10pt,after skip=8pt]
    \poppill{Dans cette leçon}{poppurple}{white}\par\vspace{5pt}
    \popsemi\fontsize{9.3}{12.6}\selectfont\color{popink}#7
  \end{tcolorbox}
  \noindent\begin{minipage}[t]{0.487\linewidth}
    \begin{tcolorbox}[enhanced,colback=popgreen,colframe=popink,boxrule=2.2pt,arc=10pt,
        drop shadow=popink,left=11pt,right=11pt,top=10pt,bottom=10pt,height=3.1cm,valign=center]
      \tcbox[colback=white,colframe=popink,boxrule=1.3pt,arc=5pt,boxsep=0pt,nobeforeafter,
        left=3pt,right=3pt,top=3pt,bottom=3pt,valign=center]{\qrcode[height=1.9cm]{https://play.google.com/store/apps/details?id=com.luvvix.onbuch&hl=fr}}%
      \hspace{8pt}\begin{minipage}[c]{0.5\linewidth}
        {\popdisplay\fontsize{11}{12.5}\selectfont\color{white}\MakeUppercase{Télécharge OnBuch}}\par\vspace{4pt}
        {\raggedright\popsemi\fontsize{8.4}{11}\selectfont\color{white}Gratuit \textbullet\ Google Play\\Tuteur IA, annales, concours, résultats\par}
      \end{minipage}
    \end{tcolorbox}
  \end{minipage}\hfill
  \begin{minipage}[t]{0.487\linewidth}
    \begin{tcolorbox}[enhanced,colback=poppurple,colframe=popink,boxrule=2.2pt,arc=10pt,
        drop shadow=popink,left=11pt,right=11pt,top=10pt,bottom=10pt,height=3.1cm,valign=center]
      \tikz[baseline=-0.7ex]\node[circle,fill=white,draw=popink,line width=1.3pt,minimum size=1.6cm,inner sep=0]{\popdisplay\fontsize{24}{24}\selectfont\color{poppurple}+};%
      \hspace{8pt}\begin{minipage}[c]{0.6\linewidth}
        {\popdisplay\fontsize{11}{12.5}\selectfont\color{white}\MakeUppercase{Passe à OnBuch+}}\par\vspace{4pt}
        {\raggedright\popsemi\fontsize{8.4}{11}\selectfont\color{white}Tous les cours signés, Tuteur IA illimité, fiches PDF — onglet OnBuch+ de l'app\par}
      \end{minipage}
    \end{tcolorbox}
  \end{minipage}
  \vspace{8pt}
  \popcontenu
  \vfill
  \signatureonbuch
  \restoregeometry
  \newpage
}

% Rangée « Ce cours contient » (page de garde)
\newcommand{\poptile}[3]{% #1 couleur #2 gros texte #3 légende
  \begin{tcolorbox}[enhanced,colback=#1,colframe=popink,boxrule=2pt,arc=9pt,shadow={2.5pt}{-2.5pt}{0pt}{popink},
      left=6pt,right=6pt,top=9pt,bottom=9pt,halign=center,valign=center,height=1.75cm,width=\linewidth,nobeforeafter]
    {\popdisplay\fontsize{12.5}{13}\selectfont\color{white}\MakeUppercase{#2}}\par\vspace{3pt}
    {\centering\popsemi\fontsize{7.8}{9.5}\selectfont\color{white}#3\par}
  \end{tcolorbox}}
\newcommand{\popcontenu}{%
  \par\noindent{\popxboldls\fontsize{7.5}{7.5}\selectfont\color{popmuted}CE COURS SIGNÉ CONTIENT}\par\vspace{7pt}
  \noindent\begin{minipage}[t]{0.235\linewidth}\poptile{popblue}{Cours}{complet, illustré, pas à pas}\end{minipage}\hfill
  \begin{minipage}[t]{0.235\linewidth}\poptile{poporange}{Méthodes}{et exercices résolus}\end{minipage}\hfill
  \begin{minipage}[t]{0.235\linewidth}\poptile{poppink}{Exercices}{type examen + corrigés}\end{minipage}\hfill
  \begin{minipage}[t]{0.235\linewidth}\poptile{popgreen}{Fiche bilan}{l'essentiel à retenir}\end{minipage}\par}

% Sommaire
\newtcolorbox{sommaire}{enhanced,colback=white,colframe=popink,boxrule=2.2pt,arc=10pt,
  drop shadow=popink,left=18pt,right=18pt,top=13pt,bottom=13pt,before skip=6pt,after skip=20pt,
  before upper={\poppill{Sommaire}{poppurple}{white}\par\vspace{9pt}\popsemi\fontsize{10.6}{14.5}\selectfont\color{popink}}}

% Signature de fin de cours — poussée tout en bas de la dernière page
\newcommand{\findecours}{%
  \par\vfill\nopagebreak
  \begin{center}\popsquiggle{poporange}\end{center}\vspace{4pt}
  \signatureonbuch
}
\AtBeginDocument{\def\llbracket{\mathopen{[\mkern-3.5mu[}}\def\rrbracket{\mathclose{]\mkern-3.5mu]}}} % intervalles d'entiers (glyphe ⟦ absent de la police)
% Glyphes absents de Fira Math : redéfinis à partir de glyphes disponibles
\AtBeginDocument{%
  \def\setminus{\mathbin{\backslash}}\let\smallsetminus\setminus
  \def\vdots{\mathord{\vbox{\baselineskip3.2pt\lineskiplimit0pt\kern4pt\hbox{.}\hbox{.}\hbox{.}}}}%
  \def\ddots{\mathinner{\mkern1mu\raise6.5pt\hbox{.}\mkern2mu\raise3.6pt\hbox{.}\mkern2mu\raise0.7pt\hbox{.}\mkern1mu}}%
  \def\triangle{\mathord{\scalebox{0.9}{$\symup{\Delta}$}}}%
  \def\nabla{\mathord{\raisebox{\depth}{\scalebox{1}[-1]{$\symup{\Delta}$}}}}%
  \def\checkmark{\popcheck{popdark}}%
  \def\star{\mathbin{\ast}}\def\bigstar{\mathord{\ast}}%
  \def\diamond{\mathbin{\scalebox{0.6}{$\Diamond$}}}%
  \def\bigcup{\mathop{\mathchoice{\raisebox{-0.3em}{\scalebox{1.7}{$\cup$}}}{\scalebox{1.25}{$\cup$}}{\cup}{\cup}}\displaylimits}%
  \def\bigcap{\mathop{\mathchoice{\raisebox{-0.3em}{\scalebox{1.7}{$\cap$}}}{\scalebox{1.25}{$\cap$}}{\cap}{\cap}}\displaylimits}%
  \def\lesseqqgtr{\mathrel{\lessgtr}}\def\bigoplus{\mathop{\oplus}\displaylimits}%
}
\providecommand{\ssi}{si et seulement si} % abréviation fréquente
\AtBeginDocument{\providecommand{\wideparen}{\overparen}} % arc de cercle

%% FIGFIX-BEGIN v3 — correctifs globaux des figures (docs/onbuch-generation/tools/figfix)
\usepackage{adjustbox}\usepackage{etoolbox}
% toute figure TikZ/pgfplots plus large que la ligne est réduite à la largeur disponible
\BeforeBeginEnvironment{tikzpicture}{\begin{adjustbox}{max width=\linewidth}}
\AfterEndEnvironment{tikzpicture}{\end{adjustbox}}
% légendes pgfplots lisibles par-dessus les courbes
\pgfplotsset{every axis/.append style={legend style={fill=white,fill opacity=0.92,text opacity=1,draw=black!15,inner sep=3pt}}}
% étiquettes d'arêtes (syntaxe edge["..."]) avec fond blanc
\tikzset{every edge quotes/.append style={fill=white,inner sep=1.2pt}}
% bulles de cartes mentales : texte à la ligne dans la bulle, bulle un peu plus grande
\tikzset{every concept/.append style={text width=2.0cm,align=center,minimum size=2.3cm,inner sep=2pt}}
%% FIGFIX-END
\begin{document}

\pagegarde{Leçon 34 — Expression écrite : économie du Cameroun (为什么喀麦隆经济比泰发展？)}{Chinois}{1ère A}{Module 4 : Activités économiques et monde du travail}{ZHA34}{2 h}{Cette leçon d'expression écrite apprend à rédiger un texte simple en chinois sur l'économie du Cameroun (agriculture, pétrole, commerce, développement), à écrire correctement les caractères du thème et à traduire des phrases dans les deux sens (thème et version).
\vspace{4pt}
\begin{multicols}{2}\small
\begin{itemize}
\item À la fin de cette leçon, je sais utiliser le vocabulaire clé de l'économie du Cameroun (经济, 发展, 农业, 石油, 出口...)
\item je sais rédiger un texte simple de 5 à 8 phrases sur l'économie du Cameroun
\item je sais structurer mon texte avec les connecteurs 因为...所以..., 而且, 但是, 首先...然后...
\item je sais écrire correctement les caractères du thème en respectant radicaux et ordre des traits
\item je sais distinguer les caractères proches (己/已, 经/轻, 发/友)
\item je sais traduire des phrases simples du français vers le chinois (thème)
\end{itemize}\end{multicols}}

\begin{sommaire}
\begin{multicols}{2}
\begin{enumerate}
\item Modèle de texte commenté
\item Vocabulaire et glossaire du thème
\item Écriture correcte des caractères
\item Structures et connecteurs pour écrire
\item Thème et version guidés
\item Production écrite et évaluation
\item Activité d'intégration
\item Méthodes et exercices résolus
\item Exercices
\item Corrigés des exercices
\item Fiche bilan
\end{enumerate}
\end{multicols}
\end{sommaire}

\begin{objectifs}
\begin{itemize}
\item À la fin de cette leçon, je sais utiliser le vocabulaire clé de l'économie du Cameroun (经济, 发展, 农业, 石油, 出口...)
\item je sais rédiger un texte simple de 5 à 8 phrases sur l'économie du Cameroun
\item je sais structurer mon texte avec les connecteurs 因为...所以..., 而且, 但是, 首先...然后...
\item je sais écrire correctement les caractères du thème en respectant radicaux et ordre des traits
\item je sais distinguer les caractères proches (己/已, 经/轻, 发/友)
\item je sais traduire des phrases simples du français vers le chinois (thème)
\item je sais traduire des phrases simples du chinois vers le français (version)
\item je sais utiliser la structure 比 pour comparer deux économies
\item je sais m'auto-évaluer à l'aide d'une grille d'évaluation
\end{itemize}
\end{objectifs}

\begin{prerequis}
\begin{itemize}
\item Structure de base Sujet + Verbe + Complément (Leçons de Seconde)
\item La comparaison avec 比 (A 比 B + adjectif)
\item Le vocabulaire des pays et des nationalités (中国, 喀麦隆, 泰国)
\item Les nombres et les quantités simples (很多, 一些)
\item Les radicaux de base (讠, 纟, 氵, 木)
\end{itemize}
\end{prerequis}

\newpage

\coursec{Modèle de texte commenté}

\courssub{Situation d'accroche et problématique}

Au lycée de Yaoundé, en classe de Première A, un journaliste chinois en visite pose une question aux élèves :

\zh{为什么喀麦隆经济发展得这么快？}

\emph{Wèishénme Kāmàilóng jīngjì fāzhǎn de zhème kuài ?}

« Pourquoi l'économie du Cameroun se développe-t-elle si vite ? »

Les élèves ont plein d'idées : le pétrole, le cacao, le café, le coton, le port de Douala, l'agriculture\ldots{} Mais comment organiser toutes ces idées dans un petit texte en chinois ? Quels connecteurs utiliser ? Comment écrire correctement les caractères comme \zh{经济} (économie) ou \zh{发展} (développement) ?

\textbf{Problématique :} comment rédiger et traduire un texte simple et bien structuré sur l'économie du Cameroun ?

Dans cette première partie de la leçon, nous allons lire, traduire et analyser un \cle{texte modèle}. Un texte modèle est un exemple à imiter : il nous montre le plan, le vocabulaire et les connecteurs à réutiliser dans notre propre production écrite.

\courssub{Lecture du texte modèle}

\begin{textebox}[喀麦隆的经济 — L'économie du Cameroun]
喀麦隆是非洲的一个国家。\\
\emph{Kāmàilóng shì Fēizhōu de yí ge guójiā.}\\
喀麦隆的经济比泰国发展得慢，但是发展得很快。\\
\emph{Kāmàilóng de jīngjì bǐ Tàiguó fāzhǎn de màn, dànshì fāzhǎn de hěn kuài.}\\
因为喀麦隆有很多石油和农业产品，所以经济越来越好。\\
\emph{Yīnwèi Kāmàilóng yǒu hěn duō shíyóu hé nóngyè chǎnpǐn, suǒyǐ jīngjì yuè lái yuè hǎo.}\\
喀麦隆人种可可、咖啡和棉花。\\
\emph{Kāmàilóng rén zhòng kěkě, kāfēi hé miánhua.}\\
杜阿拉是一个很重要的港口。\\
\emph{Dù'ālā shì yí ge hěn zhòngyào de gǎngkǒu.}\\
我觉得喀麦隆的经济以后会发展得更好。\\
\emph{Wǒ juéde Kāmàilóng de jīngjì yǐhòu huì fāzhǎn de gèng hǎo.}
\end{textebox}

\textbf{Traduction mot à mot (phrase par phrase) :}

\begin{enumerate}
\item \zh{喀麦隆是非洲的一个国家。} — Cameroun / être / Afrique / (的) / un / pays. $\rightarrow$ « Le Cameroun est un pays d'Afrique. »
\item \zh{喀麦隆的经济比泰国发展得慢，但是发展得很快。} — économie du Cameroun / comparé à / Thaïlande / se développer / lentement / mais / se développer / très / vite. $\rightarrow$ « L'économie du Cameroun se développe plus lentement que celle de la Thaïlande, mais elle se développe vite. »
\item \zh{因为喀麦隆有很多石油和农业产品，所以经济越来越好。} — parce que / le Cameroun / avoir / beaucoup de / pétrole / et / produits agricoles / donc / économie / de plus en plus / bonne. $\rightarrow$ « Parce que le Cameroun a beaucoup de pétrole et de produits agricoles, son économie va de mieux en mieux. »
\item \zh{喀麦隆人种可可、咖啡和棉花。} — les Camerounais / cultiver / cacao / café / et / coton. $\rightarrow$ « Les Camerounais cultivent le cacao, le café et le coton. »
\item \zh{杜阿拉是一个很重要的港口。} — Douala / être / un / très / important / port. $\rightarrow$ « Douala est un port très important. »
\item \zh{我觉得喀麦隆的经济以后会发展得更好。} — je / penser-trouver / économie du Cameroun / à l'avenir / (会) / se développer / encore plus / bien. $\rightarrow$ « Je pense que l'économie du Cameroun se développera encore mieux à l'avenir. »
\end{enumerate}

\textbf{Traduction naturelle complète :}

« Le Cameroun est un pays d'Afrique. Son économie se développe plus lentement que celle de la Thaïlande, mais elle se développe vite. Comme le Cameroun possède beaucoup de pétrole et de produits agricoles, son économie va de mieux en mieux. Les Camerounais cultivent le cacao, le café et le coton. Douala est un port très important. Je pense que l'économie du Cameroun se développera encore mieux à l'avenir. »

\begin{definition}[Glossaire du texte modèle]
\begin{center}
\begin{tabularx}{\linewidth}{|X|X|X|X|}
\hline
\rowcolor{popblueL} Caractères & Pinyin & Nature & Traduction \\
\hline
经济 & jīngjì & nom & économie \\
\hline
发展 & fāzhǎn & verbe / nom & (se) développer, développement \\
\hline
比 & bǐ & préposition & comparé à, que (comparaison) \\
\hline
泰国 & Tàiguó & nom propre & Thaïlande \\
\hline
石油 & shíyóu & nom & pétrole \\
\hline
农业 & nóngyè & nom & agriculture \\
\hline
产品 & chǎnpǐn & nom & produit \\
\hline
种 & zhòng & verbe & cultiver, planter \\
\hline
可可 & kěkě & nom & cacao \\
\hline
咖啡 & kāfēi & nom & café \\
\hline
棉花 & miánhua & nom & coton \\
\hline
杜阿拉 & Dù'ālā & nom propre & Douala \\
\hline
港口 & gǎngkǒu & nom & port \\
\hline
重要 & zhòngyào & adjectif & important \\
\hline
越来越 & yuè lái yuè & structure & de plus en plus \\
\hline
觉得 & juéde & verbe & penser, trouver (opinion) \\
\hline
以后 & yǐhòu & nom de temps & après, à l'avenir \\
\hline
更 & gèng & adverbe & encore plus, davantage \\
\hline
\end{tabularx}
\end{center}
\end{definition}

\begin{textebox}[Phrase clé à retenir]
喀麦隆的经济比泰国发展得慢，但是发展得很快。\\
\emph{Kāmàilóng de jīngjì bǐ Tàiguó fāzhǎn de màn, dànshì fāzhǎn de hěn kuài.}\\
« L'économie du Cameroun se développe plus lentement que celle de la Thaïlande, mais elle se développe vite. »
\end{textebox}

\courssub{Analyse du plan du texte}

Un bon petit texte chinois, comme un bon texte français, suit un plan clair. Observons le texte modèle : chaque phrase joue un rôle précis.

\begin{enumerate}
\item \textbf{Présentation (phrase 1) :} on présente le sujet — le Cameroun, pays d'Afrique.
\item \textbf{Comparaison (phrase 2) :} on situe le sujet par rapport à un autre pays avec la structure \zh{比} + connecteur d'opposition \zh{但是}.
\item \textbf{Causes (phrase 3) :} on explique pourquoi avec \zh{因为\ldots 所以\ldots}.
\item \textbf{Exemples (phrases 4 et 5) :} on illustre avec des faits concrets : cacao, café, coton, port de Douala.
\item \textbf{Opinion / conclusion (phrase 6) :} on donne son avis avec \zh{我觉得}.
\end{enumerate}

\begin{popfigure}
\begin{tikzpicture}[node distance=0.55cm,
  etape/.style={rectangle, rounded corners=3pt, draw=popblue, fill=popblueL, text width=4.6cm, align=center, font=\small, minimum height=0.9cm},
  fleche/.style={-{Stealth[length=2.5mm]}, thick, popdark}]
\node[etape, fill=popgoldL, draw=popgold, align=center] (intro){\textbf{1. Présentation}\\ 喀麦隆是非洲的一个国家。};
\node[etape, below=of intro, align=center] (comp){\textbf{2. Comparaison}\\ 比泰国发展得慢，但是很快};
\node[etape, below=of comp, align=center] (cause){\textbf{3. Causes}\\ 因为石油和农业，所以越来越好};
\node[etape, below=of cause, align=center] (ex){\textbf{4. Exemples}\\ 种可可、咖啡、棉花；杜阿拉港口};
\node[etape, below=of ex, fill=popgreenL, draw=popgreen, align=center] (concl){\textbf{5. Opinion}\\ 我觉得以后会发展得更好。};
\draw[fleche] (intro) -- (comp);
\draw[fleche] (comp) -- (cause);
\draw[fleche] (cause) -- (ex);
\draw[fleche] (ex) -- (concl);
\end{tikzpicture}
\legende{Organigramme du plan du texte modèle : cinq étapes à reproduire dans votre production.}
\end{popfigure}

\begin{methode}[Analyser un texte modèle en trois gestes]
Avant d'imiter un texte, il faut l'analyser. Voici la méthode à appliquer sur ta copie ou dans ton cahier :
\begin{enumerate}
\item \textbf{Souligner les connecteurs} (\zh{但是}, \zh{因为\ldots 所以\ldots}, \zh{我觉得}) : ce sont les « charnières » du texte, celles qui relient les idées.
\item \textbf{Encadrer le vocabulaire du thème} (ici : \zh{经济}, \zh{石油}, \zh{农业}, \zh{港口}, \zh{发展}) : ce sont les mots à réutiliser dans ta propre rédaction.
\item \textbf{Numéroter les étapes du plan} dans la marge : présentation (1), comparaison (2), causes (3), exemples (4), opinion (5).
\end{enumerate}
Ensuite seulement, tu peux rédiger ton propre texte en gardant le même squelette et en changeant les informations.
\end{methode}

\courssub{Repérage des connecteurs}

Trois connecteurs structurent le texte modèle. Les voici en action :

\begin{exemplebox}[Les trois connecteurs du texte]
\begin{itemize}
\item \textbf{Opposition :} \zh{但是} (dànshì) = « mais ».\\
\zh{发展得慢，但是发展得很快。} \emph{fāzhǎn de màn, dànshì fāzhǎn de hěn kuài.} « Ça se développe lentement, mais ça se développe vite. »
\item \textbf{Cause et conséquence :} \zh{因为\ldots 所以\ldots} (yīnwèi\ldots suǒyǐ\ldots) = « parce que\ldots donc\ldots ».\\
\zh{因为有很多石油，所以经济越来越好。} \emph{Yīnwèi yǒu hěn duō shíyóu, suǒyǐ jīngjì yuè lái yuè hǎo.} « Parce qu'il y a beaucoup de pétrole, l'économie va de mieux en mieux. »
\item \textbf{Opinion :} \zh{我觉得} (wǒ juéde) = « je pense que, je trouve que ».\\
\zh{我觉得经济以后会发展得更好。} \emph{Wǒ juéde jīngjì yǐhòu huì fāzhǎn de gèng hǎo.} « Je pense que l'économie se développera mieux à l'avenir. »
\end{itemize}
\end{exemplebox}

\begin{propriete}[Règle : 因为\ldots 所以\ldots, les deux ensemble]
En français, on ne dit jamais « parce que\ldots donc\ldots » dans la même phrase : on choisit l'un ou l'autre. En chinois, c'est le contraire : \zh{因为} et \zh{所以} se combinent naturellement dans la même phrase.

\textbf{Structure :} 因为 + cause ，所以 + conséquence 。

\zh{因为喀麦隆有很多石油，所以经济越来越好。} « Parce que le Cameroun a beaucoup de pétrole, son économie va de mieux en mieux. »

\textbf{Exception :} dans la conversation, on peut utiliser \zh{因为} seul ou \zh{所以} seul, mais à l'écrit scolaire, utilise les deux ensemble : c'est la forme attendue et la plus sûre.
\end{propriete}

\begin{attention}[Pièges fréquents des francophones]
\begin{itemize}
\item Ne traduis pas « parce que\ldots donc\ldots » par un seul mot : en chinois, garde \zh{因为} ET \zh{所以}.
\item \zh{得} (de) après le verbe introduit le complément de manière : \zh{发展得很快} = « se développe vite ». N'écris pas \zh{发展的很快} : c'est une erreur classique entre \zh{的} et \zh{得}.
\item \zh{越来越} + adjectif = « de plus en plus\ldots » : \zh{越来越好} (de mieux en mieux). Ne mets pas \zh{很} après \zh{越来越} : on ne dit pas \zh{越来越很好}.
\item L'ordre des mots de la comparaison : A + 比 + B + verbe/adjectif. \zh{喀麦隆比泰国发展得慢} : jamais \zh{喀麦隆发展得比泰国慢} au niveau HSK 2-3 (forme avancée, à éviter pour l'instant).
\item Le nom du port de Douala s'écrit \zh{杜阿拉} (Dù'ālā) en caractères chinois, pas en lettres latines.
\end{itemize}
\end{attention}

\begin{exemplebox}[Exemples traduits à réutiliser]
\zh{喀麦隆有很多石油。} \emph{Kāmàilóng yǒu hěn duō shíyóu.} « Le Cameroun a beaucoup de pétrole. »

\zh{经济越来越好。} \emph{Jīngjì yuè lái yuè hǎo.} « L'économie va de mieux en mieux. »

\zh{喀麦隆人种可可和咖啡。} \emph{Kāmàilóng rén zhòng kěkě hé kāfēi.} « Les Camerounais cultivent le cacao et le café. »

\zh{杜阿拉是一个很重要的港口。} \emph{Dù'ālā shì yí ge hěn zhòngyào de gǎngkǒu.} « Douala est un port très important. »
\end{exemplebox}

\courssub{Carte mentale du vocabulaire du thème}

Pour mémoriser le lexique, organisons-le autour du mot central \zh{经济} (économie) :

\begin{popfigure}
\begin{tikzpicture}[mindmap, grow cyclic,
  every node/.style={concept, font=\small},
  concept color=popblue,
  level 1 concept/.append style={sibling angle=90, level distance=3.2cm, font=\small}]
\node[concept, fill=popblue, text=white, align=center]{经济\\ jīngjì\\ économie}
  child[concept color=popgreen] { node[concept, align=center]{农业\\ nóngyè\\ agriculture\\ 可可 咖啡 棉花} }
  child[concept color=poporange] { node[concept, align=center]{石油\\ shíyóu\\ pétrole} }
  child[concept color=poppurple] { node[concept, align=center]{港口\\ gǎngkǒu\\ port\\ 杜阿拉} }
  child[concept color=popgold] { node[concept, align=center]{发展\\ fāzhǎn\\ développement\\ 越来越好} };
\end{tikzpicture}
\legende{Carte mentale lexicale : quatre branches autour de 经济 (économie).}
\end{popfigure}

\begin{savaistu}[Le port de Douala]
Le port de Douala (杜阿拉港, \emph{Dù'ālā gǎng}) est le plus grand port du Cameroun : une très grande partie du commerce du pays (cacao, café, bois, pétrole) passe par lui. En Chine, le plus grand port du monde est celui de Shanghai (上海港, \emph{Shànghǎi gǎng}) : des millions de conteneurs y transitent chaque année. Les ports sont des symboles du commerce international, en Chine comme au Cameroun !
\end{savaistu}

\begin{exoresolu}[Retrouver le plan d'une phrase]
\textbf{Énoncé.} Voici une phrase du texte modèle : \zh{因为喀麦隆有很多石油和农业产品，所以经济越来越好。} À quelle étape du plan appartient-elle ? Quels connecteurs contient-elle ? Traduisez-la.

\tcblower

\textbf{Solution détaillée.}
\begin{enumerate}
\item \textbf{Étape du plan :} c'est l'étape 3, les \textbf{causes} : la phrase explique \emph{pourquoi} l'économie camerounaise progresse.
\item \textbf{Connecteurs :} la paire \zh{因为\ldots 所以\ldots} (parce que\ldots donc\ldots). On repère aussi la structure \zh{越来越} (de plus en plus).
\item \textbf{Traduction mot à mot :} parce que / Cameroun / avoir / beaucoup de / pétrole / et / produits agricoles / donc / économie / de plus en plus / bonne.
\item \textbf{Traduction naturelle :} « Parce que le Cameroun possède beaucoup de pétrole et de produits agricoles, son économie va de mieux en mieux. »
\end{enumerate}
\end{exoresolu}

\begin{aretenir}[L'essentiel de la section]
\begin{itemize}
\item Un texte simple sur l'économie suit un plan en cinq étapes : \cle{présentation} $\rightarrow$ \cle{comparaison} (\zh{比}) $\rightarrow$ \cle{causes} (\zh{因为\ldots 所以\ldots}) $\rightarrow$ \cle{exemples} $\rightarrow$ \cle{opinion} (\zh{我觉得}).
\item Les connecteurs sont les charnières du texte : \zh{但是} (mais), \zh{因为\ldots 所以\ldots} (parce que\ldots donc), \zh{我觉得} (je pense que).
\item Le vocabulaire clé : \zh{经济} (économie), \zh{发展} (développement), \zh{石油} (pétrole), \zh{农业} (agriculture), \zh{港口} (port), \zh{杜阿拉} (Douala), \zh{越来越好} (de mieux en mieux).
\item Avant de rédiger, analyse toujours le modèle : souligne les connecteurs, encadre le vocabulaire, numérote le plan.
\end{itemize}
\end{aretenir}

\coursec{Vocabulaire et glossaire du thème}

Après avoir analysé le modèle de texte « Pourquoi l'économie du Cameroun se développe-t-elle plus vite que celle du Tchad ? » dans la section précédente, nous allons maintenant maîtriser le lexique indispensable pour comprendre, traduire et rédiger sur ce sujet. Ce glossaire couvre les secteurs moteurs (agriculture, pétrole), les verbes d'action économique et les connecteurs logiques pour structurer votre argumentation.

\courssub{Glossaire thématique : les secteurs clés de l'économie camerounaise}

Observez le tableau ci-dessous qui regroupe le vocabulaire de base du thème. Chaque entrée est présentée avec son caractère, sa prononciation en pinyin (tons marqués), sa nature grammaticale et sa traduction. Mémorisez les associations fréquentes (collocations) indiquées entre parenthèses.

\begin{center}
\begin{tabularx}{\linewidth}{|X|X|X|X|}
\hline
\rowcolor{popblueL} \textbf{Caractères} & \textbf{Pinyin} & \textbf{Nature} & \textbf{Traduction / Collocations} \\ \hline
经济 & jīngjì & n. & économie ; économique (ex: 经济发展 jīngjì fāzhǎn, 经济增长) \\ \hline
发展 & fāzhǎn & v./n. & (se) développer, développement ; développer (ex: 发展经济, 发展农业) \\ \hline
农业 & nóngyè & n. & agriculture (ex: 农业生产, 发展农业) \\ \hline
石油 & shíyóu & n. & pétrole, huile de pétrole (ex: 出口石油, 石油工业) \\ \hline
出口 & chūkǒu & v./n. & exporter, exportation ; port de sortie (ant. 进口 jìnkǒu) \\ \hline
产品 & chǎnpǐn & n. & produit, marchandise (ex: 农产品, 出口产品) \\ \hline
港口 & gǎngkǒu & n. & port maritime, port commercial (ex: 杜阿拉港口) \\ \hline
可可 & kěkě & n. & cacao (ex: 种可可, 可可豆) \\ \hline
咖啡 & kāfēi & n. & café (ex: 咖啡豆, 喝咖啡) \\ \hline
棉花 & miánhua & n. & coton (ex: 棉花地, 出口棉花) \\ \hline
国家 & guójiā & n. & pays, État, nation (ex: 发展中国家) \\ \hline
越来越 & yuè lái yuè & adv. & de plus en plus (structure: 越来越 + adj./v. psych.) \\ \hline
\end{tabularx}
\end{center}

\begin{definition}[Analyse morphologique : 出口 \zh{chūkǒu}]
Le mot \textbf{出口} est un composé verbe-objet typique :
\begin{itemize}
    \item \zh{出} (chū) : verbe « sortir, aller vers l'extérieur » (radical \zh{凵} kǎn, ouverture).
    \item \zh{口} (kǒu) : nom « bouche, ouverture, porte, port » (radical \zh{口} kǒu).
\end{itemize}
Littéralement : « sortir par la porte/le port » $\rightarrow$ \textbf{exporter}.
Son antonyme direct est \textbf{进口 \zh{jìnkǒu}} (importer) : \zh{进} (entrer) + \zh{口} (porte).
\emph{Attention :} \zh{出口} peut être verbe (\zh{出口产品} — exporter des produits) ou nom (\zh{出口量} — volume des exportations).
\end{definition}

\begin{exemplebox}[Exemples contextualisés]
\begin{enumerate}
    \item \zh{喀麦隆出口可可、咖啡和棉花。} (Kāmàilóng chūkǒu kěkě、kāfēi hé miánhua.) — Le Cameroun exporte du cacao, du café et du coton.
    \item \zh{农业很重要，因为国家需要发展经济。} (Nóngyè hěn zhòngyào, yīnwèi guójiā xūyào fāzhǎn jīngjì.) — L'agriculture est très importante, car le pays a besoin de développer l'économie.
    \item \zh{杜阿拉是一个大港口，石油和棉花从这里出口。} (Dùā lā shì yí gè dà gǎngkǒu, shíyóu hé miánhua cóng zhèlǐ chūkǒu.) — Douala est un grand port, le pétrole et le coton y sont exportés.
    \item \zh{这个国家的经济发展得很快。} (Zhè gè guójiā de jīngjì fāzhǎn de hěn kuài.) — L'économie de ce pays se développe très vite.
\end{enumerate}
\end{exemplebox}

\courssub{Collocations et structures lexicales utiles}

Pour écrire un texte cohérent, il ne suffit pas de connaître les mots isolés ; il faut maîtriser leurs \textbf{collocations} (associations figées). Voici les trois structures modèles du thème :

\begin{center}
\begin{tabularx}{\linewidth}{|X|X|X|}
\hline
\rowcolor{popblueL} \textbf{Collocation (Chinois)} & \textbf{Structure / Analyse} & \textbf{Traduction / Emploi} \\ \hline
\zh{发展经济} & Verbe \zh{发展} + Objet \zh{经济} & Développer l'économie (action volontariste). \\ \hline
\zh{出口产品} & Verbe \zh{出口} + Objet \zh{产品} & Exporter des produits (terme de commerce). \\ \hline
\zh{经济发展得很快} & Sujet \zh{经济} + Verbe \zh{发展} + \zh{得} + Adv. \zh{很} + Adj. \zh{快} & L'économie se développe vite (complément de degré \zh{得}). \\ \hline
\end{tabularx}
\end{center}

\begin{propriete}[Structure du complément de degré avec \zh{得}]
Lorsque l'on veut qualifier \textbf{la manière} ou \textbf{le degré} d'une action ou d'un état, on insère \zh{得} (de) après le verbe (ou l'adjectif verbe) suivi de l'adverbe de degré (\zh{很}, \zh{真}, \zh{特别}) et de l'adjectif.
\medskip
\noindent \textbf{Schéma :} Sujet + Verbe + \zh{得} + (Adv.) + Adjectif
\medskip
\noindent \textbf{Exemples :}
\begin{itemize}
    \item \zh{他跑得很快。} (Tā pǎo de hěn kuài.) — Il court très vite.
    \item \zh{经济发展得很好。} (Jīngjì fāzhǎn de hěn hǎo.) — L'économie se développe très bien.
    \item \zh{种得好。} (Zhòng de hǎo.) — Bien cultiver / La culture est bonne.
\end{itemize}
\emph{Erreur fréquente :} Oublier \zh{得} (*\zh{经济发展很快}) ou le confondre avec \zh{的} (particule nominale) ou \zh{地} (adverbial).
\end{propriete}

\courssub{Focus lexical : le verbe \zh{种} (zhòng) « cultiver, planter »}

\begin{definition}[Le verbe \zh{种} (zhòng)]
\zh{种} (zhòng, 4e ton) est un verbe d'action agricole signifiant « planter, cultiver, semer ». Il ne doit pas être confondu avec \zh{种} (zhǒng, 3e ton) qui signifie « espèce, genre, type, graine » (nom / classifieur).
\medskip
\noindent \textbf{Collocations du thème :}
\begin{itemize}
    \item \zh{种可可} (zhòng kěkě) — cultiver le cacao.
    \item \zh{种咖啡} (zhòng kāfēi) — cultiver le café.
    \item \zh{种棉花} (zhòng miánhua) — cultiver le coton.
    \item \zh{种花生} (zhòng huāshēng) — cultiver l'arachide.
\end{itemize}
\noindent \textbf{Structure :} Sujet (paysan, région, pays) + \zh{种} + Objet (culture).
\end{definition}

\begin{exemplebox}[Emploi de \zh{种} (zhòng) vs \zh{种} (zhǒng)]
\begin{itemize}
    \item \zh{喀麦隆农民种可可。} (Kāmàilóng nóngmín zhòng kěkě.) — Les paysans camerounais cultivent le cacao. \emph{(Verbe, 4e ton)}
    \item \zh{这是一种好产品。} (Zhè shì yī zhǒng hǎo chǎnpǐn.) — C'est un bon type de produit. \emph{(Classifieur/nom, 3e ton)}
    \item \zh{什么种子？} (Shénme zhǒngzi?) — Quelle graine / semence ? \emph{(Nom composé, 3e ton)}
\end{itemize}
\end{exemplebox}

\courssub{Distinction fondamentale : \zh{经济} (jīngjì) vs \zh{发展} (fāzhǎn)}

C'est une confusion majeure chez les francophones car « économie » et « développement » sont proches sémantiquement, mais leur nature syntaxique diffère.

\begin{propriete}[Distinction \zh{经济} vs \zh{发展}]
\begin{center}
\begin{tabularx}{\linewidth}{|X|X|X|}
\hline
\rowcolor{popblueL} \textbf{Critère} & \textbf{\zh{经济} (jīngjì)} & \textbf{\zh{发展} (fāzhǎn)} \\ \hline
\textbf{Nature} & Nom (souvent attributif = adjectif) & Verbe \textbf{et} Nom (verbe d'action / nom d'état) \\ \hline
\textbf{Sens principal} & « Économie » (système), « Économique » (adj.) & « Développer » (v.), « Développement » (n.) \\ \hline
\textbf{Sujet possible} & \zh{经济增长} (croissance éco.), \zh{经济中心} & \zh{国家发展}, \zh{发展迅速} \\ \hline
\textbf{Objet possible} & — & \zh{发展经济}, \zh{发展农业}, \zh{发展科技} \\ \hline
\textbf{Modifié par} & \zh{发展的经济} (l'économie qui se développe) & \zh{快速发展} (développement rapide) \\ \hline
\end{tabularx}
\end{center}

\medskip
\noindent \textbf{Règle d'or :} On dit \zh{发展经济} (verbe + objet), jamais *\zh{经济发展} (sauf comme sujet : \zh{经济发展很快} = « Le développement économique est rapide » ou « L'économie se développe vite »).
\end{propriete}

\begin{attention}[Piège de traduction : « Développement économique »]
En français, « développement économique » est un nom composé. En chinois, l'ordre des mots change selon la fonction :
\begin{itemize}
    \item Sujet/Thème : \zh{经济发展} (jīngjì fāzhǎn) — Nom + Verbe nominalisé = « Le développement économique ».
    \item Action : \zh{发展经济} (fāzhǎn jīngjì) — Verbe + Objet = « Développer l'économie ».
    \item Adjectif : \zh{经济的} (jīngjì de) — ex: \zh{经济的发展} (le développement économique [formel]).
\end{itemize}
\emph{Exercice mental :} Traduisez « Le Cameroun veut développer son économie. » $\rightarrow$ \zh{喀麦隆想发展经济。} (Pas *\zh{喀麦隆想经济发展。})
\end{attention}

\courssub{Point de grammaire : l'adverbe \zh{越来越} (yuè lái yuè) « de plus en plus »}

Cet adverbe de degré est essentiel pour décrire les tendances économiques (croissance, inflation, production).

\begin{propriete}[Structure et emploi de \zh{越来越}]
\zh{越来越} se place \textbf{devant un adjectif} ou un \textbf{verbe d'état psychologique} (aimer, détester, vouloir) pour indiquer une progression continue dans le temps.
\medskip
\noindent \textbf{Structure :} Sujet + \zh{越来越} + Adjectif / Verbe psychologique
\medskip
\noindent \textbf{Exemples :}
\begin{itemize}
    \item \zh{喀麦隆的经济越来越好。} (Kāmàilóng de jīngjì yuè lái yuè hǎo.) — L'économie du Cameroun va de mieux en mieux.
    \item \zh{出口产品越来越多。} (Chūkǒu chǎnpǐn yuè lái yuè duō.) — Les produits exportés sont de plus en plus nombreux.
    \item \zh{我越来越喜欢喝咖啡。} (Wǒ yuè lái yuè xǐhuan hē kāfēi.) — J'aime de plus en plus boire du café.
\end{itemize}
\end{propriete}

\begin{attention}[Erreur interdite : \zh{越来越} + \zh{很}]
\cle{Jamais} \zh{越来越很好} \cle{!}
En français, on dit « de plus en plus \textbf{très} bon », mais en chinois, \zh{越来越} \textbf{contient déjà l'idée de degré élevé progressif}. L'adverbe \zh{很} (très) est redondant et grammaticalement incorrect ici.
\medskip
\noindent \textbf{Correct :} \zh{越来越好} (de mieux en mieux), \zh{越来越快} (de plus en plus vite).
\noindent \textbf{Incorrect :} *\zh{越来越很好}, *\zh{越来越很快}.
\end{attention}

\coursec{Écriture correcte des caractères du thème}

Cette section vise l'acquisition de l'ordre des traits, l'identification des radicaux (clés) et la distinction des caractères proches pour le vocabulaire économique de cette leçon. Maîtriser l'écriture renforce la mémorisation du vocabulaire et évite les confusions à l'examen.

\courssub{Ordre des traits, radicaux et pièges visuels}

\begin{center}
\begin{tabularx}{\linewidth}{|X|X|X|X|X|}
\hline
\rowcolor{popblueL} \textbf{Caractère} & \textbf{Pinyin} & \textbf{Radical (Clé)} & \textbf{Nb traits} & \textbf{Ordre des traits (Règles clés) / Caractères proches} \\ \hline
经 & jīng & \zh{纟} (sī, soie) & 8 & Haut→bas, gauche→droite. \zh{纟} (3 traits) + \zh{工} (3) + \zh{巳} (2). \textbf{Proche :} \zh{径} (jìng, diamètre/chemin, radical \zh{彳}). \\ \hline
济 & jì & \zh{氵} (shuǐ, eau) & 13 & Gauche→droite. \zh{氵} (3) + \zh{齐} (qí, 6 traits : haut→bas) + \zh{攵} (pū, 4 traits : 点横竖捺). \textbf{Proche :} \zh{济} vs \zh{制} (zhì, radical \zh{刂}). \\ \hline
发 & fā & \zh{又} (yòu) & 5 & Haut→bas : \zh{癶} (bō, 5 traits) + \zh{又}. \textbf{Attention :} Ne pas confondre avec \zh{友} (yǒu, ami, radical \zh{又} mais structure diff.). \\ \hline
展 & zhǎn & \zh{尸} (shī) & 10 & Haut→bas, gauche→droite. \zh{尸} (3) + \zh{单} (dān, 7). \textbf{Proche :} \zh{尊} (zūn, respecter, haut diff.). \\ \hline
农 & nóng & \zh{曲} (qǔ) & 7 & Enclos→intérieur→fermeture. \zh{曲} (6) + \zh{寸} (cùn, 3) en bas. \textbf{Traditionnel :} \zh{農} (radical \zh{辰}). \\ \hline
业 & yè & \zh{一} (yī) & 5 & Horizontal d'abord. \zh{一} + \zh{丨} + \zh{一} + \zh{丿} + \zh{乀}. \textbf{Proche :} \zh{叶} (yè, feuille, radical \zh{艹}). \\ \hline
石 & shí & \zh{石} (shí) & 5 & Haut→bas : \zh{厂} (hǎn, 2) + \zh{口} (kǒu, 3). \textbf{Proche :} \zh{右} (yòu, droite, \zh{口} en haut à gauche). \\ \hline
油 & yóu & \zh{氵} (shuǐ) & 8 & Gauche→droite. \zh{氵} + \zh{由} (yóu). \textbf{Proche :} \zh{游} (yóu, nager, radical \zh{氵} + \zh{斿}). \\ \hline
出 & chū & \zh{凵} (kǎn) & 5 & Enclos→intérieur. \zh{凵} (2) + \zh{山} (shān, 3). \textbf{Proche :} \zh{凶} (xiōng, funeste, trait vertical traverse). \\ \hline
口 & kǒu & \zh{口} (kǒu) & 3 & Vertical gauche, haut+horizontal+bas (un trait), vertical droit. \textbf{Ne pas écrire} □ (carré fermé en 4 traits). \\ \hline
港 & gǎng & \zh{氵} (shuǐ) & 12 & Gauche→droite. \zh{氵} + \zh{共} (gòng) + \zh{页} (yè). \textbf{Proche :} \zh{江} (jiāng, fleuve, \zh{工} au lieu de \zh{共}). \\ \hline
可 & kě & \zh{口} (kǒu) & 5 & \zh{口} + \zh{丁} (dīng). \textbf{Proche :} \zh{哥} (gē, frère, \zh{可} + \zh{力}). \\ \hline
咖 & kā & \zh{口} (kǒu) & 8 & \zh{口} + \zh{加} (jiā). \textbf{Proche :} \zh{嘉} (jiā, louange, \zh{加} + \zh{口} en bas). \\ \hline
棉 & mián & \zh{木} (mù) & 11 & Gauche→droite. \zh{木} + \zh{免} (miǎn). \textbf{Proche :} \zh{绵} (mián, coton/continu, radical \zh{纟}). \\ \hline
花 & huā & \zh{艹} (cǎo) & 7 & Haut→bas. \zh{艹} (3) + \zh{化} (huà, 4). \textbf{Proche :} \zh{华} (huá, splendeur/Chine, \zh{十} + \zh{匕} + \zh{人}). \\ \hline
种 & zhòng/zhǒng & \zh{禾} (hé) & 9 & Gauche→droite. \zh{禾} (5) + \zh{中} (zhōng, 4). \textbf{Double lecture :} zhòng (v. cultiver) / zhǒng (n. espèce). \\ \hline
\end{tabularx}
\end{center}

\begin{methode}[Entraînement à l'écriture (15 min)]
\begin{enumerate}
    \item \textbf{Repérage :} Sur votre cahier, tracez la grille de 9 cases (米字格). Identifiez le radical et la structure (gauche-droite, haut-bas, enclos).
    \item \textbf{Traçage guidé :} Écrivez chaque caractère 5 fois en respectant scrupuleusement l'ordre des traits indiqué ci-dessus. Dites le nom des traits à voix haute (ex: \emph{héng, shù, piě, nà}).
    \item \textbf{Dictée de caractères :} Fermez le tableau. Le professeur dicte : « \zh{经济}, \zh{发展}, \zh{农业}, \zh{石油}, \zh{出口}, \zh{港口}, \zh{可可}, \zh{咖啡}, \zh{棉花}, \zh{种} (zhòng) ». Écrivez sans modèle.
    \item \textbf{Jeu des paires :} Associez le caractère à son radical : \zh{经}–\zh{纟}, \zh{济}–\zh{氵}, \zh{油}–\zh{氵}, \zh{港}–\zh{氵}, \zh{农}–\zh{曲}, \zh{棉}–\zh{木}, \zh{花}–\zh{艹}, \zh{种}–\zh{禾}.
\end{enumerate}
\end{methode}

\courssub{Texte d'application : L'économie du Cameroun en chiffres}

Le texte suivant réemploie le vocabulaire et les structures vus ci-dessus. Lisez-le à haute voix, identifiez les collocations et répondez aux questions.

\begin{textebox}[Texte original : 喀麦隆的经济概况]
喀麦隆位于中非，被称为“非洲缩影”。这里的经济以农业为主，也有石油工业。

喀麦隆的农业很发达。南方气候湿热，农民种可可、咖啡和香蕉。北方适合种棉花和花生。这些农产品是主要的出口产品，卖到欧洲和中国。

除了农业，石油也是重要收入。石油出口量很大，港口很忙。杜阿拉港口是最大的港口，很多船从那里出口石油和木材。

近年来，国家想发展工业和服务业。基础设施建设越来越快，修路、建桥、建电厂。经济发展得很快，人民生活越来越好。

但是，挑战也很多。比如：电力不够，运输成本高，青年失业率高。政府正在努力解决这些问题，希望喀麦隆成为中非经济强国。
\end{textebox}

\begin{center}
\begin{tabularx}{\linewidth}{|X|X|X|}
\hline
\rowcolor{popblueL} \textbf{Caractères} & \textbf{Pinyin} & \textbf{Traduction} \\ \hline
中非 & Zhōngfēi & Afrique centrale \\ \hline
缩影 & suōyǐng & miniature, reflet condensé \\ \hline
以...为主 & yǐ... wéi zhǔ & avoir pour base principale... \\ \hline
气候 & qìhòu & climat \\ \hline
湿热 & shīrè & chaud et humide \\ \hline
适合 & shìhé & convenir, être adapté \\ \hline
花生 & huāshēng & arachide \\ \hline
收入 & shōurù & revenu, recette \\ \hline
近年来 & jìnniánlái & ces dernières années \\ \hline
工业 & gōngyè & industrie \\ \hline
服务业 & fúwùyè & secteur des services \\ \hline
基础设施 & jīchǔ shèshī & infrastructures \\ \hline
建设 & jiànshè & construction, bâtir \\ \hline
修路 & xiū lù & construire des routes \\ \hline
建桥 & jiàn qiáo & construire des ponts \\ \hline
电厂 & diànchǎng & centrale électrique \\ \hline
挑战 & tiǎozhàn & défi \\ \hline
电力 & diànlì & électricité \\ \hline
运输成本 & yùnshū chéngběn & coût de transport \\ \hline
青年 & qīngnián & jeunesse, jeunes \\ \hline
失业率 & shīyèlǜ & taux de chômage \\ \hline
政府 & zhèngfǔ & gouvernement \\ \hline
努力 & nǔlì & s'efforcer, faire des efforts \\ \hline
解决 & jiějué & résoudre \\ \hline
强国 & qiángguó & grande puissance, pays fort \\ \hline
\end{tabularx}
\end{center}

\begin{exercice}[Questions de compréhension et d'analyse]
\begin{enumerate}
    \item Relevez dans le texte trois collocations \zh{动词 + 名词} (Verbe + Objet) liées à l'économie.
    \item Trouvez deux phrases utilisant la structure \zh{越来越 + Adjectif}. Traduisez-les.
    \item Identifiez la phrase contenant le complément de degré \zh{得}. Analysez sa structure (V + 得 + Adv + Adj).
    \item Pourquoi le texte utilise-t-il \zh{种} (zhòng) et non \zh{种} (zhǒng) ? Justifiez par deux exemples du texte.
\end{enumerate}
\end{exercice}

\begin{corrige}[Exercice de compréhension et d'analyse]
\begin{enumerate}
    \item \zh{种可可} (cultiver le cacao), \zh{种咖啡} (cultiver le café), \zh{种棉花} (cultiver le coton), \zh{种花生} (cultiver l'arachide), \zh{出口石油} (exporter le pétrole), \zh{出口产品} (exporter des produits), \zh{发展工业} (développer l'industrie), \zh{发展服务业} (développer les services), \zh{解决问题} (résoudre les problèmes).
    \item 1. \zh{基础设施建设越来越快} (La construction d'infrastructures est de plus en plus rapide). 2. \zh{人民生活越来越好} (La vie du peuple va de mieux en mieux).
    \item \zh{经济发展得很快} (L'économie se développe très vite). Structure : Sujet \zh{经济} + Verbe \zh{发展} + \zh{得} + Adv. \zh{很} + Adj. \zh{快}.
    \item Parce que \zh{种} est utilisé comme \textbf{verbe} (4e ton) signifiant « cultiver ». Exemples du texte : \zh{农民种可可} (les paysans cultivent le cacao), \zh{种棉花和花生} (cultiver le coton et l'arachide). \zh{种} (zhǒng, 3e ton) est un nom/classifieur (ex: \zh{一种} — un type).
\end{enumerate}
\end{corrige}

\courssub{Exercice résolu : Thème guidé (Français $\rightarrow$ Chinois)}

\begin{exoresolu}[Traduire les phrases suivantes en chinois]
\textbf{Énoncé :}
\begin{enumerate}
    \item Le Cameroun exporte du cacao, du café et du coton.
    \item L'agriculture est la base de l'économie.
    \item Ces dernières années, le développement économique est de plus en plus rapide.
    \item Les paysans du Nord cultivent le coton et les arachides.
    \item Le port de Douala est très grand ; beaucoup de pétrole y est exporté.
\end{enumerate}

\tcblower

\textbf{Solution détaillée :}

\begin{enumerate}
    \item \zh{喀麦隆出口可可、咖啡和棉花。}
    \begin{itemize}
        \item \zh{喀麦隆} (Sujet) + \zh{出口} (Verbe) + \zh{可可、咖啡和棉花} (Objets coordonnés par \zh{和} / \zh{、}).
        \item \emph{Note :} Pas de classifieur nécessaire pour les matières premières en vrac ici.
    \end{itemize}
    \item \zh{农业是经济的基础。} \emph{ou} \zh{经济以农业为主。}
    \begin{itemize}
        \item Structure \zh{是} (shì) pour définir : A est la base de B $\rightarrow$ A \zh{是} B \zh{的} 基础 \zh{(jīchǔ)}.
        \item Structure \zh{以...为主} (yǐ... wéi zhǔ) : « Prendre ... comme principal » $\rightarrow$ \zh{经济以农业为主。}
    \end{itemize}
    \item \zh{近年来，经济发展越来越快。}
    \begin{itemize}
        \item \zh{近年来} (Complément de temps en tête de phrase).
        \item \zh{经济发展} (Sujet « Développement économique » / Verbe « L'économie se développe »).
        \item \zh{越来越快} (yuè lái yuè kuài) — \textbf{Pas de \zh{很} !} Voir \textbf{Attention} ci-dessus.
    \end{itemize}
    \item \zh{北方的农民种棉花和花生。}
    \begin{itemize}
        \item \zh{北方} (Nord) + \zh{的} + \zh{农民} (Paysans).
        \item \zh{种} (zhòng, 4e ton) = cultiver.
        \item \zh{花生} (huāshēng) = arachides.
    \end{itemize}
    \item \zh{杜阿拉港口很大，很多石油从那里出口。}
    \begin{itemize}
        \item \zh{杜阿拉港口} (Sujet 1) + \zh{很大} (Adj. prédicatif).
        \item \zh{很多石油} (Sujet 2, objet logique) + \zh{从那里} (De là-bas) + \zh{出口} (Verbe passif sans \zh{被} bèi, structure courante : Objet + Lieu + Verbe).
        \item Alternative passive formelle : \zh{很多石油被从那里出口。} (Moins naturel à l'oral).
    \end{itemize}
\end{enumerate}
\end{exoresolu}

\courssub{Production écrite guidée : Mini-paragraphe}

\begin{methode}[Rédiger un paragraphe de 5-6 phrases sur l'économie camerounaise]
\begin{enumerate}
    \item \textbf{Situation :} Présentez le pays et son statut (\zh{位于}, \zh{被称为}).
    \item \textbf{Secteur principal :} Agriculture (\zh{以...为主}, \zh{种}, \zh{出口}).
    \item \textbf{Second secteur :} Pétrole / Port (\zh{石油}, \zh{港口}, \zh{出口}).
    \item \textbf{Tendance :} Progrès récent (\zh{近年来}, \zh{越来越}, \zh{发展得...}).
    \item \textbf{Défi / Avenir :} Problème + Action gouvernementale (\zh{但是}, \zh{挑战}, \zh{政府}, \zh{解决}, \zh{希望}).
\end{enumerate}
\emph{Mots-clés à utiliser :} \zh{喀麦隆, 农业, 可可, 咖啡, 棉花, 石油, 港口, 出口, 发展, 越来越, 经济, 国家, 政府}.
\end{methode}

\begin{experience}[Débat oral : « Cameroun vs Tchad / Cameroun vs Chine »]
\textbf{Consigne :} En binôme, préparez un dialogue de 2 minutes. L'élève A est un expert économique camerounais, l'élève B un journaliste chinois.
\textbf{Sujet :} \zh{为什么喀麦隆经济比泰发展？} (Pourquoi l'économie du Cameroun se développe-t-elle plus vite que celle du Tchad ?) \emph{ou} \zh{喀麦隆和中国经济合作的未来} (Avenir de la coopération économique Cameroun-Chine).
\textbf{Structures imposées :} \zh{因为...所以...}, \zh{虽然...但是...}, \zh{越来越...}, \zh{发展得...}, \zh{出口...}.
\textbf{Évaluation :} Prononciation des tons (zhòng/zhǒng, jīngjì/fāzhǎn), fluidité, richesse lexicale.
\end{experience}

\begin{aretenir}[Synthèse de la leçon 34 — Points clés pour l'examen]
\begin{itemize}
    \item \textbf{Vocabulaire :} Maîtriser l'orthographe (caractères), le pinyin (tons) et le sens des 16 mots du glossaire + les 15 mots du texte.
    \item \textbf{Grammaire :} 1) Complément de degré \zh{V + 得 + Adv + Adj} (\zh{发展得很快}). 2) Adverbe de progression \zh{越来越 + Adj} (sans \zh{很}). 3) Distinction \zh{发展经济} (V+O) vs \zh{经济发展} (Sujet). 4) Double lecture \zh{种} (zhòng/zhǒng).
    \item \textbf{Écriture :} Ordre des traits des 17 caractères du tableau ; radicaux \zh{纟, 氵, 木, 艹, 禾, 口, 扌} (pour \zh{挑} dans \zh{挑战}).
    \item \textbf{Expression écrite :} Structure en 5 étapes (Situation, Agriculture, Pétrole, Tendance, Défi) + connecteurs logiques (\zh{除了...}, \zh{近年来}, \zh{但是}, \zh{希望}).
    \item \textbf{Traduction :} Savoir passer du thème (Fr$\rightarrow$Ch) et version (Ch$\rightarrow$Fr) sur les phrases types du \texttt{exoresolu}.
\end{itemize}
\end{aretenir}

\begin{savaistu}[Le Cameroun : « L'Afrique en miniature »]
Le surnom \zh{非洲缩影} (Fēizhōu suōyǐng) — « L'Afrique en miniature » ou « Afrique condensée » — n'est pas seulement touristique, il est \textbf{géographiquement et économiquement exact}.
\begin{itemize}
    \item \textbf{Géographie :} Le Cameroun possède \textbf{tous} les grands types de climats et paysages africains : côte océane (Sud-Ouest), forêt équatoriale (Sud/Est), savane herbeuse (Adamaoua), zone sahélienne (Extrême-Nord), montagnes volcaniques (Mont Cameroun, 4095 m).
    \item \textbf{Agriculture :} Cette diversité permet de produire \textbf{presque toutes les cultures d'Afrique} : cultures d'exportation (cacao, café, coton, banane, caoutchouc, huile de palme) ET vivrières (manioc, maïs, mil, sorgho, igname, plantain, riz, pomme de terre).
    \item \textbf{Comparaison Chine :} La Chine a aussi une grande diversité (riz au sud, blé au nord), mais le Cameroun concentre cette diversité sur une surface 4,5 fois plus petite (475 000 km² vs 9,6 M km²), ce qui en fait un « laboratoire » agricole unique.
\end{itemize}
\emph{Ce contexte justifie pourquoi le vocabulaire \zh{种可可, 种咖啡, 种棉花, 种花生} est si central dans votre programme HSK 3-4 camerounais !}
\end{savaistu}

\coursec{Écriture correcte des caractères}

Dans la section précédente, nous avons appris le vocabulaire de l'économie du Cameroun : \zh{经济} (jīngjì, « économie »), \zh{发展} (fāzhǎn, « développement »), \zh{石油} (shíyóu, « pétrole »), \zh{港口} (gǎngkǒu, « port »), \zh{棉花} (miánhua, « coton »). Savoir lire ces mots ne suffit pas : à l'examen comme dans la vie réelle, il faut aussi savoir les \cle{écrire correctement}, trait par trait. Un caractère mal écrit peut devenir un autre caractère et changer complètement le sens de la phrase. Cette section vous donne la méthode, les règles et l'entraînement nécessaires.

\courssub{Observer avant d'écrire : les radicaux du thème}

Regardez attentivement les caractères du titre de notre leçon : \zh{为什么喀麦隆经济比较发展？} (Wèishénme Kāmàilóng jīngjì bǐjiào fāzhǎn ? — « Pourquoi l'économie du Cameroun se développe-t-elle relativement bien ? »). Vous remarquez que certains petits éléments reviennent souvent : \zh{纟} dans \zh{经}, \zh{氵} dans \zh{济}, \zh{油} et \zh{港}, \zh{木} dans \zh{棉}. Ces éléments sont les \cle{radicaux} (ou clés, \zh{部首} bùshǒu) : ils donnent une indication de \textbf{sens}.

\begin{definition}[Les radicaux clés de la leçon]
Un \cle{radical} (\zh{部首} bùshǒu) est la partie d'un caractère qui indique sa famille de sens. Le reste du caractère donne souvent une indication de \textbf{prononciation} (partie phonétique). Connaître les radicaux permet de deviner le sens d'un caractère inconnu et de le mémoriser plus facilement.
\end{definition}

\begin{center}
\begin{tabularx}{\linewidth}{|X|X|X|X|}
\hline
\rowcolor{popblueL} Caractère & Radical & Sens du radical & Partie phonétique \\
\hline
\zh{经} jīng & \zh{纟} (soie) & fil, tissu → « organiser, économie » & \zh{巠} jīng \\
\hline
\zh{济} jì & \zh{氵} (eau) & eau, rivière → « traverser, aider » & \zh{齐} qí \\
\hline
\zh{发} fā & \zh{又} (main) & action de la main → « émettre, développer » & (caractère simplifié, phonétique non apparent) \\
\hline
\zh{展} zhǎn & \zh{尸} (corps/toit) & corps, toit → « déployer, étendre » & \zh{单} dān/shàn \\
\hline
\zh{油} yóu & \zh{氵} (eau) & liquide → « huile, pétrole » & \zh{由} yóu \\
\hline
\zh{港} gǎng & \zh{氵} (eau) & eau → « port » & \zh{巷} xiàng \\
\hline
\zh{棉} mián & \zh{木} (bois) & plante, arbre → « coton » & \zh{帛} bó \\
\hline
\end{tabularx}
\end{center}

Observez la logique : \zh{油} (huile) et \zh{港} (port) ont tous les deux le radical de l'eau \zh{氵}, car l'huile est un liquide et le port est au bord de l'eau. \zh{棉} (coton) a le radical du bois \zh{木}, car le coton est une plante. \zh{经} a le radical de la soie \zh{纟}, car autrefois « organiser les fils » voulait dire « gérer » : de là vient le sens moderne d'« économie ».

\begin{exemplebox}[经济 : un mot, deux radicaux]
Le mot \zh{经济} (jīngjì, « économie ») s'écrit avec le radical de la soie \zh{纟} à gauche de \zh{经} et le radical de l'eau \zh{氵} à gauche de \zh{济}.\\
\zh{喀麦隆的经济在发展。}\\
\emph{Kāmàilóng de jīngjì zài fāzhǎn.}\\
« L'économie du Cameroun est en train de se développer. »
\end{exemplebox}

\courssub{La règle générale d'écriture}

\begin{propriete}[Ordre général des traits]
Les caractères chinois s'écrivent selon des règles fixes :
\begin{enumerate}
\item de \cle{haut en bas} : on écrit d'abord la partie supérieure ;
\item de \cle{gauche à droite} : on écrit d'abord le radical de gauche (\zh{氵}, \zh{纟}, \zh{木}) ;
\item l'\cle{extérieur avant l'intérieur} : pour les caractères avec un élément qui enveloppe (comme \zh{尸} dans \zh{展}), on trace d'abord l'enveloppe, puis l'intérieur ;
\item l'horizontal \zh{一} avant le vertical \zh{丨} quand ils se croisent (exemple : \zh{十}).
\end{enumerate}
Exception fréquente : pour les caractères entièrement fermés comme \zh{国} (pays), on trace l'extérieur, l'intérieur, puis on « ferme la porte » par le trait horizontal du bas en dernier.
\end{propriete}

Comparaison avec le français : en français, vous écrivez les lettres d'un mot de gauche à droite, une par une. En chinois, chaque caractère est un petit dessin construit selon un plan précis : respecter l'ordre des traits rend l'écriture plus rapide, plus belle et plus facile à mémoriser, car la main garde le « chemin » en mémoire.

\courssub{Ordre des traits : 经 (8 traits) et 济 (9 traits)}

Le caractère \zh{经} (jīng) s'écrit en 8 traits, en deux temps : d'abord le radical \zh{纟} à gauche (3 traits : deux \zh{撇折} et un \zh{提}), puis la partie droite \zh{又} (2 traits : \zh{横撇} puis \zh{捺}) et enfin \zh{工} en bas à droite (3 traits : horizontal, vertical, horizontal).

Pour \zh{济} (jì), on écrit d'abord les trois gouttes d'eau \zh{氵} à gauche (3 traits : point, point, \zh{提} vers le haut), puis la partie droite \zh{齐} (6 traits), de haut en bas. Total : 9 traits. C'est la règle « gauche avant droite ».

\begin{popfigure}
\begin{center}
\begin{tikzpicture}[scale=0.9]
% Grille 田字格 pour 经
\draw[popline, thick] (0,0) rectangle (3,3);
\draw[popline, dashed] (1.5,0) -- (1.5,3);
\draw[popline, dashed] (0,1.5) -- (3,1.5);
% Trait 1 : 撇折
\draw[->, popink, very thick] (0.55,2.75) -- (0.85,2.15) -- (0.6,1.95);
\node[popink, font=\small] at (0.35,2.8) {1};
% Trait 2 : 撇折
\draw[->, poporange, very thick] (0.5,1.85) -- (0.8,1.3) -- (0.55,1.1);
\node[poporange, font=\small] at (0.3,1.85) {2};
% Trait 3 : 提
\draw[->, popgreen, very thick] (0.5,0.7) -- (1.05,1.0);
\node[popgreen, font=\small] at (0.4,0.5) {3};
% Trait 4 : 横撇 de 又
\draw[->, popblue, very thick] (1.7,2.7) -- (2.5,2.7) -- (2.0,2.05);
\node[popblue, font=\small] at (1.65,2.9) {4};
% Trait 5 : 捺
\draw[->, poppurple, very thick] (2.15,2.55) -- (2.75,1.95);
\node[poppurple, font=\small] at (2.85,2.2) {5};
% Trait 6 : 横 de 工
\draw[->, popgold, very thick] (1.75,1.45) -- (2.65,1.45);
\node[popgold, font=\small] at (2.2,1.65) {6};
% Trait 7 : 竖
\draw[->, poppink, very thick] (2.2,1.45) -- (2.2,0.6);
\node[poppink, font=\small] at (2.4,1.0) {7};
% Trait 8 : 横 bas
\draw[->, popdark, very thick] (1.7,0.4) -- (2.7,0.4);
\node[popdark, font=\small] at (2.2,0.2) {8};
\node[below] at (1.5,-0.15) {\zh{经} jīng — 8 traits};
\end{tikzpicture}
\hspace{1cm}
\begin{tikzpicture}[scale=0.9]
% Grille 田字格 pour 济
\draw[popline, thick] (0,0) rectangle (3,3);
\draw[popline, dashed] (1.5,0) -- (1.5,3);
\draw[popline, dashed] (0,1.5) -- (3,1.5);
% Radical 氵 (3 traits)
\draw[->, popink, very thick] (0.5,2.6) -- (0.7,2.3);
\node[popink, font=\small] at (0.3,2.7) {1};
\draw[->, poporange, very thick] (0.9,2.5) -- (1.1,2.2);
\node[poporange, font=\small] at (1.2,2.6) {2};
\draw[->, popgreen, very thick] (0.6,1.8) -- (1.2,2.1);
\node[popgreen, font=\small] at (0.4,1.7) {3};
% Partie droite 齐 (6 traits) - simplification visuelle
% Trait 4: 横 (top of 齐)
\draw[->, popblue, very thick] (1.7,2.6) -- (2.7,2.6);
\node[popblue, font=\small] at (2.2,2.8) {4};
% Trait 5: 竖 (left down)
\draw[->, poppurple, very thick] (1.8,2.6) -- (1.8,1.6);
\node[poppurple, font=\small] at (1.6,2.1) {5};
% Trait 6: 横 (middle)
\draw[->, popgold, very thick] (1.8,2.1) -- (2.6,2.1);
\node[popgold, font=\small] at (2.3,2.25) {6};
% Trait 7: 竖 (middle down)
\draw[->, poppink, very thick] (2.2,2.1) -- (2.2,1.1);
\node[poppink, font=\small] at (2.4,1.6) {7};
% Trait 8: 横 (bottom of top part)
\draw[->, popdark, very thick] (1.8,1.6) -- (2.6,1.6);
\node[popdark, font=\small] at (2.2,1.4) {8};
% Trait 9: 竖 (main stem) - simplified representation
\draw[->, popink, very thick] (2.2,1.6) -- (2.2,0.4);
\node[poporange, font=\small] at (2.5,1.0) {9};
\node[below] at (1.5,-0.15) {\zh{济} jì — 9 traits};
\end{tikzpicture}
\end{center}
\legende{Ordre des traits de \zh{经} (8 traits) : d'abord \zh{纟} à gauche, puis \zh{又}, puis \zh{工} ; et de \zh{济} (9 traits) : d'abord \zh{氵} à gauche (3 traits), puis \zh{齐} à droite (6 traits, de haut en bas).}
\end{popfigure}

\begin{popfigure}
\begin{center}
\begin{tikzpicture}[scale=0.9]
% Grille 田字格 pour 展
\draw[popline, thick] (0,0) rectangle (3,3);
\draw[popline, dashed] (1.5,0) -- (1.5,3);
\draw[popline, dashed] (0,1.5) -- (3,1.5);
% Trait 1 : 横折 de 尸
\draw[->, popink, very thick] (0.7,2.75) -- (2.3,2.75) -- (2.2,2.35);
\node[popink, font=\small] at (0.55,2.9) {1};
% Trait 2 : 横 de 尸
\draw[->, poporange, very thick] (0.75,2.35) -- (2.15,2.35);
\node[poporange, font=\small] at (1.45,2.5) {2};
% Trait 3 : 撇 de 尸
\draw[->, popgreen, very thick] (0.75,2.75) -- (0.45,0.7);
\node[popgreen, font=\small] at (0.3,1.7) {3};
% Traits 4-5-6 : intérieur haut (单)
\draw[->, popblue, very thick] (1.0,1.95) -- (2.3,1.95);
\node[popblue, font=\small] at (2.45,2.05) {4};
\draw[->, poppurple, very thick] (1.35,2.05) -- (1.35,1.55);
\node[poppurple, font=\small] at (1.2,1.8) {5};
\draw[->, popgold, very thick] (1.95,2.05) -- (1.95,1.55);
\node[popgold, font=\small] at (2.15,1.8) {6};
% Trait 7 : 横
\draw[->, poppink, very thick] (0.95,1.35) -- (2.35,1.35);
\node[poppink, font=\small] at (2.5,1.45) {7};
% Trait 8 : 竖提
\draw[->, popdark, very thick] (1.15,1.2) -- (1.15,0.45) -- (1.5,0.6);
\node[popdark, font=\small] at (0.95,0.8) {8};
% Trait 9 : 撇
\draw[->, popblue, very thick] (1.75,1.15) -- (1.35,0.35);
\node[popblue, font=\small] at (1.85,0.75) {9};
% Trait 10 : 捺
\draw[->, poporange, very thick] (1.85,1.1) -- (2.55,0.35);
\node[poporange, font=\small] at (2.6,0.8) {10};
\node[below] at (1.5,-0.15) {\zh{展} zhǎn — 10 traits};
\end{tikzpicture}
\end{center}
\legende{Ordre des traits de \zh{展} (10 traits) : d'abord l'enveloppe \zh{尸} (traits 1-3), puis l'intérieur de haut en bas (traits 4-10). Règle : extérieur avant intérieur.}
\end{popfigure}

\begin{methode}[Comment mémoriser un caractère en 4 étapes]
\begin{enumerate}
\item \cle{Identifier le radical} : il donne le sens (\zh{氵} = eau dans \zh{油} et \zh{港}).
\item \cle{Identifier la partie phonétique} : elle donne le son (\zh{由} yóu dans \zh{油} yóu ; \zh{巷} xiàng dans \zh{港} gǎng ; \zh{单} dān dans \zh{展} zhǎn).
\item \cle{Découper le caractère} en blocs : \zh{经} = \zh{纟} + \zh{又} + \zh{工} ; \zh{济} = \zh{氵} + \zh{齐} ; \zh{展} = \zh{尸} + \zh{单}.
\item \cle{Écrire 5 fois} le caractère en respectant l'ordre des traits, à voix haute en comptant les traits, dans une case \zh{田字格}.
\end{enumerate}
Cette méthode « sens + son + geste » est bien plus efficace que la copie mécanique.
\end{methode}

\courssub{Caractères proches à ne pas confondre}

Certains caractères se ressemblent énormément mais ont un sens et une prononciation différents. Les francophones, habitués à un alphabet de 26 lettres bien distinctes, les confondent souvent. Voici les paires dangereuses de notre thème.

\begin{center}
\begin{tabularx}{\linewidth}{|X|X|X|X|}
\hline
\rowcolor{popblueL} Caractères & Pinyin & Différence graphique & Traduction \\
\hline
\zh{己} / \zh{已} / \zh{巳} & jǐ / yǐ / sì & ouvert en bas / à moitié fermé / totalement fermé & soi-même / déjà / (branche terrestre) \\
\hline
\zh{经} / \zh{轻} & jīng / qīng & radical \zh{纟} (soie) / radical \zh{车} (voiture) & économie / léger \\
\hline
\zh{发} / \zh{友} & fā / yǒu & \zh{发} a un trait horizontal + point en haut à droite ; \zh{友} a deux points (\zh{丷}) & développer / ami \\
\hline
\zh{油} / \zh{由} & yóu / yóu & \zh{油} a le radical \zh{氵} en plus & huile / depuis, par \\
\hline
\end{tabularx}
\end{center}

\begin{popfigure}
\begin{center}
\begin{tikzpicture}[scale=1.0]
% 己 ouvert
\draw[popline, thick] (0,0) rectangle (2,2);
\draw[popline, dashed] (1,0) -- (1,2);
\draw[popline, dashed] (0,1) -- (2,1);
\draw[popink, very thick] (0.5,1.6) -- (1.5,1.6) -- (1.45,1.15);
\draw[popink, very thick] (0.55,1.15) -- (1.45,1.15);
\draw[popink, very thick] (0.55,1.15) -- (0.55,0.45) -- (1.35,0.45);
\node[below, align=center] at (1,-0.1) {\zh{己} jǐ\\ ouvert};
% 已 moitié fermé
\draw[popline, thick] (3,0) rectangle (5,2);
\draw[popline, dashed] (4,0) -- (4,2);
\draw[popline, dashed] (3,1) -- (5,1);
\draw[poporange, very thick] (3.5,1.6) -- (4.5,1.6) -- (4.45,1.15);
\draw[poporange, very thick] (3.55,1.15) -- (4.45,1.15);
\draw[poporange, very thick] (3.55,1.15) -- (3.55,0.45) -- (4.25,0.45);
\node[below, align=center] at (4,-0.1) {\zh{已} yǐ\\ à moitié fermé};
% 巳 fermé
\draw[popline, thick] (6,0) rectangle (8,2);
\draw[popline, dashed] (7,0) -- (7,2);
\draw[popline, dashed] (6,1) -- (8,1);
\draw[popgreen, very thick] (6.5,1.6) -- (7.5,1.6) -- (7.45,1.15);
\draw[popgreen, very thick] (6.55,1.15) -- (7.45,1.15);
\draw[popgreen, very thick] (6.55,1.15) -- (6.55,0.45) -- (7.45,0.45) -- (7.45,1.05);
\node[below, align=center] at (7,-0.1) {\zh{巳} sì\\ fermé};
\end{tikzpicture}
\end{center}
\legende{Comparaison de \zh{己} (ouvert), \zh{已} (à moitié fermé) et \zh{巳} (fermé) : seul le bas du caractère change.}
\end{popfigure}

\begin{attention}[己 / 已 : l'erreur classique des francophones]
\zh{己} (jǐ, « soi-même », comme dans \zh{自己} zìjǐ) est \textbf{ouvert en bas} : le trait vertical ne remonte pas. \zh{已} (yǐ, « déjà », comme dans \zh{已经} yǐjing) est \textbf{à moitié fermé} : le trait remonte un peu. \zh{巳} (sì) est \textbf{totalement fermé}. Astuce de mémorisation : « \zh{己} jǐ, je m'ouvre aux autres ; \zh{已} yǐ, c'est déjà presque fini, la porte se ferme ». Attention aussi à \zh{已经} yǐjing : il s'écrit avec \zh{已} (fermé à moitié) et non \zh{己} !
\end{attention}

Autres pièges : \zh{发} (fā, développer) a un petit trait horizontal et un point en plus par rapport à \zh{友} (yǒu, ami) — ne transformez pas « développement » en « amitié » ! Et \zh{油} (yóu, huile, pétrole) se prononce exactement comme \zh{由} (yóu) : seul le radical \zh{氵} les distingue à l'écrit. Dans \zh{石油} (shíyóu, pétrole), n'oubliez jamais les trois gouttes d'eau.

\courssub{Exercice résolu et copie guidée}

\begin{exoresolu}[Identifier radicaux et ordre des traits]
Énoncé : pour les caractères \zh{港} et \zh{棉}, indiquez le radical, la partie phonétique et le côté par lequel on commence à écrire.
\tcblower
Solution détaillée :
\begin{itemize}
\item \zh{港} (gǎng, port) : radical \zh{氵} (eau) à gauche, partie phonétique \zh{巷} à droite. On commence par les trois gouttes d'eau à gauche (règle « gauche avant droite »), puis on écrit la partie droite de haut en bas.
\item \zh{棉} (mián, coton) : radical \zh{木} (bois) à gauche, partie phonétique \zh{帛} à droite. On commence par \zh{木} (horizontal, vertical, puis les deux traits obliques), puis la partie droite.
\end{itemize}
Dans les deux cas, le radical de gauche donne le sens : l'eau pour le port, la plante pour le coton.
\end{exoresolu}

\begin{exercice}[Copie guidée en cases 田字格]
Sur votre cahier, tracez des cases \zh{田字格} (un carré divisé en quatre par des pointillés). Copiez chaque caractère 5 fois en comptant les traits à voix haute : \zh{经} (8 traits), \zh{济} (9 traits), \zh{展} (10 traits), \zh{油} (8 traits), \zh{港} (12 traits), \zh{棉} (12 traits). Puis écrivez le mot complet \zh{经济发展} (jīngjì fāzhǎn, « développement économique ») sans faute, trois fois.
\end{exercice}

\begin{corrige}[Exercice de copie guidée]
Vérifiez votre copie point par point : (1) dans \zh{经}, le radical \zh{纟} doit rester étroit à gauche, et \zh{工} se trouve en bas à droite, sous \zh{又} ; (2) dans \zh{济}, les trois gouttes \zh{氵} sont alignées en diagonale, le dernier trait monte vers la droite ; (3) dans \zh{展}, le grand trait oblique \zh{撇} de \zh{尸} descend bien à gauche et l'intérieur reste compact ; (4) dans \zh{油}, le radical \zh{氵} ne doit pas toucher \zh{由}. Si un caractère « déborde » de sa case, recommencez en réduisant le radical de gauche : il n'occupe qu'un tiers de la largeur.
\end{corrige}

\begin{aretenir}[L'essentiel de l'écriture des caractères du thème]
\begin{itemize}
\item Le \cle{radical} donne le sens : \zh{纟} soie (\zh{经}), \zh{氵} eau (\zh{济}, \zh{油}, \zh{港}), \zh{木} bois (\zh{棉}), \zh{尸} enveloppe (\zh{展}).
\item Ordre des traits : gauche avant droite, haut avant bas, extérieur avant intérieur.
\item \zh{经} : 8 traits (\zh{纟} + \zh{又} + \zh{工}) ; \zh{济} : 9 traits (\zh{氵} + \zh{齐}) ; \zh{展} : 10 traits (\zh{尸} + intérieur).
\item Ne pas confondre : \zh{己} ouvert / \zh{已} à moitié fermé ; \zh{发} / \zh{友} ; \zh{油} / \zh{由} ; \zh{经} / \zh{轻}.
\item Méthode : sens + son + 5 écritures comptées en case \zh{田字格}.
\end{itemize}
\end{aretenir}

Maintenant que vous savez écrire correctement les caractères du thème, la section suivante vous donnera les structures et les connecteurs pour les assembler en phrases et en texte sur l'économie du Cameroun.

\coursec{Structures et connecteurs pour écrire}

Dans la section précédente, nous avons appris à écrire correctement les caractères du thème économique (\zh{经济}, \zh{发展}, \zh{石油}, \zh{农业}…). Savoir écrire les mots ne suffit pas : il faut maintenant les \textbf{relier entre eux} pour construire un texte cohérent. Un bon paragraphe chinois, comme un bon paragraphe français, repose sur des \cle{connecteurs logiques} : la cause, l'addition, l'opposition, l'opinion, la comparaison, la progression. C'est l'objet de cette section : vous donner la « boîte à outils » grammaticale pour rédiger votre texte sur l'économie du Cameroun.

\courssub{Observer avant de comprendre : un mini-texte}

\begin{textebox}[L'économie du Cameroun : un exemple de paragraphe]
喀麦隆有很多石油，也有很多农业产品。\\
Kāmàilóng yǒu hěn duō shíyóu, yě yǒu hěn duō nóngyè chǎnpǐn.\\
« Le Cameroun a beaucoup de pétrole et aussi beaucoup de produits agricoles. »\\[4pt]
因为喀麦隆有很多资源，所以经济发展得越来越快。\\
Yīnwèi Kāmàilóng yǒu hěn duō zīyuán, suǒyǐ jīngjì fāzhǎn de yuè lái yuè kuài.\\
« Parce que le Cameroun a beaucoup de ressources, son économie se développe de plus en plus vite. »\\[4pt]
喀麦隆经济比泰国发展得慢，但是越来越好。\\
Kāmàilóng jīngjì bǐ Tàiguó fāzhǎn de màn, dànshì yuè lái yuè hǎo.\\
« L'économie du Cameroun se développe plus lentement que celle de la Thaïlande, mais elle va de mieux en mieux. »\\[4pt]
我觉得喀麦隆的经济会越来越好。\\
Wǒ juéde Kāmàilóng de jīngjì huì yuè lái yuè hǎo.\\
« Je pense que l'économie du Cameroun ira de mieux en mieux. »
\end{textebox}

Observez les mots en gras logique : \zh{因为…所以…}, \zh{也}, \zh{比…得…}, \zh{但是}, \zh{越来越}, \zh{我觉得}. Ce sont exactement les outils de cette leçon. Étudions-les un par un.

\courssub{Le connecteur de cause : 因为…所以…}

\begin{propriete}[Exprimer la cause et la conséquence]
Structure : \cle{因为} + cause，\cle{所以} + résultat。\\
En chinois, contrairement au français, on peut (et on préfère souvent) utiliser \zh{因为} et \zh{所以} \textbf{ensemble dans la même phrase}. En français, « parce que… donc… » est une faute de syntaxe ; en chinois, c'est la construction la plus naturelle et la plus élégante pour marquer le lien logique fort.
\begin{itemize}
\item \zh{因为喀麦隆有很多石油，所以经济发展得很快。} (\emph{Yīnwèi Kāmàilóng yǒu hěn duō shíyóu, suǒyǐ jīngjì fāzhǎn de hěn kuài.}) « Parce que le Cameroun a beaucoup de pétrole, son économie se développe vite. »
\item On peut aussi n'utiliser que \zh{所以} seul si la cause est déjà connue ou évoquée juste avant : \zh{喀麦隆有很多资源，所以经济越来越好。} (\emph{Kāmàilóng yǒu hěn duō zīyuán, suǒyǐ jīngjì yuè lái yuè hǎo.}) « Le Cameroun a beaucoup de ressources, donc son économie va de mieux en mieux. »
\end{itemize}
\end{propriete}

\begin{popfigure}
\begin{center}
\begin{tikzpicture}[node distance=1.2cm]
\node[draw=popblue, fill=popblueL, rounded corners, thick, align=center, minimum width=4.2cm, minimum height=1.1cm] (cause) {因为 + cause\\ \emph{yīnwèi} « parce que »};
\node[draw=poporange, fill=poporangeL, rounded corners, thick, align=center, minimum width=4.2cm, minimum height=1.1cm, right=3.2cm of cause] (res) {所以 + résultat\\ \emph{suǒyǐ} « donc »};
\draw[-{Stealth[length=3mm]}, very thick, popdark] (cause) -- node[above, align=center]{\small la cause entraîne\\ \small le résultat} (res);
\end{tikzpicture}
\end{center}
\legende{Schéma de la relation de cause : 因为 [cause] → 所以 [résultat]}
\end{popfigure}

\begin{exemplebox}[Exemples traduits]
\begin{itemize}
\item \zh{因为喀麦隆有很多农业产品，所以经济越来越好。} (\emph{Yīnwèi Kāmàilóng yǒu hěn duō nóngyè chǎnpǐn, suǒyǐ jīngjì yuè lái yuè hǎo.}) « Parce que le Cameroun a beaucoup de produits agricoles, son économie va de mieux en mieux. »
\item \zh{因为杜阿拉有一个大港口，所以很多公司在那里。} (\emph{Yīnwèi Dù'ālā yǒu yí ge dà gǎngkǒu, suǒyǐ hěn duō gōngsī zài nàlǐ.}) « Parce que Douala a un grand port, beaucoup d'entreprises s'y trouvent. »
\item \zh{因为年轻人学习汉语，所以他们可以找到好工作。} (\emph{Yīnwèi niánqīngrén xuéxí Hànyǔ, suǒyǐ tāmen kěyǐ zhǎodào hǎo gōngzuò.}) « Parce que les jeunes apprennent le chinois, ils peuvent trouver un bon travail. »
\end{itemize}
\end{exemplebox}

\courssub{Le connecteur d'addition : 而且 et 也}

\begin{propriete}[Ajouter une information]
\begin{itemize}
\item \cle{而且} (\emph{érqiě}) « de plus, et aussi » : se place \textbf{en tête de la deuxième proposition}, jamais entre deux noms. Structure : phrase 1，\textbf{而且} + phrase 2。
\item \cle{也} (\emph{yě}) « aussi » : se place \textbf{avant le verbe} (ou l'auxiliaire), à l'intérieur de la proposition.
\end{itemize}
\zh{喀麦隆有石油，而且有很多农业产品。} (\emph{Kāmàilóng yǒu shíyóu, érqiě yǒu hěn duō nóngyè chǎnpǐn.}) « Le Cameroun a du pétrole, et de plus beaucoup de produits agricoles. »\\
\zh{喀麦隆有石油，也有很多农业产品。} (\emph{Kāmàilóng yǒu shíyóu, yě yǒu hěn duō nóngyè chǎnpǐn.}) « Le Cameroun a du pétrole et aussi beaucoup de produits agricoles. »
\end{propriete}

\begin{attention}[而且 n'est pas « et » entre noms]
Pour relier deux \textbf{noms}, on utilise \zh{和} (\emph{hé}) : \zh{石油和农业产品} « le pétrole et les produits agricoles ». Pour relier deux \textbf{phrases/propositions} et insister sur l'addition de deux faits, on utilise \zh{而且} ou la simple virgule chinoise \zh{，}. Ne dites jamais \zh{喀麦隆有石油和很多农业产品} si vous voulez faire deux affirmations distinctes : préférez \zh{而且} ou deux phrases.
\end{attention}

\courssub{Le connecteur d'opposition : 但是}

\begin{propriete}[Exprimer le contraste]
\cle{但是} (\emph{dànshì}) « mais, cependant » se place en tête de la deuxième proposition : phrase 1，\textbf{但是} + phrase 2。 Sa forme courte \zh{但} (\emph{dàn}) est plus littéraire ; à l'écrit scolaire courant, utilisez \zh{但是}.\\
\zh{喀麦隆经济比泰国发展得慢，但是发展得越来越好。} (\emph{Kāmàilóng jīngjì bǐ Tàiguó fāzhǎn de màn, dànshì fāzhǎn de yuè lái yuè hǎo.}) « L'économie du Cameroun se développe plus lentement que celle de la Thaïlande, mais elle se développe de mieux en mieux. »
\end{propriete}

\begin{exemplebox}[Autres exemples]
\begin{itemize}
\item \zh{雅温得很大，但是杜阿拉是经济首都。} (\emph{Yǎwēndé hěn dà, dànshì Dù'ālā shì jīngjì shǒudū.}) « Yaoundé est grande, mais Douala est la capitale économique. »
\item \zh{石油很重要，但是农业更重要。} (\emph{Shíyóu hěn zhòngyào, dànshì nóngyè gèng zhòngyào.}) « Le pétrole est important, mais l'agriculture l'est davantage. »
\end{itemize}
\end{exemplebox}

\courssub{Le connecteur d'opinion : 我觉得 / 我认为}

\begin{propriete}[Donner son avis]
Structure : \cle{我觉得} / \cle{我认为} + phrase complète (sans \zh{que} !).\\
\zh{我觉得} (\emph{wǒ juéde}) « je pense, je trouve, j'ai l'impression » est le plus courant à l'oral et à l'écrit courant ; \zh{我认为} (\emph{wǒ rènwéi}) « je considère, j'estime » est plus soutenu, idéal pour conclure une rédaction formelle.
\begin{itemize}
\item \zh{我觉得喀麦隆的经济会越来越好。} (\emph{Wǒ juéde Kāmàilóng de jīngjì huì yuè lái yuè hǎo.}) « Je pense que l'économie du Cameroun ira de mieux en mieux. »
\item \zh{我认为农业是喀麦隆经济的未来。} (\emph{Wǒ rènwéi nóngyè shì Kāmàilóng jīngjì de wèilái.}) « Je considère que l'agriculture est l'avenir de l'économie camerounaise. »
\end{itemize}
Attention : après \zh{我觉得} ou \zh{我认为}, on met directement la phrase, sans mot de liaison : \zh{我觉得喀麦隆很美} et non « 我觉得\emph{que}喀麦隆很美 ».
\end{propriete}

\courssub{La comparaison d'action : A 比 B + verbe + 得 + adverbe}

\begin{propriete}[Comparer la manière dont deux sujets font une action]
Pour comparer la \textbf{manière} ou la \textbf{vitesse} dont deux sujets accomplissent une action :\\
\textbf{A + 比 + B + verbe + 得 + adverbe}\\
\zh{喀麦隆比泰国发展得慢。} (\emph{Kāmàilóng bǐ Tàiguó fāzhǎn de màn.}) « Le Cameroun se développe plus lentement que la Thaïlande. »\\
\zh{中国比喀麦隆发展得快。} (\emph{Zhōngguó bǐ Kāmàilóng fāzhǎn de kuài.}) « La Chine se développe plus vite que le Cameroun. »\\
\zh{杜阿拉比雅温得发展得快。} (\emph{Dù'ālā bǐ Yǎwēndé fāzhǎn de kuài.}) « Douala se développe plus vite que Yaoundé. »
\end{propriete}

\begin{popfigure}
\begin{center}
\begin{tikzpicture}
\node[draw=popgreen, fill=popgreenL, rounded corners, thick, minimum width=1.6cm, minimum height=0.9cm] (a) {A};
\node[draw=popgreen, fill=popgreenL, rounded corners, thick, minimum width=1.6cm, minimum height=0.9cm, right=0.5cm of a] (bi) {比 \emph{bǐ}};
\node[draw=popgreen, fill=popgreenL, rounded corners, thick, minimum width=1.6cm, minimum height=0.9cm, right=0.5cm of bi] (b) {B};
\node[draw=popblue, fill=popblueL, rounded corners, thick, minimum width=1.8cm, minimum height=0.9cm, right=0.5cm of b] (v) {verbe};
\node[draw=poporange, fill=poporangeL, rounded corners, thick, minimum width=1.6cm, minimum height=0.9cm, right=0.5cm of v] (de) {得 \emph{de}};
\node[draw=poppurple, fill=poppurpleL, rounded corners, thick, minimum width=2cm, minimum height=0.9cm, right=0.5cm of de] (adj) {adverbe};
\node[below=0.35cm of a, align=center] {\small 喀麦隆};
\node[below=0.35cm of bi, align=center] {\small 比};
\node[below=0.35cm of b, align=center] {\small 泰国};
\node[below=0.35cm of v, align=center] {\small 发展};
\node[below=0.35cm of de, align=center] {\small 得};
\node[below=0.35cm of adj, align=center] {\small 慢};
\draw[-{Stealth[length=2.5mm]}, thick, popdark] (a) -- (bi);
\draw[-{Stealth[length=2.5mm]}, thick, popdark] (bi) -- (b);
\draw[-{Stealth[length=2.5mm]}, thick, popdark] (b) -- (v);
\draw[-{Stealth[length=2.5mm]}, thick, popdark] (v) -- (de);
\draw[-{Stealth[length=2.5mm]}, thick, popdark] (de) -- (adj);
\end{tikzpicture}
\end{center}
\legende{Schéma de la comparaison d'action : A 比 B + V + 得 + Adj — « Le Cameroun se développe plus lentement que la Thaïlande. »}
\end{popfigure}

\begin{attention}[Comparaison d'action vs comparaison d'état]
Si l'on compare un \textbf{adjectif d'état} (taille, beauté, importance) et non une action, on n'utilise \textbf{pas} \zh{得} :\\
\zh{中国比喀麦隆大。} (\emph{Zhōngguó bǐ Kāmàilóng dà.}) « La Chine est plus grande que le Cameroun. » (Pas de \zh{得}).\\
\zh{喀麦隆比泰国热。} « Le Cameroun est plus chaud que la Thaïlande. »
\end{attention}

\courssub{Le complément de degré avec 得}

\begin{propriete}[Dire « comment » se fait l'action]
Structure : \textbf{verbe + 得 + adverbe}. Le mot-outil \cle{得} (\emph{de}, ton neutre) introduit un complément qui évalue la manière, le degré ou le résultat de l'action.
\begin{itemize}
\item \zh{发展得很快} (\emph{fāzhǎn de hěn kuài}) « se développer (très) vite »
\item \zh{发展得很慢} (\emph{fāzhǎn de hěn màn}) « se développer lentement »
\item \zh{他说汉语说得很好。} (\emph{Tā shuō Hànyǔ shuō de hěn hǎo.}) « Il parle très bien chinois. » (Règle : si le verbe a un objet \zh{汉语}, on répète le verbe \zh{说} avant \zh{得}).
\end{itemize}
\end{propriete}

\begin{attention}[Les trois \emph{de} : ne les confondez pas !]
C'est l'erreur n°1 des francophones.
\begin{itemize}
\item \zh{的} (\emph{de}) : \textbf{avant un nom} (possession, attribut). \zh{喀麦隆\underline{的}经济} (l'économie \underline{du} Cameroun).
\item \zh{得} (\emph{de}) : \textbf{après un verbe} (complément de degré). \zh{发展\underline{得}快} (se développer \underline{vite}).
\item \zh{地} (\emph{de}) : \textbf{avant un verbe} (complément de manière). \zh{慢慢\underline{地}发展} (se développer \underline{lentement}).
\end{itemize}
Astuce mnémotechnique : \zh{的} a \zh{白} (blanc) à droite $\rightarrow$ nom (chose concrète) ; \zh{得} a \zh{彳} (pas) en bas à gauche $\rightarrow$ action (verbe) ; \zh{地} a \zh{土} (terre) en bas à gauche $\rightarrow$ manière (sur le terrain).
\end{attention}

\begin{attention}[« Très » après 得 : un seul 很 !]
Après \zh{得}, l'adverbe \zh{很} ne signifie plus vraiment « très », il sert surtout de support rythmique obligatoire pour l'adjectif monosyllabique. On dit \zh{发展得很快}, jamais \zh{发展得很快很} ni \zh{发展得很很}.
\end{attention}

\courssub{La progression : 越来越 + adjectif}

\begin{propriete}[Exprimer « de plus en plus »]
Structure : \cle{越来越} (\emph{yuè lái yuè}) + adjectif (ou verbe d'état psychologique).\\
\zh{越来越好} (\emph{yuè lái yuè hǎo}) « de mieux en mieux » ; \zh{越来越快} « de plus en plus vite » ; \zh{越来越多} « de plus en plus nombreux » ; \zh{越来越重要} « de plus en plus important ».\\
\zh{喀麦隆的经济越来越好。} (\emph{Kāmàilóng de jīngjì yuè lái yuè hǎo.}) « L'économie du Cameroun va de mieux en mieux. »\\
\zh{学汉语的人越来越多。} (\emph{Xué Hànyǔ de rén yuè lái yuè duō.}) « Les gens qui apprennent le chinois sont de plus en plus nombreux. »
\end{propriete}

\begin{attention}[Pas de 很 après 越来越]
On ne dit jamais \zh{越来越很好} : après \zh{越来越}, l'adjectif seul suffit, \zh{很} est interdit car \zh{越来越} porte déjà l'idée de degré élevé.
\end{attention}

\courssub{Tableau récapitulatif des connecteurs et structures}

\begin{center}
\begin{tabularx}{\linewidth}{|X|X|X|X|}
\hline
\rowcolor{popblueL} Connecteur / Structure & Pinyin & Emploi & Exemple traduit \\
\hline
因为…所以… & yīnwèi… suǒyǐ… & cause $\rightarrow$ conséquence & 因为有很多石油，所以经济发展得快。« Parce qu'il y a beaucoup de pétrole, l'économie se développe vite. » \\
\hline
而且 & érqiě & addition de phrases (tête de prop. 2) & 有石油，而且有农业产品。« Il y a du pétrole, et de plus des produits agricoles. » \\
\hline
也 & yě & addition (avant le verbe) & 也有很多资源。« Il y a aussi beaucoup de ressources. » \\
\hline
但是 & dànshì & opposition / concession & 比泰国慢，但是越来越好。« Plus lent que la Thaïlande, mais de mieux en mieux. » \\
\hline
我觉得 / 我认为 & wǒ juéde / wǒ rènwéi & opinion (suivi phrase directe) & 我觉得经济会越来越好。« Je pense que l'économie ira de mieux en mieux. » \\
\hline
A 比 B + V + 得 + Adj & bǐ & comparaison d'action & 喀麦隆比泰国发展得慢。« Le Cameroun se développe plus lentement que la Thaïlande. » \\
\hline
越来越 + Adj & yuè lái yuè & progression & 越来越好。« De mieux en mieux. » \\
\hline
V + 得 + Adv & de & complément de degré & 发展得快。« Se développer vite. » \\
\hline
\end{tabularx}
\end{center}

\courssub{Méthode : construire un paragraphe complet}

\begin{methode}[Le plan en quatre phrases type]
Pour rédiger votre paragraphe sur l'économie du Cameroun (ou tout sujet similaire), suivez toujours cette architecture logique :
\begin{enumerate}
\item \textbf{Phrase d'introduction (Sujet + Ressources)} : présentez le sujet avec \zh{有} ou \zh{是}. 
Exemple : \zh{喀麦隆是一个有很多资源的国家。} (\emph{Kāmàilóng shì yí ge yǒu hěn duō zīyuán de guójiā.}) « Le Cameroun est un pays qui a beaucoup de ressources. »
\item \textbf{Deux phrases de causes} avec \zh{因为…所以…} : expliquez les moteurs économiques.
Exemple : \zh{因为喀麦隆有石油和农业产品，所以经济发展得越来越快。}
\item \textbf{Une nuance ou une comparaison} avec \zh{但是} ou \zh{比} : montrez l'esprit critique.
Exemple : \zh{喀麦隆经济比泰国发展得慢，但是越来越好。}
\item \textbf{Une opinion personnelle} avec \zh{我觉得} ou \zh{我认为} : concluez.
Exemple : \zh{我觉得喀麦隆的经济会越来越好。}
\end{enumerate}
\cle{Relecture obligatoire} : vérifiez que chaque verbe a son sujet, que les \zh{得} sont bien placés après les verbes d'action, que \zh{越来越} n'est pas suivi de \zh{很}, et que les tons du pinyin sont notés correctement.
\end{methode}

\begin{exoresolu}[Assembler les phrases avec le bon connecteur]
Énoncé : Reliez les deux phrases suivantes avec le connecteur qui convient, écrivez en caractères, pinyin et traduisez.\\
a) 喀麦隆有很多石油。／ 经济发展得很快。 (relation de cause)\\
b) 喀麦隆经济比泰国发展得慢。／ 经济越来越好。 (relation d'opposition)
\tcblower
Solution détaillée :\\
a) La deuxième phrase est la conséquence de la première $\rightarrow$ \zh{因为…所以…}.\\
\zh{因为喀麦隆有很多石油，所以经济发展得很快。}\\
\emph{Yīnwèi Kāmàilóng yǒu hěn duō shíyóu, suǒyǐ jīngjì fāzhǎn de hěn kuài.}\\
« Parce que le Cameroun a beaucoup de pétrole, son économie se développe vite. »\\
b) Les deux idées se contrastent (lent vs bien) $\rightarrow$ \zh{但是}.\\
\zh{喀麦隆经济比泰国发展得慢，但是经济越来越好。}\\
\emph{Kāmàilóng jīngjì bǐ Tàiguó fāzhǎn de màn, dànshì jīngjì yuè lái yuè hǎo.}\\
« L'économie du Cameroun se développe plus lentement que celle de la Thaïlande, mais elle va de mieux en mieux. »\\
Vérification : \zh{得} est bien placé après le verbe \zh{发展}, et \zh{越来越} est suivi de l'adjectif seul \zh{好}, sans \zh{很}.
\end{exoresolu}

\begin{exercice}[À vous de jouer : Thème (Français $\rightarrow$ Chinois)]
Traduisez en chinois, en utilisant obligatoirement les connecteurs et structures de la leçon. Écrivez les caractères, le pinyin (avec tons) et respectez la ponctuation chinoise.
\begin{enumerate}
\item Parce que Douala a un grand port, beaucoup d'entreprises s'y trouvent.
\item Le Cameroun a du pétrole, et de plus beaucoup de produits agricoles.
\item Je pense que l'économie du Cameroun ira de mieux en mieux.
\item La Chine se développe plus vite que le Cameroun, mais le Cameroun se développe de mieux en mieux.
\end{enumerate}
\end{exercice}

\begin{corrige}[Exercice « À vous de jouer »]
\begin{enumerate}
\item \zh{因为杜阿拉有一个大港口，所以很多公司在那里。}\\
\emph{Yīnwèi Dù'ālā yǒu yí ge dà gǎngkǒu, suǒyǐ hěn duō gōngsī zài nàlǐ.}
\item \zh{喀麦隆有石油，而且有很多农业产品。}\\
\emph{Kāmàilóng yǒu shíyóu, érqiě yǒu hěn duō nóngyè chǎnpǐn.}
\item \zh{我觉得喀麦隆的经济会越来越好。}\\
\emph{Wǒ juéde Kāmàilóng de jīngjì huì yuè lái yuè hǎo.}
\item \zh{中国比喀麦隆发展得快，但是喀麦隆发展得越来越好。}\\
\emph{Zhōngguó bǐ Kāmàilóng fāzhǎn de kuài, dànshì Kāmàilóng fāzhǎn de yuè lái yuè hǎo.}
\end{enumerate}
\end{corrige}

\begin{aretenir}[L'essentiel de la section]
\begin{itemize}
\item \zh{因为…所以…} s'utilise en entier dans une même phrase (contrairement au français « parce que… donc… » incorrect).
\item \zh{而且} relie des phrases (tête de proposition) ; \zh{也} se place avant le verbe ; \zh{但是} exprime l'opposition.
\item \zh{我觉得} / \zh{我认为} + phrase directe, sans « que » (\zh{ que} n'existe pas en chinois ici).
\item Comparaison d'action : A \zh{比} B + verbe + \zh{得} + adverbe. Comparaison d'état (adjectif seul) : A \zh{比} B + adjectif (sans \zh{得}).
\item Complément de degré : Verbe + \zh{得} + Adverbe. Attention aux trois \emph{de} : \zh{的} (nom), \zh{得} (verbe), \zh{地} (verbe, manière).
\item Après \zh{得}, un seul \zh{很} (support rythmique) ; après \zh{越来越}, adjectif seul (pas de \zh{很}).
\item Structure type d'un paragraphe argumentatif simple : Introduction (\zh{有} / \zh{est}) $\rightarrow$ Causes (\zh{因为…所以}) $\rightarrow$ Nuance/Comparaison (\zh{但是}/\zh{比}) $\rightarrow$ Opinion (\zh{我觉得}).
\end{itemize}
\end{aretenir}

\coursec{Modèle de texte commenté}

Nous commençons par l'étude d'un texte modèle original, synthétisant le vocabulaire et les structures cibles de la leçon. Ce texte répond à la question du programme : \zh{为什么喀麦隆经济比泰发展？} (Pourquoi l'économie du Cameroun se développe-t-elle [plus que la Thaïlande] ?). Nous l'analysons phrase par phrase.

\begin{textebox}[Texte modèle : Pourquoi l'économie du Cameroun se développe-t-elle ?]
\zh{喀麦隆是非洲中部的一个国家。} \\
\zh{这个国家有很多自然资源。} \\
\zh{例如，喀麦隆出口石油、可可、咖啡和木材。} \\
\zh{因为石油卖得很好，所以国家有钱修路、建学校和建医院。} \\
\zh{杜阿拉是最大的港口，对贸易非常重要。} \\
\zh{我觉得喀麦隆的经济越来越好。} \\
\zh{以后，如果农业和工业更发展，喀麦隆会更强。} \\

\textbf{Pinyin :} \\
Kāmàilóng shì Fēizhōngbù de yí gè guójiā. \\
Zhè gè guójiā yǒu hěn duō zìrán zīyuán. \\
Lìrú, Kāmàilóng chūkǒu shíyóu, kěkě, kāfēi hé múcái. \\
Yīnwèi shíyóu mài de hěn hǎo, suǒyǐ guójiā yǒu qián xiū lù, jiàn xuéxiào hé jiàn yīyuàn. \\
Dù'ālā shì zuì dà de gǎngkǒu, duì màoyì fēicháng zhòngyào. \\
Wǒ juéde Kāmàilóng de jīngjì yuè lái yuè hǎo. \\
Yǐhòu, rúguǒ nóngyè hé gōngyè gèng fāzhǎn, Kāmàilóng huì gèng qiáng.

\textbf{Traduction française :}
Le Cameroun est un pays d'Afrique centrale. \\
Ce pays a beaucoup de ressources naturelles. \\
Par exemple, le Cameroun exporte du pétrole, du cacao, du café et du bois. \\
Parce que le pétrole se vend très bien, donc le pays a de l'argent pour construire des routes, bâtir des écoles et des hôpitaux. \\
Douala est le plus grand port, il est très important pour le commerce. \\
Je pense que l'économie du Cameroun va de mieux en mieux. \\
À l'avenir, si l'agriculture et l'industrie se développent davantage, le Cameroun sera plus fort.
\end{textebox}

\begin{exemplebox}[Analyse structurelle du texte modèle]
\begin{enumerate}
\item \zh{喀麦隆是非洲中部的一个国家。} \textbf{Structure :} S (\zh{喀麦隆}) + \zh{是} + Nom prédicatif (\zh{非洲中部的一个国家}). \emph{Note :} \zh{中部} (centre) localise \zh{非洲}.
\item \zh{这个国家有很多自然资源。} \textbf{Structure :} S (\zh{这个国家}) + \zh{有} + COD (\zh{很多自然资源}). \zh{很多} = beaucoup de (quantifieur).
\item \zh{例如，喀麦隆出口石油、可可、咖啡和木材。} \textbf{Connecteur :} \zh{例如} (par exemple) en tête. \textbf{Verbe :} \zh{出口} (exporter, vt). \textbf{Énumération :} \zh{、} (séparateur liste), \zh{和} (et) avant dernier élément.
\item \zh{因为石油卖得很好，所以国家有钱修路、建学校和建医院。} \textbf{Causalité :} \zh{因为}... \zh{所以}... \textbf{Degré :} \zh{卖得很好} (V + \zh{得} + Adj). \textbf{Série verbale :} \zh{有钱} (avoir de l'argent) + but \zh{修路、建学校...} (construire routes, écoles...). \emph{Distinction :} \zh{修} (routes/réparer) vs \zh{建} (bâtiments).
\item \zh{杜阿拉是最大的港口，对贸易非常重要。} \textbf{Superlatif :} \zh{最} + Adj (\zh{大}) + \zh{的} + Nom. \textbf{Prédication adjectivale :} Sujet \zh{杜阿拉} (topic) + \zh{对贸易非常重要} (commentaire).
\item \zh{我觉得喀麦隆的经济越来越好。} \textbf{Verbe mental :} \zh{觉得} + Proposition. \textbf{Évolution :} \zh{越来越} + Adj (\zh{好}).
\item \zh{以后，如果农业和工业更发展，喀麦隆会更强。} \textbf{Temps :} \zh{以后} (à l'avenir) en tête. \textbf{Condition :} \zh{如果}... (pas de \zh{则} nécessaire ici). \textbf{Comparatif implicite :} \zh{更发展}, \zh{更强} (plus développé, plus fort qu'actuellement). \textbf{Modal :} \zh{会} (futur/certitude).
\end{enumerate}
\end{exemplebox}

\coursec{Vocabulaire et glossaire du thème}

Le tableau ci-dessous regroupe le lexique essentiel de la leçon (niveau HSK 3-4), classé par nature grammaticale. Maîtrisez l'écriture, le pinyin (tons) et le sens.

\begin{center}
\begin{tabularx}{\linewidth}{|X|X|X|X|}
\hline \rowcolor{popblueL} \textbf{Caractères} & \textbf{Pinyin} & \textbf{Nature} & \textbf{Traduction} \\ \hline
经济 & jīngjì & Nom & Économie \\ \hline
发展 & fāzhǎn & Verbe & Se développer, développer \\ \hline
资源 & zīyuán & Nom & Ressource(s) \\ \hline
自然资源 & zìrán zīyuán & Nom & Ressources naturelles \\ \hline
石油 & shíyóu & Nom & Pétrole \\ \hline
可可 & kěkě & Nom & Cacao \\ \hline
咖啡 & kāfēi & Nom & Café \\ \hline
木材 & múcái & Nom & Bois, bois d'œuvre \\ \hline
出口 & chūkǒu & Verbe / Nom & Exporter / Exportation(s) \\ \hline
进口 & jìnkǒu & Verbe / Nom & Importer / Importation(s) \\ \hline
贸易 & màoyì & Nom & Commerce, échanges commerciaux \\ \hline
港口 & gǎngkǒu & Nom & Port \\ \hline
基础设施 & jīchǔ shèshī & Nom & Infrastructures \\ \hline
农业 & nóngyè & Nom & Agriculture \\ \hline
工业 & gōngyè & Nom & Industrie \\ \hline
旅游业 & lǚyóuyè & Nom & Tourisme \\ \hline
重要 & zhòngyào & Adj (V.état) & Important \\ \hline
发达 & fādá & Adj (V.état) & Développé (économie, transport, technologie) \\ \hline
强 & qiáng & Adj (V.état) & Fort, puissant \\ \hline
快 & kuài & Adj (V.état) & Rapide, vite \\ \hline
慢 & màn & Adj (V.état) & Lent \\ \hline
卖 & mài & Verbe & Vendre \\ \hline
买 & mǎi & Verbe & Acheter \\ \hline
修 & xiū & Verbe & Construire (route), réparer \\ \hline
建 & jiàn & Verbe & Construire, bâtir (bâtiment) \\ \hline
因为 & yīnwèi & Conjonction & Parce que \\ \hline
所以 & suǒyǐ & Conjonction & Donc, par conséquent \\ \hline
如果 & rúguǒ & Conjonction & Si \\ \hline
例如 & lìrú & Conjonction & Par exemple \\ \hline
越来越 & yuè lái yuè & Adverbe & De plus en plus \\ \hline
更 & gèng & Adverbe & Plus, davantage (comparatif implicite) \\ \hline
最 & zuì & Adverbe & Le plus (superlatif) \\ \hline
以后 & yǐhòu & Nom temps / Adv & À l'avenir, après, plus tard \\ \hline
觉得 & juéde & Verbe & Penser, trouver (avis subjectif) \\ \hline
认为 & rènwéi & Verbe & Considérer, estimer (plus formel) \\ \hline
\end{tabularx}
\end{center}

\coursec{Écriture correcte des caractères}

Cette section cible l'écriture normée (ordre des traits, radicaux, structure) des caractères clés du thème. Respectez les règles : \textbf{haut $\rightarrow$ bas, gauche $\rightarrow$ droite, horizontal avant vertical, trait central avant ailes, enveloppe avant contenu, fermer en dernier}.

\begin{propriete}[Analyse graphémique : caractères à maîtriser]
\begin{enumerate}
\item \zh{经} (jīng) : Radical \zh{纟} (sī, soie, 3 traits). Structure gauche-droite. Ordre : \zh{纟} (3) $\rightarrow$ \zh{工} (3) $\rightarrow$ \zh{土} (3) = 8 traits. \emph{Sens :} trame, passer par, gérer $\rightarrow$ économie (\zh{经济}).
\item \zh{济} (jì) : Radical \zh{氵} (shuǐ, eau, 3 traits). Structure gauche-droite. Phonétique \zh{齐} (qí). Ordre : \zh{氵} $\rightarrow$ \zh{齐} (丿 $\rightarrow$ 亠 $\rightarrow$ 八 $\rightarrow$ 米) = 13 traits. \emph{Sens :} traverser l'eau, secourir $\rightarrow$ économie, finances.
\item \zh{资} (zī) : Radical \zh{亻} (rén, homme, 2 traits). Structure gauche-droite. Phonétique \zh{次} (cì). Ordre : \zh{亻} $\rightarrow$ \zh{二} $\rightarrow$ \zh{欠} = 10 traits. \emph{Sens :} capital, ressources.
\item \zh{源} (yuán) : Radical \zh{氵} (eau). Structure gauche-droite. Phonétique \zh{原} (yuán). Ordre : \zh{氵} $\rightarrow$ \zh{厂} $\rightarrow$ \zh{白} $\rightarrow$ \zh{小} $\rightarrow$ \zh{水} = 13 traits. \emph{Sens :} source, origine.
\item \zh{港} (gǎng) : Radical \zh{氵}. Structure gauche-droite. Phonétique \zh{共} (gòng) + \zh{口} (kǒu). Ordre : \zh{氵} $\rightarrow$ \zh{共} (一 $\rightarrow$ 口 $\rightarrow$ 八) $\rightarrow$ \zh{口} = 12 traits. \emph{Sens :} port (lieu d'eau commun).
\item \zh{贸} (mào) : Radical \zh{贝} (bèi, coquillage/monnaie, 4 traits). Structure haut-bas. Haut : \zh{戊} (wù). Ordre : \zh{戊} (5) $\rightarrow$ \zh{贝} (4) = 9 traits. \emph{Sens :} commerce (échange de biens/monnaie).
\item \zh{修} (xiū) : Radical \zh{亻} (homme). Structure gauche-droite. Droite : \zh{彡} (shān, ornement/cheveux, 3 traits) + \zh{卩} (jié, sceau, 2 traits). Ordre : \zh{亻} $\rightarrow$ \zh{彡} $\rightarrow$ \zh{卩} = 9 traits. \emph{Sens :} cultiver, réparer, construire (route).
\item \zh{建} (jiàn) : Radical \zh{廴} (yǐn, marcher/bâtiment, 3 traits). Structure bas-gauche (enveloppe). Dedans : \zh{笔} (bǐ, pinceau) simplifié en \zh{聿}? Non, \zh{建} = \zh{廴} + \zh{笔} (traditionnel) $\rightarrow$ simplifié \zh{廴} + \zh{笔} (partie sup) + \zh{廴}? Simplifié : \zh{廴} (3) + \zh{笔} (10) = 13 traits? Non, \zh{建} simplifié 9 traits. Décomposition : \zh{廴} + \zh{笔} (partie haute \zh{竹} simplifiée \zh{⺮} + \zh{毛}). \emph{Ordre pratique :} haut (\zh{竹}/\zh{聿}) $\rightarrow$ \zh{廴}. \emph{Sens :} établir, construire (bâtiment).
\item \zh{强} (qiáng) : Radical \zh{弓} (gōng, arc, 3 traits). Structure haut-bas. Milieu : \zh{虫} (chóng, insecte, 6 traits). Bas : \zh{刂} (dāo, couteau, 2 traits). Ordre : \zh{弓} $\rightarrow$ \zh{虫} $\rightarrow$ \zh{刂} = 11 traits. \emph{Sens :} fort, puissant (arc + insecte + force).
\end{enumerate}
\end{propriete}

\begin{attention}[Caractères proches / Pièges visuels]
\begin{itemize}
\item \zh{经} (jīng, économie) vs \zh{径} (jìng, diamètre/chemin) : même phonétique \zh{巠}, radical différent (soie vs fil/silk).
\item \zh{济} (jì, économie/aider) vs \zh{祭} (jì, sacrifier) : même phonétique \zh{齐}, radical eau vs montrer/viande.
\item \zh{修} (xiū, réparer) vs \zh{徐} (xú, lent) : même composant droit, radical homme vs marche (彳).
\item \zh{建} (jiàn, construire) vs \zh{健} (jiàn, en bonne santé) : même phonétique \zh{建}, radical enveloppe vs personne (\zh{亻}).
\item \zh{资} (zī, ressource) vs \zh{咨} (zī, consulter) : même phonétique \zh{次}, radical homme vs bouche.
\end{itemize}
\end{attention}

\begin{experience}[Atelier écriture : Dictée guidée \& Mémorisation]
\textbf{Consigne :} L'enseignant dicte les mots suivants (pinyin $\rightarrow$ caractères). Les élèves écrivent sur grille \zh{田字格} en respectant l'ordre des traits.
\begin{enumerate}
\item \zh{jīngjì} (économie) \item \zh{zīyuán} (ressources) \item \zh{chūkǒu} (exporter) \item \zh{gǎngkǒu} (port)
\item \zh{jiànshè} (construction/bâtir -- \zh{建设}) \item \zh{qiáng} (fort) \item \zh{fāzhǎn} (développer)
\end{enumerate}
\textbf{Auto-correction :} Vérifier la proportion des composants (ex: \zh{纟} étroit à gauche), la fermeture de \zh{廴}, la symétrie de \zh{弓}.
\end{experience}

\coursec{Structures et connecteurs pour écrire}

Cette section synthétise les outils grammaticaux indispensables pour produire un texte économique cohérent (niveau HSK 3-4). Chaque structure est illustrée d'exemples contextualisés Cameroun/Chine.

\begin{propriete}[1. Causalité : \textbf{因为} (yīnwèi) ... \textbf{所以} (suǒyǐ) ...]
\zh{Structure :} \zh{因为} + Cause (Proposition 1) ， \zh{所以} + Conséquence (Proposition 2).
\emph{Règle :} \zh{所以} est \textbf{obligatoire} à l'écrit formel pour marquer la conséquence logique.
\begin{center}
\begin{tabularx}{\linewidth}{|X|X|X|}
\hline \rowcolor{popblueL} Chinois & Pinyin & Français \\ \hline
\zh{因为喀麦隆有石油，所以经济发展得快。} & Yīnwèi Kāmàilóng yǒu shíyóu, suǒyǐ jīngjì fāzhǎn de kuài. & Parce que le Cameroun a du pétrole, l'économie se développe vite. \\ \hline
\zh{因为基础设施不好，所以运输很慢。} & Yīnwèi jīchǔ shèshī bù hǎo, suǒyǐ yùnshū hěn màn. & Comme les infrastructures sont mauvaises, le transport est lent. \\ \hline
\end{tabularx}
\end{center}
\emph{Variante orale :} \zh{因为}... (sans \zh{所以}) ou \zh{所以}... \zh{因为}... (conséquence d'abord).
\end{propriete}

\begin{propriete}[2. Comparatif explicite : \textbf{比} (bǐ)]
\zh{Structure :} A + \zh{比} + B + Adj (sans \zh{很}, \zh{非常} etc. directement après l'adj si non modifié).
\emph{Négation :} A + \zh{不比} + B + Adj (A n'est pas plus... que B / A est aussi... que B).
\emph{Degré :} A + \zh{比} + B + \zh{得多} / \zh{多了} / \zh{一点儿} + Adj.
\begin{center}
\begin{tabularx}{\linewidth}{|X|X|X|}
\hline \rowcolor{popblueL} Chinois & Pinyin & Français \\ \hline
\zh{喀麦隆的经济比泰国发展得快。} & Kāmàilóng de jīngjì bǐ Tàiguó fāzhǎn de kuài. & L'éco du Cameroun se développe plus vite que la Thaïlande. \\ \hline
\zh{杜阿拉的港口比雅温得大得多。} & Dù'ālā de gǎngkǒu bǐ Yǎwēndé dà de duō. & Le port de Douala est bien plus grand que Yaoundé. \\ \hline
\zh{中国的工业不比喀麦隆落后。} & Zhōngguó de gōngyè bù bǐ Kāmàilóng luòhòu. & L'industrie chinoise n'est pas moins avancée que le Cameroun. \\ \hline
\end{tabularx}
\end{center}
\end{propriete}

\begin{propriete}[3. Superlatif : \textbf{最} (zuì)]
\zh{Structure :} Sujet + \zh{最} + Adj (+ \zh{的} + Nom).
\begin{center}
\begin{tabularx}{\linewidth}{|X|X|X|}
\hline \rowcolor{popblueL} Chinois & Pinyin & Français \\ \hline
\zh{杜阿拉是最大的港口。} & Dù'ālā shì zuì dà de gǎngkǒu. & Douala est le plus grand port. \\ \hline
\zh{石油是最重要的出口产品。} & Shíyóu shì zuì zhòngyào de chūkǒu chǎnpǐn. & Le pétrole est le produit d'exportation le plus important. \\ \hline
\end{tabularx}
\end{center}
\end{propriete}

\begin{propriete}[4. Comparatif implicite / Progression : \textbf{更} (gèng) \& \textbf{越来越} (yuè lái yuè)]
\begin{itemize}
\item \zh{更} (plus, davantage) : compare l'état actuel à un état futur/antérieur implicite. \zh{Structure :} S + \zh{更} + Adj/V.
\zh{Ex:} \zh{以后经济会更好。} (L'éco ira mieux à l'avenir.)
\item \zh{越来越} (de plus en plus) : processus progressif. \zh{Structure :} S + \zh{越来越} + Adj/V.
\zh{Ex:} \zh{喀麦隆的贸易越来越发达。} (Le commerce du Cameroun est de plus en plus développé.)
\end{itemize}
\end{propriete}

\begin{propriete}[5. Complément de degré : \textbf{V + 得 + Adj}]
\zh{Structure :} Verbe + \zh{得} + Adjectif (souvent précédé de \zh{很}, \zh{真}, \zh{非常}).
\emph{Règle d'or :} Verbe $\rightarrow$ \zh{得}. Nom/Adjectif attributif $\rightarrow$ \zh{的}.
\begin{center}
\begin{tabularx}{\linewidth}{|X|X|X|}
\hline \rowcolor{popblueL} Chinois & Structure & Français \\ \hline
\zh{石油卖得很好。} & V (\zh{卖}) + \zh{得} + Adv+Adj & Le pétrole se vend très bien. \\ \hline
\zh{经济发展得很快。} & V (\zh{发展}) + \zh{得} + Adv+Adj & L'économie se développe très vite. \\ \hline
\zh{路修得很好。} & V (\zh{修}) + \zh{得} + Adv+Adj & La route est bien construite. \\ \hline
\end{tabularx}
\end{center}
\end{propriete}

\begin{propriete}[6. Les trois particules « de » : déterminant, degré, manière]
\begin{center}
\begin{tabularx}{\linewidth}{|X|X|X|X|}
\hline \rowcolor{popblueL} Particule & Position & Fonction & Exemple \\ \hline
\zh{的} & Dét + \zh{的} + Nom & Possession, attribution, relativisation & \zh{喀麦隆的经济}, \zh{非常重要的港口} \\ \hline
\zh{得} & V + \zh{得} + Adj & Complément de degré (résultat/manière) & \zh{发展得快}, \zh{卖得好} \\ \hline
\zh{地} & Adv + \zh{地} + V & Complément de manière (adverbe en -ment) & \zh{认真地学习} (étudier sérieusement) \\ \hline
\end{tabularx}
\end{center}
\end{propriete}

\begin{attention}[Piège majeur] \zh{发展得快} (se développe vite) vs \zh{快速发展} (se développer rapidement -- V composé). \zh{发展快} (sans \zh{得}) = phrase nominale "développement rapide" ou oral très familier. \textbf{À l'écrit : toujours \zh{得}.}
\end{attention}


\begin{propriete}[7. Temps futur / Modal : \textbf{会} (huì) + \textbf{以后} (yǐhòu)]
\zh{Structure :} (Sujet) + \zh{以后} + \zh{会} + Verbe/Adj.
\emph{Place de \zh{以后} :} En tête de phrase (Thème) ou après Sujet. \textbf{Jamais en fin de phrase.}
\begin{center}
\begin{tabularx}{\linewidth}{|X|X|X|}
\hline \rowcolor{popblueL} Chinois & Structure & Français \\ \hline
\zh{以后喀麦隆会更强。} & Temps (Thème) + S + Modal + Adj & À l'avenir, le Cameroun sera plus fort. \\ \hline
\zh{喀麦隆以后会更强。} & S + Temps + Modal + Adj & Le Cameroun sera plus fort à l'avenir. \\ \hline
\end{tabularx}
\end{center}
\end{propriete}

\begin{propriete}[8. Condition : \textbf{如果} (rúguǒ) ... (就 \zh{jiù}) ...]
\zh{Structure :} \zh{如果} + Condition, (Sujet +) \zh{就} + Conséquence. \zh{就} marque la conséquence logique, souvent omis si \zh{会} présent.
\zh{Ex:} \zh{如果农业更发展，喀麦隆就更强。} / \zh{如果农业更发展，喀麦隆会更强。}
\end{propriete}

\begin{aretenir}[Résumé des connecteurs logiques pour l'expression écrite]
\begin{center}
\begin{tabularx}{\linewidth}{|X|X|X|}
\hline \rowcolor{popblueL} Fonction & Connecteurs chinois & Français \\ \hline
Énumération & \zh{和}、\zh{与}、\zh{跟}、\zh{、} (liste) & Et \\ \hline
Exemplification & \zh{例如} (lìrú), \zh{比如} (bǐrú) & Par exemple, comme \\ \hline
Cause & \zh{因为} (yīnwèi), \zh{由于} (yóuyú - plus formel) & Parce que, en raison de \\ \hline
Conséquence & \zh{所以} (suǒyǐ), \zh{因此} (yīncǐ - formel) & Donc, par conséquent \\ \hline
Condition & \zh{如果} (rúguǒ), \zh{要是} (yàoshi - oral) & Si \\ \hline
Concession & \zh{虽然} (suīrán) ... \zh{但是} (dànshì) ... & Bien que ... mais ... \\ \hline
But & \zh{为了} (wèile) + Verbe, \zh{以便} (yǐbiàn) & Pour, afin de \\ \hline
Temps (futur) & \zh{以后} (yǐhòu), \zh{将来} (jiānglái) & À l'avenir, plus tard \\ \hline
Progression & \zh{越来越} (yuè lái yuè), \zh{日益} (rìyì - écrit) & De plus en plus \\ \hline
Comparaison & \zh{比} (bǐ), \zh{更} (gèng), \zh{最} (zuì) & Plus que, plus, le plus \\ \hline
\end{tabularx}
\end{center}
\end{aretenir}

\coursec{Thème et version guidés}

Application concrète des structures ci-dessus. Nous détaillons le raisonnement pas à pas pour trois phrases types du programme.

\begin{methode}[Méthode en 4 étapes : du français au chinois (Thème)]
\textbf{Étape 1 : Analyser la phrase française.} Identifier le Sujet (S), le Verbe (V), les Compléments (COD, COI, Circonstanciels) et les connecteurs logiques.
\textbf{Étape 2 : Traduire les blocs lexicaux.} Chercher le vocabulaire chinois exact (nature grammaticale : nom, verbe, adjectif statif). Attention aux classificateurs.
\textbf{Étape 3 : Assembler selon l'ordre chinois canonique.} \textbf{S + (Temps/Adv/Connecteur) + V + (Particule) + C}. Compléments lieu/temps \textbf{AVANT} le verbe. \zh{的} (dét.), \zh{得} (degré), \zh{了} (aspect).
\textbf{Étape 4 : Vérifier et relire.} Tons, ordre des mots, classificateurs (Nombre + Classif + Nom), cohérence connecteurs (\zh{因为}...\zh{所以}).
\end{methode}

\begin{popfigure}
\begin{tikzpicture}[
    node distance=1.2cm and 1.5cm,
    box/.style={rectangle, rounded corners, draw=popblue, thick, fill=popblueL, text width=3.2cm, align=center, font=\small},
    arrow/.style={-Stealth, thick, popdark}
]
\node[box, align=center] (lire){1. Lire \& Analyser\\ (S, V, C, Connecteurs)};
\node[box, below=of lire, align=center] (blocs){2. Traduire les\\ blocs lexicaux};
\node[box, below=of blocs, align=center] (assembler){3. Assembler selon\\ l'ordre S-V-C chinois};
\node[box, below=of assembler, align=center] (verifier){4. Vérifier\\ (Tons, 的/得/了, Ordre)};
\draw[arrow] (lire) -- (blocs);
\draw[arrow] (blocs) -- (assembler);
\draw[arrow] (assembler) -- (verifier);
\draw[arrow, poporange] (verifier.east) -- ++(1.5,0) |- (lire.east) node[midway, right, font=\tiny, text=poporange, align=center] {Si erreur :\\ revenir à l'étape 1};
\end{tikzpicture}
\legende{Organigramme de la méthode de traduction (Thème)}
\end{popfigure}

\begin{methode}[Méthode en 4 étapes : du chinois au français (Version)]
\textbf{Étape 1 : Segmenter la phrase chinoise.} Repérer les frontières : Sujet (position initiale, avant \zh{是}/\zh{有}), Verbe (non conjugué), Particules (\zh{的}/\zh{得}/\zh{了}/\zh{过}/\zh{着}), Compléments.
\textbf{Étape 2 : Traduire les unités de sens.} Mot à mot puis regroupement. Attention aux mots « vides » (\zh{会}, \zh{能}, \zh{要}, \zh{以后}, \zh{越来越}).
\textbf{Étape 3 : Restructurer en français correct.} Conjugaison (temps, mode), accords, ordre S-V-C, prépositions (\zh{的}$\rightarrow$de/possessif, \zh{在}$\rightarrow$à/dans).
\textbf{Étape 4 : Polir le style.} Naturel : \zh{经济越来越好} $\rightarrow$ « L'économie va de mieux en mieux » (pas « est de plus en plus bonne »).
\end{methode}

\courssub{Thème guidé : Français $\rightarrow$ Chinois}

\begin{propriete}[Rappel : La particule \textbf{的} (possession / détermination / relativisation)]
\zh{Structure :} Déterminant + \zh{的} + Nom noyau.
\zh{Ex:} \zh{喀麦隆} + \zh{的} + \zh{经济} = \zh{喀麦隆的经济} (L'économie \textbf{du} Cameroun).
\zh{Ex:} \zh{非常重要} + \zh{的} + \zh{港口} = \zh{非常重要的港口} (Un port très important).
\end{propriete}

\begin{exemplebox}[Thème 1 : Phrase simple d'identification / classification]
\textbf{Français :} « Le Cameroun est un pays d'Afrique. »
\begin{enumerate}
\item \textbf{Analyse :} S = Le Cameroun ; V = est (identité/classification) ; Atribut = un pays d'Afrique.
\item \textbf{Blocs lexicaux :} Cameroun = \zh{喀麦隆} ; est = \zh{是} ; pays = \zh{国家} ; Afrique = \zh{非洲} ; central = \zh{中部} ; un = \zh{一} + classif. \zh{个} (omis en énoncé générique).
\item \textbf{Assemblage :} S (\zh{喀麦隆}) + \zh{是} + Atribut (\zh{非洲中部的国家}). \emph{Ordre déterminants :} Grand $\rightarrow$ Petit (\zh{非洲} $\rightarrow$ \zh{中部} $\rightarrow$ \zh{的} $\rightarrow$ \zh{国家}).
\item \textbf{Résultat :} \zh{喀麦隆是非洲中部的国家。}
\end{enumerate}
\textbf{Pinyin :} Kāmàilóng shì Fēizhōngbù de guójiā.
\end{exemplebox}

\begin{attention}[Piège : L'article indéfini « un » / Classificateur]
En chinois, \zh{一个} individualise/compte. Dans une définition générique (\emph{Le Cameroun est un pays...}), on omet \zh{一个}. \zh{喀麦隆是一个非洲国家} met l'accent sur « \textbf{un} (seul) pays ». La forme sans \zh{一个} est la norme pour les vérités générales.
\end{attention}

\begin{exemplebox}[Thème 2 : Verbe d'action transitif + COD (Énumération)]
\textbf{Français :} « Le Cameroun exporte du cacao, du café et du bois. »
\begin{enumerate}
\item \textbf{Analyse :} S = Le Cameroun ; V = exporte ; COD = cacao, café, bois (objets indéfinis, quantité non spécifiée).
\item \textbf{Blocs lexicaux :} Exporter = \zh{出口} (vt) ; Cacao = \zh{可可} ; Café = \zh{咖啡} ; Bois = \zh{木材} ; Et = \zh{和} / \zh{与} (écrit).
\item \textbf{Assemblage :} S (\zh{喀麦隆}) + V (\zh{出口}) + COD (\zh{可可、咖啡和木材}). Pas de \zh{了} (fait habituel/gnomique). Pas de \zh{的} après verbe.
\item \textbf{Résultat :} \zh{喀麦隆出口可可、咖啡和木材。}
\end{enumerate}
\textbf{Pinyin :} Kāmàilóng chūkǒu kěkě, kāfēi hé múcái.
\end{exemplebox}

\begin{exemplebox}[Thème 3 : Phrase complexe causale + Complément de degré + Série verbale]
\textbf{Français :} « Parce que le pétrole se vend bien, le pays a de l'argent pour construire des routes et des écoles. »
\begin{enumerate}
\item \textbf{Analyse :} Prop 1 (Cause) : S1=Le pétrole, V1=se vend, Degré=bien. Prop 2 (Conséquence) : S2=Le pays, V2=a, COD2=de l'argent, But=construire routes/écoles.
\item \textbf{Bloc Cause :} \zh{因为} + S1(\zh{石油}) + V1(\zh{卖}) + \zh{得} + Adj(\zh{好}) $\rightarrow$ \zh{因为石油卖得好} (ou \zh{卖得很好}).
\item \textbf{Bloc Conséquence :} \zh{所以} + S2(\zh{国家}) + V2(\zh{有}) + COD(\zh{钱}) + But(\zh{修路、建学校}).
\item \textbf{Assemblage :} \zh{因为石油卖得很好，所以国家有钱修路、建学校。}
\end{enumerate}
\textbf{Pinyin :} Yīnwèi shíyóu mài de hěn hǎo, suǒyǐ guójiā yǒu qián xiū lù, jiàn xuéxiào.
\end{exemplebox}

\begin{attention}[Erreurs fréquentes francophones sur Thème 3]
\begin{itemize}
\item \textbf{Oublier \zh{所以} :} *\zh{因为石油卖得好，国家有钱...} $\rightarrow$ Incorrect à l'écrit formel.
\item \textbf{Confusion \zh{得}/\zh{的} :} *\zh{卖的好}, *\zh{修的路} (si verbe). \textbf{Règle :} V + \zh{得} + Adj ; V + \zh{了} + \zh{的} + Nom (aspect accompli + déterminant).
\item \textbf{Verbes « construire » :} \zh{修} (xiū) pour routes (\zh{修路}), rails ; \zh{建} (jiàn) pour bâtiments (\zh{建学校}, \zh{建医院}, \zh{建厂}). Ne pas confondre.
\item \textbf{« Se vend bien » :} \zh{卖得好} (V + \zh{得} + Adj). Pas \zh{很好卖} (adj + V -- structure potentielle/résultative différente).
\end{itemize}
\end{attention}

\begin{exoresolu}[Thème supplémentaire : « Douala est un port très important pour l'économie. »]
\textbf{Énoncé :} Traduire en chinois (caractères + pinyin).
\textbf{Solution détaillée :}
\begin{enumerate}
\item \textbf{Analyse :} S = Douala ; V = est (classification) ; Atribut = un port très important ; Compl. finalité = pour l'économie.
\item \textbf{Vocabulaire :} Douala = \zh{杜阿拉} ; Port = \zh{港口} ; Important = \zh{重要} ; Très = \zh{非常} / \zh{很} ; Pour = \zh{对}... \zh{来说} / \zh{对于}... ; Économie = \zh{经济}.
\item \textbf{Structure cible :} Relativisation du complément « pour l'économie » avant le nom « port ».
\zh{Structure :} S + \zh{是} + [ \zh{对} + Éco + \zh{非常重要} + \zh{的} ] + \zh{港口}.
\item \textbf{Proposition 1 (Standard, élégante) :} \zh{杜阿拉是一个对经济非常重要的港口。}
\item \textbf{Proposition 2 (Focus importance) :} \zh{杜阿拉的港口对经济非常重要。} (Le port de Douala est très important pour l'économie.)
\item \textbf{Choix pédagogique :} Proposition 1 (intègre la relative \zh{对...重要的}).
\end{enumerate}
\textbf{Réponse attendue :} \zh{杜阿拉是一个对经济非常重要的港口。} (Dù'ālā shì yí gè duì jīngjì fēicháng zhòngyào de gǎngkǒu.)
\end{exoresolu}

\courssub{Version guidée : Chinois $\rightarrow$ Français}

\begin{exemplebox}[Version 1 : Structure \zh{越来越} + Adjectif verbal]
\textbf{Chinois :} \zh{喀麦隆的经济越来越好。}
\textbf{Pinyin :} Kāmàilóng de jīngjì yuè lái yuè hǎo.
\begin{enumerate}
\item \textbf{Segmentation :} \zh{喀麦隆的经济} (Sujet) | \zh{越来越} (Adv degré) | \zh{好} (V.état).
\item \textbf{Littéral :} Cameroun de économie de-plus-en-plus bon.
\item \textbf{Restructuration :} Sujet « L'économie du Cameroun ». \zh{好} = aller bien / être bon. \zh{越来越好} = évolution positive $\rightarrow$ « va de mieux en mieux » / « s'améliore de plus en plus ».
\item \textbf{Résultat :} \textbf{« L'économie du Cameroun va de mieux en mieux. »} (Ou : « ... s'améliore de plus en plus. »)
\end{enumerate}
\end{exemplebox}

\begin{propriete}[Structure \textbf{越来越} + Adj / Verbe mental]
\zh{越来越} (yuè lái yuè) = « de plus en plus ». Place : \textbf{devant} l'adjectif ou le verbe d'état.
\zh{Ex:} \zh{天气越来越冷。} (Il fait de plus en plus froid.)
\zh{Ex:} \zh{我越来越喜欢中文。} (J'aime de plus en plus le chinois.)
\end{propriete}

\begin{attention}[Piège version : \zh{好} = « bon » ou « bien » ?]
\zh{好} (hǎo) est un verbe adjectival (état).
\begin{itemize}
\item \zh{经济好} = L'économie est bonne / va bien.
\item \zh{越来越好} = « De mieux en mieux » (évolution) et non « de plus en plus bon » (maladroit pour une économie).
\item \textbf{Astuce :} États/évolutions $\rightarrow$ « aller bien/mieux », « s'améliorer », « se porter bien ».
\end{itemize}
\end{attention}

\begin{exemplebox}[Version 2 : Modal \zh{会} + Temps \zh{以后} + Comparatif \zh{更}]
\textbf{Chinois :} \zh{我觉得喀麦隆的经济以后会更发展。}
\textbf{Pinyin :} Wǒ juéde Kāmàilóng de jīngjì yǐhòu huì gèng fāzhǎn.
\begin{enumerate}
\item \textbf{Segmentation :} \zh{我} (S1) | \zh{觉得} (V1) | \zh{喀麦隆的经济} (S2) | \zh{以后} (Temps) | \zh{会} (Modal) | \zh{更} (Degré) | \zh{发展} (V2).
\item \textbf{Ordre chinois :} [S2] + [Temps \zh{以后}] + [Modal \zh{会}] + [Degré \zh{更}] + [V].
\item \textbf{Ordre français :} [S2] + [Verbe Futur] + [Degré] + [Temps].
\item \textbf{Blocs :} « Je pense que » + « l'économie du Cameroun » + « à l'avenir » + « sera » + « plus développée / se développera davantage ».
\item \textbf{Résultat :} \textbf{« Je pense que l'économie du Cameroun se développera davantage à l'avenir. »}
\end{enumerate}
\end{exemplebox}

\begin{attention}[Point clé : Place de \zh{以后} (yǐhòu)]
En français : « à l'avenir » souvent en fin. En chinois : \textbf{AVANT le verbe/modal}.
\begin{center}
\begin{tabularx}{\linewidth}{|X|X|X|}
\hline \rowcolor{popblueL} Chinois & Structure & Français \\ \hline
\zh{以后会更发展} & Temps + Modal + Degré + V & Se développera davantage à l'avenir \\ \hline
\zh{以后去北京} & Temps + V & Ira à Pékin à l'avenir \\ \hline
\end{tabularx}
\end{center}
\textbf{Erreur fréquente :} *\zh{会更发展以后} (Calque français). \textbf{Interdit.}
\end{attention}

\begin{propriete}[Comparatif \textbf{更} (gèng) = « plus / davantage »]
\zh{更} s'utilise \textbf{sans \zh{比}} (comparaison implicite avec état actuel/passé).
\zh{Structure :} Sujet + \zh{更} + Adj / Verbe.
\zh{Ex:} \zh{这个更贵。} (Celui-ci est plus cher.)
\zh{Ex:} \zh{我们要更努力。} (Nous devons travailler plus dur.)
\emph{Différence :} \zh{比} compare A et B explicitement. \zh{更} compare à un état de référence implicite.
\end{propriete}

\coursec{Production écrite et évaluation}

\courssub{Entraînement mixte : Thème \& Version}

\begin{exercice}[Entraînement Thème (Français $\rightarrow$ Chinois)]
Traduire les phrases suivantes en chinois (caractères + pinyin). Utilisez le vocabulaire de la leçon.
\begin{enumerate}
\item Le Cameroun a beaucoup de ressources naturelles. (ressource naturelle = \zh{自然资源} zìrán zīyuán)
\item Parce que le port de Douala est grand, le commerce est développé. (port = \zh{港口}, grand = \zh{大}, commerce = \zh{贸易}, développé = \zh{发达})
\item Je pense que l'agriculture est très importante. (agriculture = \zh{农业}, important = \zh{重要})
\item Le pétrole et le bois sont les principales exportations. (pétrole = \zh{石油}, bois = \zh{木材}, principal = \zh{主要}, exportation = \zh{出口} [nom])
\item À l'avenir, le Cameroun sera plus fort. (fort = \zh{强})
\end{enumerate}
\end{exercice}

\begin{corrige}[Exercice Thème]
\begin{enumerate}
\item \zh{喀麦隆有很多自然资源。} (Kāmàilóng yǒu hěn duō zìrán zīyuán.)
\item \zh{因为杜阿拉的港口很大，所以贸易很发达。} (Yīnwèi Dù'ālā de gǎngkǒu hěn dà, suǒyǐ màoyì hěn fādá.) \emph{Note : \zh{发达} s'utilise pour économie/commerce/transport/technologie.}
\item \zh{我觉得农业非常重要。} (Wǒ juéde nóngyè fēicháng zhòngyào.)
\item \zh{石油和木材是主要的出口。} (Shíyóu hé múcái shì zhǔyào de chūkǒu.) \emph{Variante plus précise :} \zh{石油和木材是主要出口产品。} (sont les principaux produits d'exportation).
\item \zh{以后喀麦隆会更强。} (Yǐhòu Kāmàilóng huì gèng qiáng.) \emph{Place de \zh{以后} : en tête ou après sujet (\zh{喀麦隆以后会更强}).}
\end{enumerate}
\end{corrige}

\begin{exercice}[Entraînement Version (Chinois $\rightarrow$ Français)]
Traduire les phrases suivantes en français naturel et correct.
\begin{enumerate}
\item \zh{喀麦隆的农业比工业重要。}
\item \zh{因为基础设施不好，所以运输很慢。} (infrastructure = \zh{基础设施}, transport = \zh{运输}, lent = \zh{慢})
\item \zh{我们觉得以后旅游业会更好。} (tourisme = \zh{旅游业})
\item \zh{喀麦隆出口的产品越来越多。} (produit = \zh{产品})
\item \zh{杜阿拉和雅温得是最重要的城市。} (Yaoundé = \zh{雅温得})
\end{enumerate}
\end{exercice}

\begin{corrige}[Exercice Version]
\begin{enumerate}
\item \textbf{L'agriculture du Cameroun est plus importante que l'industrie.} (Structure \zh{比} : A \zh{比} B Adj).
\item \textbf{Parce que les infrastructures ne sont pas bonnes, le transport est très lent.} (\zh{不好} = pas bon/mauvais ; \zh{慢} = lent).
\item \textbf{Nous pensons que le tourisme sera meilleur à l'avenir. / ... se développera mieux à l'avenir.} (\zh{会更好} $\rightarrow$ futur + comparatif ; \zh{以后} $\rightarrow$ à l'avenir).
\item \textbf{Les produits exportés par le Cameroun sont de plus en plus nombreux.} (\zh{出口的产品} = produits [que le Cameroun] exporte -- relative sans \zh{的} visible mais implicite ; \zh{越来越多} = de plus en plus nombreux).
\item \textbf{Douala et Yaoundé sont les villes les plus importantes.} (\zh{最} + Adj + \zh{的} + Nom = superlatif).
\end{enumerate}
\end{corrige}

\courssub{Production écrite guidée : « Mon avis sur l'économie camerounaise »}

\begin{methode}[Consignes de rédaction (100-150 caractères chinois)]
\begin{enumerate}
\item \textbf{Préparer :} Faire un plan en 4 idées : 1. Situation géo/ressources. 2. Atouts (port, exportations). 3. Problèmes (infrastructures, industrie). 4. Avenir (conditionnel \zh{如果}...\zh{会}...).
\item \textbf{Rédiger :} Utiliser \textbf{au moins 5 connecteurs} de la boîte \emph{À retenir} (\zh{因为}...\zh{所以}, \zh{例如}, \zh{比}, \zh{越来越}, \zh{如果}...\zh{会}, \zh{最}, \zh{更}).
\item \textbf{Vérifier (Check-list) :}
\begin{itemize}
\item Ordre des mots (Temps/Lieu avant V) ? \checkmark
\item \zh{的}/\zh{得}/\zh{地} corrects ? \checkmark
\item Classificateurs si Nombre + Nom ? \checkmark
\item Ponctuation chinoise (， 。) ? \checkmark
\item Pinyin mental correct (tons) ? \checkmark
\end{itemize}
\end{enumerate}
\end{methode}

\begin{exemplebox}[Modèle de production élève (Niveau visé)]
\zh{喀麦隆有很多石油和木材，所以出口赚很多钱。} \\
\zh{因为杜阿拉港口最大，所以贸易很发达。} \\
\zh{但是基础设施不太好，运输比较慢。} \\
\zh{我觉得经济越来越好。} \\
\zh{如果以后工业更发展，喀麦隆会更强。} \\

\textbf{Pinyin :} Kāmàilóng yǒu hěn duō shíyóu hé múcái, suǒyǐ chūkǒu zhuàn hěn duō qián. Yīnwèi Dù'ālā gǎngkǒu zuì dà, suǒyǐ màoyì hěn fādá. Dànshì jīchǔ shèshī bù tài hǎo, yùnshū bǐjiào màn. Wǒ juéde jīngjì yuè lái yuè hǎo. Rúguǒ yǐhòu gōngyè gèng fāzhǎn, Kāmàilóng huì gèng qiáng.

\textbf{Traduction :} Le Cameroun a beaucoup de pétrole et de bois, donc les exportations rapportent beaucoup d'argent. Parce que le port de Douala est le plus grand, le commerce est très développé. Mais les infrastructures ne sont pas très bonnes, le transport est assez lent. Je pense que l'économie va de mieux en mieux. Si l'industrie se développe davantage à l'avenir, le Cameroun sera plus fort.
\end{exemplebox}

\begin{experience}[Débat oral guidé : « Cameroun vs Chine : modèles de développement »]
\textbf{Consigne :} En binôme, préparez un dialogue de 2 minutes. L'élève A est un étudiant camerounais, l'élève B un étudiant chinois. Utilisez \textbf{impérativement} : \zh{因为}...\zh{所以}..., \zh{比}, \zh{越来越}, \zh{以后会}..., \zh{如果}...\zh{就/会}...
\textbf{Rôles \& Pistes :}
\begin{itemize}
\item \textbf{A (Cameroun) :} \zh{喀麦隆有石油、木材，资源比中国丰富。} (Ressources plus abondantes.)
\item \textbf{B (Chine) :} \zh{中国的基础设施比喀麦隆好，高铁、公路很发达。} (Infra meilleures, TGV, routes développées.)
\item \textbf{A :} \zh{我觉得喀麦隆以后会更发展，因为人口年轻，市场大。} (Je pense développement futur : population jeune, grand marché.)
\item \textbf{B :} \zh{如果中喀合作更多，两国经济都会更强。} (Si coopération sino-camerounaise plus forte, les deux économies plus fortes.)
\end{itemize}
\textbf{Objectif :} Réinvestir lexique économie + connecteurs logiques à l'oral. Évaluation : fluidité, justesse tons/structures, richesse vocabulaire.
\end{experience}

\begin{savaistu}[Le saviez-vous ? : Douala, poumon économique \& Coopération sino-camerounaise]
\zh{杜阿拉} (Dù'ālā) est la capitale économique du Cameroun. Son port autonome gère plus de \textbf{90\%} des échanges commerciaux extérieurs du pays (cacao, café, coton, bois, pétrole, aluminium). La ville concentre aussi la majorité des industries (brasseries, raffinerie, textile, agro-alimentaire).

Comparée à \zh{上海} (Shànghǎi) en Chine, Douala joue un rôle similaire de « porte d'entrée » maritime (\zh{门户} ménhù), bien que Shanghai soit une mégapole mondiale de 24 millions d'habitants contre environ 3,5 millions pour Douala.

Le projet de port en eau profonde de \zh{克里比} (Kèlǐbǐ, Kribi), financé en partie par des investisseurs chinois (\zh{中国港湾} Zhōngguó Gǎngwān), illustre la coopération sino-camerounaise dans les infrastructures (\zh{基础设施合作} jīchǔ shèshī hézuò). La Chine est le premier partenaire commercial du Cameroun et un financeur majeur de ses routes, barrages (\zh{水电站} shuǐdiànzhàn) et stades (\zh{体育场} tǐyùchǎng).
\end{savaistu}

\coursec{Production écrite et évaluation}

Après avoir traduit des phrases en thème et en version, vous êtes maintenant prêts pour l'étape la plus importante de la leçon : \cle{produire votre propre texte} en chinois sur l'économie du Cameroun. C'est exactement ce type de tâche qui vous sera demandé à l'examen. Nous allons procéder méthodiquement : d'abord observer un texte modèle complet, puis comprendre la tâche, puis rédiger, relire et évaluer.

\courssub{Un texte modèle complet à observer}

Avant d'écrire, lisez attentivement ce texte d'élève de Première. Il respecte exactement le plan et les contraintes de la tâche finale. Repérez les connecteurs et les mots du glossaire.

\begin{textebox}[喀麦隆的经济 — Texte modèle d'élève]
喀麦隆是非洲中部的一个国家。\\
Kāmàilóng shì Fēizhōu zhōngbù de yí ge guójiā.\\
« Le Cameroun est un pays d'Afrique centrale. »\\[4pt]
喀麦隆的经济比一些非洲国家发展得快，但是比中国发展得慢。\\
Kāmàilóng de jīngjì bǐ yìxiē Fēizhōu guójiā fāzhǎn de kuài, dànshì bǐ Zhōngguó fāzhǎn de màn.\\
« L'économie du Cameroun se développe plus vite que celle de certains pays africains, mais plus lentement que celle de la Chine. »\\[4pt]
因为喀麦隆有很多自然资源，比如石油、木材和可可，所以它的经济一直在发展。\\
Yīnwèi Kāmàilóng yǒu hěn duō zìrán zīyuán, bǐrú shíyóu, mùcái hé kěkě, suǒyǐ tā de jīngjì yìzhí zài fāzhǎn.\\
« Parce que le Cameroun a beaucoup de ressources naturelles, comme le pétrole, le bois et le cacao, son économie se développe continuellement. »\\[4pt]
而且，喀麦隆的农业也很重要：很多农民种咖啡和香蕉。\\
Érqiě, Kāmàilóng de nóngyè yě hěn zhòngyào : hěn duō nóngmín zhòng kāfēi hé xiāngjiāo.\\
« De plus, l'agriculture du Cameroun est aussi très importante : beaucoup de paysans cultivent du café et des bananes. »\\[4pt]
但是，喀麦隆的工业还不太发达，很多产品要从外国进口。\\
Dànshì, Kāmàilóng de gōngyè hái bù tài fādá, hěn duō chǎnpǐn yào cóng wàiguó jìnkǒu.\\
« Cependant, l'industrie du Cameroun n'est pas encore très développée, et beaucoup de produits doivent être importés de l'étranger. »\\[4pt]
我觉得喀麦隆的经济以后会更好，因为年轻人越来越努力。\\
Wǒ juéde Kāmàilóng de jīngjì yǐhòu huì gèng hǎo, yīnwèi niánqīngrén yuè lái yuè nǔlì.\\
« Je pense que l'économie du Cameroun sera meilleure à l'avenir, parce que les jeunes travaillent de plus en plus dur. »
\end{textebox}

\begin{exemplebox}[Observation guidée]
Observez le texte modèle et répondez :
\begin{enumerate}
\item Combien de phrases contient-il ? (6 phrases, chacune terminée par \zh{。})
\item Quels connecteurs voyez-vous ? (\zh{但是} « mais », \zh{因为...所以...} « parce que... donc... », \zh{而且} « de plus »)
\item Quels mots du glossaire sont utilisés ? (\zh{经济}, \zh{发展}, \zh{石油}, \zh{农业}, \zh{工业}, \zh{进口}, \zh{自然资源}...)
\item Où se trouve l'opinion personnelle ? (à la fin, avec \zh{我觉得...})
\end{enumerate}
Remarquez la progression logique : présentation générale → comparaison → cause → exemple → limite → opinion. C'est exactement le plan imposé de votre tâche finale.
\end{exemplebox}

\courssub{La tâche finale : consignes et plan imposé}

\begin{definition}[Tâche finale de la leçon]
Rédigez un texte de \cle{6 à 8 phrases} en chinois sur le sujet : « L'économie du Cameroun » (\zh{喀麦隆的经济}).
\end{definition}

\begin{propriete}[Plan imposé en cinq étapes]
Votre texte doit suivre obligatoirement cette progression :
\begin{enumerate}
\item \textbf{Introduction} : présenter le Cameroun et son économie en une phrase générale.\\ Modèle : \zh{喀麦隆是非洲中部的一个国家，它的经济在发展。} (\emph{Kāmàilóng shì Fēizhōu zhōngbù de yí ge guójiā, tā de jīngjì zài fāzhǎn.}) « Le Cameroun est un pays d'Afrique centrale, son économie est en développement. »
\item \textbf{Comparaison avec un pays} : utiliser la structure comparative avec \zh{比}.\\ Modèle : \zh{喀麦隆的经济比乍得发展得快。} (\emph{Kāmàilóng de jīngjì bǐ Zhàdé fāzhǎn de kuài.}) « L'économie du Cameroun se développe plus vite que celle du Tchad. »
\item \textbf{Causes} : expliquer avec \zh{因为...所以...}.\\ Modèle : \zh{因为喀麦隆有石油和可可，所以经济一直在发展。} « Parce que le Cameroun a du pétrole et du cacao, son économie se développe continuellement. »
\item \textbf{Exemples} : citer des produits ou secteurs avec \zh{比如} « par exemple ».\\ Modèle : \zh{喀麦隆出口很多产品，比如石油、木材和咖啡。} « Le Cameroun exporte beaucoup de produits, par exemple le pétrole, le bois et le café. »
\item \textbf{Opinion personnelle} : conclure avec \zh{我觉得...}.\\ Modèle : \zh{我觉得喀麦隆的经济以后会更好。} « Je pense que l'économie du Cameroun sera meilleure à l'avenir. »
\end{enumerate}
\end{propriete}

\begin{propriete}[Contraintes obligatoires]
\begin{itemize}
\item \textbf{Au moins 3 connecteurs} parmi : \zh{因为...所以...} (cause-conséquence), \zh{但是} (opposition), \zh{而且} (addition).
\item \textbf{Au moins 6 mots du glossaire} du thème : \zh{经济}, \zh{发展}, \zh{石油}, \zh{农业}, \zh{工业}, \zh{出口}, \zh{进口}, \zh{产品}, \zh{资源}, \zh{市场}, \zh{农民}, \zh{国家}...
\item \textbf{6 à 8 phrases}, chacune terminée par le point plein \zh{。}.
\item Caractères soignés, traits corrects.
\end{itemize}
\end{propriete}

\begin{center}
\begin{tabularx}{\linewidth}{|X|X|X|X|}\hline
\rowcolor{popblueL} Étape du plan & Structure-clé & Exemple en caractères + pinyin & Traduction \\ \hline
Introduction & \zh{是} + nom & \zh{喀麦隆是一个非洲国家。} \emph{Kāmàilóng shì yí ge Fēizhōu guójiā.} & Le Cameroun est un pays africain. \\ \hline
Comparaison & A + \zh{比} + B + adj. & \zh{经济比乍得发展得快。} \emph{jīngjì bǐ Zhàdé fāzhǎn de kuài.} & L'économie se développe plus vite que celle du Tchad. \\ \hline
Cause & \zh{因为}... \zh{所以}... & \zh{因为有石油，所以经济在发展。} \emph{Yīnwèi yǒu shíyóu, suǒyǐ jīngjì zài fāzhǎn.} & Parce qu'il y a du pétrole, l'économie se développe. \\ \hline
Exemple & \zh{比如} + liste & \zh{比如可可和咖啡。} \emph{bǐrú kěkě hé kāfēi.} & par exemple le cacao et le café. \\ \hline
Opinion & \zh{我觉得}... & \zh{我觉得以后会更好。} \emph{Wǒ juéde yǐhòu huì gèng hǎo.} & Je pense que ce sera meilleur à l'avenir. \\ \hline
\end{tabularx}
\end{center}

\courssub{Méthode de rédaction en quatre temps}

\begin{methode}[Réussir sa production écrite : brouillon → copie au propre]
\begin{enumerate}
\item \textbf{Brouillon} : écrivez vos 6 à 8 phrases rapidement, en suivant le plan imposé. Ne cherchez pas la perfection : notez d'abord les idées. Vérifiez que chaque étape du plan (introduction, comparaison, causes, exemples, opinion) est présente.
\item \textbf{Vérification des caractères} : relisez caractère par caractère. Contrôlez surtout \zh{经济} (clé 纟 pour \zh{经}, clé 氵pour \zh{济}), \zh{发展} (ne pas confondre \zh{发} et \zh{友}), \zh{石油} (\zh{油} avec la clé de l'eau 氵). Un caractère mal écrit coûte des points.
\item \textbf{Vérification des connecteurs} : comptez vos connecteurs. Avez-vous bien \zh{因为...所以...}, \zh{但是}, \zh{而且} ? Vérifiez que \zh{因为} et \zh{所以} sont bien dans la même phrase, séparés par une virgule chinoise \zh{，}.
\item \textbf{Copie au propre} : recopiez lentement, traits nets, chaque phrase terminée par \zh{。}. Une copie propre et lisible rapporte des points à l'examen.
\end{enumerate}
\end{methode}

\begin{attention}[La ponctuation chinoise : le point plein]
Chaque phrase chinoise se termine par \cle{。} (point plein, petit cercle), \textbf{jamais} par le point occidental « . ». De même, la virgule chinoise est \zh{，} et non « , », et les deux-points sont \zh{：}.
\begin{itemize}
\item Correct : \zh{我觉得喀麦隆的经济以后会更好。}
\item Incorrect : \zh{我觉得喀麦隆的经济以后会更好.}
\end{itemize}
Erreur fréquente des francophones : oublier le point final ou mettre une virgule occidentale entre \zh{因为} et \zh{所以}. Écrivez : \zh{因为喀麦隆有石油，所以经济在发展。}
\end{attention}

\begin{exemplebox}[La phrase modèle de conclusion à retenir pour l'examen]
\zh{我觉得喀麦隆的经济以后会更好。}\\
\emph{Wǒ juéde Kāmàilóng de jīngjì yǐhòu huì gèng hǎo.}\\
« Je pense que l'économie du Cameroun sera meilleure à l'avenir. »\\[4pt]
Structure : \zh{我觉得} (je pense que) + sujet + \zh{以后} (à l'avenir) + \zh{会} (futur/probabilité) + \zh{更} + adjectif (encore plus...). Variantes utiles :
\begin{itemize}
\item \zh{我觉得喀麦隆的农业以后会更发达。} « Je pense que l'agriculture du Cameroun sera plus développée à l'avenir. »
\item \zh{我觉得以后会有更多中国公司来喀麦隆。} « Je pense qu'à l'avenir, plus d'entreprises chinoises viendront au Cameroun. »
\end{itemize}
\end{exemplebox}

\courssub{La grille d'évaluation}

Votre texte sera noté sur 20 points selon cinq critères. Lisez bien cette grille \textbf{avant} de rédiger : elle vous dit exactement ce que le correcteur attend.

\begin{popfigure}
\begin{tikzpicture}[x=1cm,y=1cm]
\fill[popblueL] (0,0) rectangle (10.5,0.9);
\node[anchor=west, font=\bfseries] at (0.15,0.45) {Critère};
\node[anchor=west, font=\bfseries] at (6.2,0.45) {Points};
\draw[popline] (0,0) rectangle (10.5,-5.6);
\foreach \y in {-0.9,-1.85,-2.8,-3.75,-4.7}{\draw[popline] (0,\y) -- (10.5,\y);}
\draw[popline] (6.1,0) -- (6.1,-5.6);
\node[anchor=west, align=left] at (0.15,-0.45) {1. Vocabulaire : au moins 6 mots du glossaire, bien employés};
\node at (6.85,-0.45) {\textbf{4 pts}};
\node[anchor=west, align=left] at (0.15,-1.4) {2. Grammaire et connecteurs : \zh{比}, \zh{因为...所以}, \zh{但是}, \zh{而且}};
\node at (6.85,-1.4) {\textbf{4 pts}};
\node[anchor=west, align=left] at (0.15,-2.35) {3. Écriture des caractères : traits corrects, lisible, \zh{。} final};
\node at (6.85,-2.35) {\textbf{4 pts}};
\node[anchor=west, align=left] at (0.15,-3.3) {4. Structure du texte : plan en 5 étapes respecté, 6 à 8 phrases};
\node at (6.85,-3.3) {\textbf{4 pts}};
\node[anchor=west, align=left] at (0.15,-4.25) {5. Traduction : version française correcte de votre texte};
\node at (6.85,-4.25) {\textbf{4 pts}};
\node[anchor=west, font=\bfseries] at (0.15,-5.15) {TOTAL};
\node[font=\bfseries] at (6.85,-5.15) {20 pts};
\end{tikzpicture}
\legende{Grille d'évaluation de la production écrite (5 critères × 4 points = 20 points)}
\end{popfigure}

\begin{aretenir}[Ce qui fait gagner ou perdre des points]
\begin{itemize}
\item \textbf{Vocabulaire (4 pts)} : un mot du glossaire mal utilisé = point perdu ; un mot français non traduit = point perdu.
\item \textbf{Grammaire (4 pts)} : attention à l'ordre des mots dans la comparaison : A + \zh{比} + B + adjectif. Ne dites jamais \zh{经济快比乍得} !
\item \textbf{Caractères (4 pts)} : un caractère inventé ou avec un trait manquant est compté faux. En cas de doute, écrivez un mot plus simple que vous maîtrisez.
\item \textbf{Structure (4 pts)} : une étape du plan manquante = points perdus, même si les phrases sont correctes.
\item \textbf{Traduction (4 pts)} : vous devez pouvoir traduire votre propre texte en français ; relisez donc chaque phrase en vous demandant : « Qu'est-ce que cela veut dire ? »
\end{itemize}
\end{aretenir}

\courssub{Entraînement guidé avant la rédaction}

\begin{exoresolu}[Assembler une phrase complète avec connecteurs]
Énoncé : combinez les éléments suivants en une seule phrase correcte avec \zh{因为...所以...} : \zh{喀麦隆 / 有 / 很多 / 自然资源 / 经济 / 发展}.\\
\tcblower
Solution détaillée :
\begin{enumerate}
\item Identifier la cause : \zh{喀麦隆有很多自然资源} « le Cameroun a beaucoup de ressources naturelles ».
\item Identifier la conséquence : \zh{经济在发展} « l'économie se développe ».
\item Assembler avec \zh{因为} au début de la cause et \zh{所以} au début de la conséquence, séparées par \zh{，} :
\end{enumerate}
\zh{因为喀麦隆有很多自然资源，所以经济在发展。}\\
\emph{Yīnwèi Kāmàilóng yǒu hěn duō zìrán zīyuán, suǒyǐ jīngjì zài fāzhǎn.}\\
« Parce que le Cameroun a beaucoup de ressources naturelles, son économie se développe. »\\
Piège évité : ne pas mettre \zh{所以} avant le sujet de la conséquence si le sujet change ; ici le sujet \zh{经济} vient après \zh{所以}, c'est correct.
\end{exoresolu}

\begin{exercice}[À vous : préparez votre brouillon]
\begin{enumerate}
\item Écrivez une phrase de comparaison entre le Cameroun et un pays de votre choix avec \zh{比}.
\item Écrivez une phrase de cause avec \zh{因为...所以...} sur les ressources du Cameroun.
\item Écrivez une phrase d'opposition avec \zh{但是} (par exemple : agriculture forte / industrie faible).
\item Écrivez votre conclusion personnelle avec \zh{我觉得...以后会更...}.
\end{enumerate}
\end{exercice}

\begin{corrige}[Exercice : préparez votre brouillon]
Propositions de réponses (d'autres réponses correctes sont possibles) :
\begin{enumerate}
\item \zh{喀麦隆的经济比加蓬发展得慢。} \emph{Kāmàilóng de jīngjì bǐ Jiāpéng fāzhǎn de màn.} « L'économie du Cameroun se développe plus lentement que celle du Gabon. »
\item \zh{因为喀麦隆有石油和可可，所以经济一直在发展。} « Parce que le Cameroun a du pétrole et du cacao, son économie se développe continuellement. »
\item \zh{喀麦隆的农业很重要，但是工业还不太发达。} « L'agriculture du Cameroun est très importante, mais l'industrie n'est pas encore très développée. »
\item \zh{我觉得喀麦隆的经济以后会更好。} « Je pense que l'économie du Cameroun sera meilleure à l'avenir. »
\end{enumerate}
\end{corrige}

\courssub{Relecture en binôme et correction collective}

\begin{experience}[Relecture croisée avec la grille]
\begin{enumerate}
\item Échangez votre copie au propre avec votre voisin.
\item Notez le texte de votre camarade avec la grille d'évaluation : attribuez de 0 à 4 points pour chacun des 5 critères, en justifiant chaque point retiré (caractère faux, connecteur absent, étape du plan manquante...).
\item Entourez en rouge : les caractères douteux, les phrases sans \zh{。}, les connecteurs mal placés.
\item Rendez la copie avec un conseil d'amélioration écrit en français.
\item Correction collective au tableau : le professeur réécrit deux ou trois phrases fautives fréquentes et la classe les corrige ensemble.
\end{enumerate}
\end{experience}

\begin{attention}[Erreurs fréquentes relevées en correction]
\begin{itemize}
\item Oublier \zh{的} dans \zh{喀麦隆的经济} : sans \zh{的}, le possessif est faux.
\item Mettre l'adjectif avant \zh{比} : on dit \zh{喀麦隆比乍得大}, jamais \zh{喀麦隆大比乍得}.
\item Traduire mot à mot « beaucoup de » par \zh{很多 de} : en chinois, pas de préposition : \zh{很多产品}.
\item Utiliser \zh{是} devant un adjectif : on dit \zh{经济很好}, pas \zh{经济是好}.
\end{itemize}
\end{attention}

\begin{savaistu}[À l'examen, l'écriture soignée rapporte des points]
À l'examen de chinois LV2, les correcteurs attribuent des points à la \cle{qualité graphique} : une écriture lisible et des traits corrects font partie de la note. Entraînez-vous spécialement sur les caractères du thème économique : \zh{经济} (\zh{经} : clé de la soie 纟 à gauche, écrire les trois traits de la clé de haut en bas ; \zh{济} : clé de l'eau 氵, trois gouttes à gauche), \zh{发展} (\zh{发} : cinq traits, attention au dernier point en haut à droite ; \zh{展} : dix traits, l'extérieur avant l'intérieur), \zh{石油} (\zh{石} : cinq traits, horizontal avant vertical ; \zh{油} : huit traits, la clé de l'eau d'abord). Un caractère bien tracé montre que vous connaissez l'ordre des traits : c'est une compétence évaluée !
\end{savaistu}

\begin{aretenir}[L'essentiel de la section]
\begin{itemize}
\item Tâche finale : 6 à 8 phrases sur \zh{喀麦隆的经济}, plan en 5 étapes (introduction → comparaison → causes → exemples → opinion).
\item Contraintes : 3 connecteurs minimum (\zh{因为...所以}, \zh{但是}, \zh{而且}) et 6 mots du glossaire minimum.
\item Méthode : brouillon → caractères → connecteurs → copie au propre.
\item Chaque phrase se termine par \zh{。}.
\item Conclusion type : \zh{我觉得喀麦隆的经济以后会更好。}
\item Évaluation sur 20 points : vocabulaire, grammaire, caractères, structure, traduction (4 pts chacun).
\end{itemize}
\end{aretenir}

\newpage

\coursec{Activité d'intégration}

\coursec{Leçon 34 — Expression écrite : économie du Cameroun (为什么喀麦隆经济比泰国发展得快？)}

\courssub{Objectifs de la leçon}

À la fin de cette leçon, tu seras capable de :
\begin{itemize}
\item \textbf{Comprendre} un texte informatif simple sur l'économie du Cameroun (vocabulaire HSK 2-3 : \zh{经济}, \zh{发展}, \zh{出口}, \zh{进口}, \zh{资源}, \zh{港口}, \zh{农业}, \zh{工业}) ;
\item \textbf{Écrire correctement} les caractères clés du thème : ordre des traits, radicaux, décomposition, discrimination des caractères proches (\zh{出}/\zh{山}, \zh{石}/\zh{右}, \zh{港}/\zh{江}) ;
\item \textbf{Mobiliser} les structures syntaxiques pour \cle{comparer} (\zh{A 比 B + Adj.}), \cle{expliquer une cause} (\zh{因为……所以……}), \cle{exprimer une opinion} (\zh{我觉得……}) et \cle{articuler} un propos (\zh{还有}, \zh{但是}, \zh{而且}) ;
\item \textbf{Rédiger} un article court (6-8 phrases) structuré, titré \zh{喀麦隆的经济}, destiné au journal scolaire \zh{《友谊报》} ;
\item \textbf{Traduire} (thème version) des phrases et un court texte sur ce thème, en soignant l'accord des classificateurs et la place des compléments.
\end{itemize}

\courssub{Situation de communication}

Tu es élève en Première A au lycée de Yaoundé. Ton établissement publie chaque trimestre un journal scolaire bilingue sino-camerounais, \zh{《友谊报》} (\emph{Yǒuyì Bào}, « Le Journal de l'Amitié »), lu par les élèves camerounais et par les élèves chinois de l'Institut Confucius. Pour le prochain numéro, consacré au « Monde du travail », la rédaction cherche de \cle{jeunes reporters économiques}. Ton professeur de chinois te propose, avec ton binôme, d'écrire un article intitulé \zh{喀麦隆的经济} (« L'économie du Cameroun »). Les meilleurs articles seront affichés au mur de la classe et publiés dans le journal.

\courssub{Étape 1 : Modèle de texte commenté}

\begin{textebox}[Document d'appui : 喀麦隆经济简介]
喀麦隆是一个非洲国家，经济正在发展。\\
\emph{Kāmàilóng shì yí ge Fēizhōu guójiā, jīngjì zhèngzài fāzhǎn.}\\
« Le Cameroun est un pays africain dont l'économie est en développement. »\\
喀麦隆出口可可、咖啡和石油。\\
\emph{Kāmàilóng chūkǒu kěkě, kāfēi hé shíyóu.}\\
« Le Cameroun exporte du cacao, du café et du pétrole. »\\
杜阿拉有一个很大的港口。\\
\emph{Dù'ālā yǒu yí ge hěn dà de gǎngkǒu.}\\
« Douala possède un très grand port. »
\end{textebox}

\begin{aretenir}[Analyse du texte support]
\begin{itemize}
\item \textbf{Structure globale} : Présentation du pays (Sujet + \zh{是} + Nom) $\rightarrow$ État de l'économie (\zh{正在} + V) $\rightarrow$ Activités principales (Sujet + V + Objets) $\rightarrow$ Infrastructure clé (Lieu + \zh{有} + Objet).
\item \textbf{Grammaire repérée} : \zh{正在} (action en cours) + Verbe ; \zh{和} (énumération exhaustive) ; \zh{一个} (classificateur générique) ; \zh{很大的} (adjectif attributif + \zh{的}).
\item \textbf{Vocabulaire clé} : \zh{非洲} (Afrique), \zh{国家} (pays), \zh{出口} (exporter), \zh{可可} (cacao), \zh{咖啡} (café), \zh{石油} (pétrole), \zh{杜阿拉} (Douala), \zh{港口} (port).
\end{itemize}
\end{aretenir}

\courssub{Étape 2 : Vocabulaire et glossaire du thème}

\begin{center}
\begin{tabularx}{\linewidth}{|X|X|X|X|}
\hline
\rowcolor{popblueL} Caractères & Pinyin & Nature & Traduction \\
\hline
经济 & jīngjì & nom & économie \\
\hline
发展 & fāzhǎn & verbe / nom & (se) développer, développement \\
\hline
出口 & chūkǒu & verbe & exporter \\
\hline
进口 & jìnkǒu & verbe & importer \\
\hline
可可 & kěkě & nom & cacao \\
\hline
咖啡 & kāfēi & nom & café \\
\hline
石油 & shíyóu & nom & pétrole \\
\hline
港口 & gǎngkǒu & nom & port (maritime) \\
\hline
资源 & zīyuán & nom & ressource(s) \\
\hline
农业 & nóngyè & nom & agriculture \\
\hline
工业 & gōngyè & nom & industrie \\
\hline
农民 & nóngmín & nom & paysan, agriculteur \\
\hline
工人 & gōngrén & nom & ouvrier, travailleur \\
\hline
泰国 & Tàiguó & nom propre & Thaïlande \\
\hline
中部 & zhōngbù & nom de lieu & centre, partie centrale \\
\hline
越来越 & yuè lái yuè & adv. & de plus en plus \\
\hline
会 & huì & verbe aux. & akan (futur probable), savoir \\
\hline
\end{tabularx}
\end{center}

\courssub{Étape 3 : Écriture correcte des caractères clés}

\begin{methode}[Rappel méthodologique : écrire un caractère]
\begin{enumerate}
\item Identifier le \cle{radical} (clé) et le nombre de traits restants.
\item Respecter l'ordre des traits : \textbf{haut $\rightarrow$ bas, gauche $\rightarrow$ droite, horizontal $\rightarrow$ vertical, centre $\rightarrow$ côtés, fermer en dernier}.
\item Vérifier la structure (haut-bas, gauche-droite, enveloppement) et l'équilibre dans le carré imaginaire.
\end{enumerate}
\end{methode}

\begin{propriete}[Fiches caractères : analyse graphique et pièges]
\begin{center}
\begin{tabularx}{\linewidth}{|X|X|X|X|}
\hline
\rowcolor{popblueL} Caractère & Radical (Clé) & Nb traits & Pièges / Caractères proches \\
\hline
\zh{经} & \zh{纟} (soie, 3 tr.) & 8 au total & Ne pas confondre \zh{巳} (sì, 6e branche terrestre) et \zh{已} (yǐ, déjà). Ici : \zh{巳} (boucle fermée). \\
\hline
\zh{济} & \zh{氵} (eau, 3 tr.) & 13 au total & Droite : \zh{齐} (qí, aligné). Ne pas écrire \zh{戊} (wù) en bas à droite. \\
\hline
\zh{发} & \zh{又} (main droite, 2 tr.) & 5 au total & Forme simplifiée. Traditionnel : \zh{髮} (cheveu + ami). Haut : \zh{癶} (bō) + \zh{一} + \zh{又}. \\
\hline
\zh{展} & \zh{尸} (corps, 3 tr.) & 10 au total & Bas : \zh{单} (dān, simple) et non \zh{千} (qiān). Attention au trait horizontal long final. \\
\hline
\zh{港} & \zh{氵} (eau, 3 tr.) & 12 au total & Droite : \zh{巷} (xiàng, ruelle). Ne pas confondre avec \zh{江} (jiāng, fleuve, droite : \zh{工}). \\
\hline
\zh{口} & \zh{口} (bouche, 3 tr.) & 3 au total & Carré fermé. Différent de \zh{囗} (wéi, enceinte, plus grand) et \zh{日} (rì, soleil, traits verticaux intérieurs ne touchent pas le bas). \\
\hline
\zh{石} & \zh{石} (pierre, 5 tr.) & 5 au total & Horizontal du milieu ne traverse pas. Proche : \zh{右} (yòu, droite : \zh{口} + \zh{工}). \\
\hline
\zh{油} & \zh{氵} (eau, 3 tr.) & 8 au total & Droite : \zh{由} (yóu, raison). Haut de \zh{由} : \zh{田} (champ) pas \zh{由} (différent de \zh{甲} jiǎ). \\
\hline
\zh{出} & \zh{凵} (réceptacle, 2 tr.) & 5 au total & Les deux traits verticaux ne touchent pas le bas. Proche : \zh{山} (shān, montagne, 3 pics verticaux). \\
\hline
\zh{可} & \zh{口} (bouche, 3 tr.) & 5 au total & Droite : \zh{丁} (dīng). Ne pas ajouter de trait horizontal en bas. \\
\hline
\end{tabularx}
\end{center}
\end{propriete}

\begin{experience}[Atelier écriture : « Caractères en miroir »]
En binôme. L'élève A dicte un mot du glossaire (ex. \zh{石油}), l'élève B l'écrit au tableau ou sur ardoise en \textbf{nommant l'ordre des traits à voix haute} (« horizontal, vertical, horizontal... »). Inverser les rôles. Vérifier ensemble les 10 caractères du tableau ci-dessus. Durée : 10 min.
\end{experience}

\courssub{Étape 4 : Structures et connecteurs pour écrire}

\begin{propriete}[Structure 1 : La comparaison avec \cle{比} (bǐ)]
\textbf{Schéma} : \cle{A + 比 + B + Adjectif (sv) (+ 得 + Complément de degré)} \\
\textbf{Règle} : L'adjectif fonctionne comme un verbe d'état. On n'utilise \textbf{pas} \zh{很} (hěn) après \zh{比}. Pour intensifier, on ajoute \zh{得} + complément (vite, beaucoup, un peu).
\begin{center}
\begin{tabularx}{\linewidth}{|X|X|X|}
\hline
\rowcolor{popblueL} Structure & Exemple (Caractères + Pinyin) & Traduction \\
\hline
A 比 B Adj. & \zh{泰国经济比喀麦隆好。} \emph{Tàiguó jīngjì bǐ Kāmàilóng hǎo.} & L'économie thaïlandaise est meilleure que celle du Cameroun. \\
\hline
A 比 B Adj. 得 多/快/早 & \zh{泰国经济比喀麦隆发展得快。} \emph{Tàiguó jīngjì bǐ Kāmàilóng fāzhǎn de kuài.} & L'économie thaïlandaise se développe plus vite que celle du Cameroun. \\
\hline
Négation : A 没有 B 这么/那么 Adj. & \zh{喀麦隆经济没有泰国那么发达。} \emph{Kāmàilóng jīngjì méiyǒu Tàiguó nàme fādá.} & L'économie camerounaise n'est pas aussi développée que la thaïlandaise. \\
\hline
\end{tabularx}
\end{center}
\end{propriete}

\begin{attention}[Erreur fréquente : \zh{*喀麦隆经济比泰国发展}]
\zh{发展} est un verbe (ou nom). Après \zh{比}, il faut un \textbf{adjectif / verbe d'état}. Dire \zh{比泰国发展} est incorrect (comparer à un verbe d'action).
\emph{Correct} : \zh{比泰国发展得快} (plus vite) / \zh{比泰国发达} (plus développé).
\end{attention}


\begin{propriete}[Structure 2 : La cause-conséquence avec \cle{因为……所以……} (yīnwèi... suǒyǐ...)]
\textbf{Schéma} : \cle{因为 + Cause (Sujet + Prédicat)，所以 + Conséquence (Sujet + Prédicat).} \\
La virgule est obligatoire avant \zh{所以}. Le sujet peut être omis s'il est identique dans les deux propositions.
\end{propriete}

\begin{exemplebox}[Exemples]
\zh{因为喀麦隆出口石油，所以经济有收入。} \\
\emph{Yīnwèi Kāmàilóng chūkǒu shíyóu, suǒyǐ jīngjì yǒu shōurù.} \\
« Parce que le Cameroun exporte du pétrole, l'économie a des revenus. »
\end{exemplebox}
\begin{attention}[Piège francophone]
Ne pas mettre un simple nom après \zh{因为}. \textbf{Incorrect} : \zh{*因为石油，所以...} \textbf{Correct} : \zh{因为喀麦隆有石油 / 因为出口石油...} (Sujet + Verbe requis).
\end{attention}


\begin{propriete}[Structure 3 : L'opinion avec \cle{我觉得} (wǒ juéde) et le futur probable \cle{会} (huì)]
\begin{itemize}
\item \zh{我觉得} + Proposition complète. Ex. : \zh{我觉得喀麦隆的未来很好。}
\item \zh{会} + Verbe / Adj. : exprime une prédiction, une certitude faible. Ex. : \zh{经济会越来越好。} (L'économie ira de mieux en mieux).
\item \zh{越来越} + Adj. : « de plus en plus... ». Structure figée.
\end{itemize}
\end{propriete}

\begin{propriete}[Connecteurs d'articulation (niveau HSK 2-3)]
\begin{center}
\begin{tabularx}{\linewidth}{|X|X|X|}
\hline
\rowcolor{popblueL} Connecteur & Pinyin & Emploi / Exemple \\
\hline
\zh{还有} & hái yǒu & Addition (idée supplémentaire). \zh{杜阿拉有港口，还有机场。} \\
\hline
\zh{但是} & dànshì & Opposition / Concession. \zh{经济发展得慢，但是资源多。} \\
\hline
\zh{而且} & érqiě & Addition renforcée (souvent après une qualité positive). \zh{资源多，而且质量好。} \\
\hline
\zh{所以} & suǒyǐ & Conséquence (toujours après \zh{因为}). \\
\hline
\end{tabularx}
\end{center}
\end{propriete}

\courssub{Étape 5 : Thème et version guidés}

\begin{exercice}[Thème guidé : Français $\rightarrow$ Chinois]
Traduis les phrases suivantes en chinois (caractères + pinyin). Utilise le vocabulaire et les structures de la leçon.
\begin{enumerate}
\item Le Cameroun est un pays d'Afrique centrale. (\textit{Indice : 中部})
\item Son économie se développe de plus en plus vite. (\textit{Indice : 越来越快})
\item Le Cameroun exporte du cacao et du pétrole. (\textit{Indice : 和})
\item Parce qu'il y a un grand port à Douala, le commerce est pratique. (\textit{Indice : 因为……所以……，方便 fāngbiàn})
\item Je pense que l'économie sera meilleure qu'avant. (\textit{Indice : 我觉得……会……比以前 yǐqián 好})
\end{enumerate}
\end{exercice}

\begin{corrige}[Exercice Thème guidé]
\begin{enumerate}
\item \zh{喀麦隆是非洲中部的一个国家。} (\emph{Kāmàilóng shì Fēizhōu zhōngbù de yí ge guójiā.})
\item \zh{它的经济发展得越来越快。} (\emph{Tā de jīngjì fāzhǎn de yuè lái yuè kuài.})
\item \zh{喀麦隆出口可可和石油。} (\emph{Kāmàilóng chūkǒu kěkě hé shíyóu.})
\item \zh{因为杜阿拉有一个大港口，所以贸易很方便。} (\emph{Yīnwèi Dù'ālā yǒu yí ge dà gǎngkǒu, suǒyǐ màoyì hěn fāngbiàn.})
\item \zh{我觉得经济会比以前好。} (\emph{Wǒ juéde jīngjì huì bǐ yǐqián hǎo.})
\end{enumerate}
\end{corrige}

\begin{exercice}[Version guidée : Chinois $\rightarrow$ Français]
Traduis en français le court paragraphe suivant.
\begin{textebox}[Texte à traduire]
喀麦隆的农业很重要。因为很多农民种可可和咖啡，所以出口赚很多钱。但是工业比较少。我觉得将来工业会发展得很快，因为国家在建新港口和路。
\end{textebox}
\end{exercice}

\begin{corrige}[Exercice Version guidée]
« L'agriculture du Cameroun est très importante. Parce que beaucoup de paysans cultivent du cacao et du café, l'exportation rapporte beaucoup d'argent. Mais l'industrie est relativement peu développée. Je pense qu'à l'avenir, l'industrie se développera très vite, car le pays construit de nouveaux ports et routes. »
\begin{itemize}
\item \zh{种} (zhòng) : cultiver, planter.
\item \zh{赚} (zhuàn) : gagner (argent).
\item \zh{比较} (bǐjiào) : relativement, assez.
\item \zh{将来} (jiānglái) : à l'avenir, dans le futur.
\item \zh{建} (jiàn) : construire, bâtir.
\item \zh{路} (lù) : route.
\end{itemize}
\end{corrige}

\courssub{Étape 6 : Production écrite et évaluation (Activité intégrée)}

\courssub{Consignes de l'activité finale}

Travail en binôme. Durée totale : 2 h (réparties selon le tableau ci-dessous).

\begin{center}
\begin{tabularx}{\linewidth}{|X|X|X|}
\hline
\rowcolor{popblueL} Phase & Durée & Description \\
\hline
1. Lecture \& Repérage & 10 min & Lire le document d'appui (Étape 1). Souligner mots du glossaire, entourer connecteurs. Choisir 6 mots à utiliser. \\
\hline
2. Planification & 10 min & Rédiger un plan en français (4 parties : Introduction, Comparaison, Causes, Opinion). \\
\hline
3. Rédaction (Brouillon) & 25 min & Écrire l'article en chinois (6-8 phrases). Titre obligatoire : \zh{喀麦隆的经济}. Contraintes : 3 connecteurs, 6 mots glossaire, 1 \zh{比}, 1 \zh{因为……所以……}, 1 \zh{我觉得}. Soigner l'écriture. \\
\hline
4. Échange \& Version & 15 min & Échanger l'article avec un autre binôme. Traduire l'article reçu en français (version) sur feuille à part. \\
\hline
5. Évaluation croisée & 10 min & Noter l'article reçu avec la grille (sur 20). Écrire un conseil d'amélioration. \\
\hline
6. Mise au propre & 20 min & Corriger son article selon les retours. Recopier proprement pour affichage / remise. \\
\hline
\end{tabularx}
\end{center}

\begin{attention}[Rappel des pièges à éviter avant de rendre la copie]
\begin{itemize}
\item \textbf{Comparaison} : \zh{A 比 B Adj.} Pas de \zh{很}. Pas de verbe d'action seul après \zh{比}. (\zh{*比泰国发展} $\rightarrow$ \zh{比泰国发展得快}).
\item \textbf{Cause} : \zh{因为} + Sujet + Verbe. (\zh{*因为石油} $\rightarrow$ \zh{因为出口石油}).
\item \textbf{Aspect} : Fait général / vérité permanente $\rightarrow$ \textbf{pas de \zh{了}} après \zh{发展} ou \zh{出口}.
\item \textbf{Caractères} : \zh{出} (pied sort) $\neq$ \zh{山} (3 pics). \zh{石} (pierre, trait central ne traverse pas) $\neq$ \zh{右} (main droite + bouche). \zh{港} (eau + ruelle) $\neq$ \zh{江} (eau + travail).
\item \textbf{Ponctuation} : Virgule chinoise \zh{，}, point \zh{。}, points de suspension \zh{……}, guillemets \zh{「」}. Pas d'espace entre caractères.
\end{itemize}
\end{attention}

\begin{exoresolu}[Article modèle commenté : 喀麦隆的经济]
Rédige un article de 6 à 8 phrases répondant à toutes les consignes.
\tcblower
\textbf{Article modèle :}\\
喀麦隆的经济\\
\emph{Kāmàilóng de jīngjì}\\
喀麦隆是非洲中部的一个国家，经济正在发展。\\
\emph{Kāmàilóng shì Fēizhōu zhōngbù de yí ge guójiā, jīngjì zhèngzài fāzhǎn.}\\
« Le Cameroun est un pays d'Afrique centrale dont l'économie est en développement. »\\
泰国经济比喀麦隆发展得快，但是喀麦隆有很多资源。\\
\emph{Tàiguó jīngjì bǐ Kāmàilóng fāzhǎn de kuài, dànshì Kāmàilóng yǒu hěn duō zīyuán.}\\
« L'économie thaïlandaise se développe plus vite que celle du Cameroun, mais le Cameroun a beaucoup de ressources. »\\
因为喀麦隆出口可可、咖啡和石油，所以很多农民和工人有工作。\\
\emph{Yīnwèi Kāmàilóng chūkǒu kěkě, kāfēi hé shíyóu, suǒyǐ hěn duō nóngmín hé gōngrén yǒu gōngzuò.}\\
« Parce que le Cameroun exporte du cacao, du café et du pétrole, beaucoup de paysans et d'ouvriers ont du travail. »\\
杜阿拉还有一个很大的港口。\\
\emph{Dù'ālā hái yǒu yí ge hěn dà de gǎngkǒu.}\\
« Douala possède aussi un très grand port. »\\
我觉得喀麦隆的经济会越来越好。\\
\emph{Wǒ juéde Kāmàilóng de jīngjì huì yuè lái yuè hǎo.}\\
« Je pense que l'économie du Cameroun ira de mieux en mieux. »\\[0.3em]
\textbf{Commentaire détaillé des choix linguistiques :}
\begin{itemize}
\item \textbf{Plan respecté} : Phrase 1 (Introduction : localisation + état), Phrase 2 (Comparaison \zh{比} + \zh{发展得快} + opposition \zh{但是}), Phrase 3 (Cause \zh{因为……所以……} + vocabulaire thème), Phrase 4 (Exemple infrastructure + \zh{还有}), Phrase 5 (Opinion \zh{我觉得} + futur \zh{会} + \zh{越来越}).
\item \textbf{Connecteurs (4/3 min)} : \zh{但是}, \zh{因为……所以……}, \zh{还有} (dans \zh{还有}), \zh{会} (marque futur).
\item \textbf{Vocabulaire thème (8/6 min)} : \zh{经济}, \zh{发展}, \zh{出口}, \zh{可可}, \zh{咖啡}, \zh{石油}, \zh{港口}, \zh{资源}, \zh{农民}, \zh{工人}, \zh{国家}.
\item \textbf{Grammaire} : \zh{正在} (progressif) ; \zh{比} + V + \zh{得} + Adv ; \zh{因为/所以} (sujets explicites) ; \zh{会} + \zh{越来越} + Adj (prédiction) ; \zh{的} (attributif) correctement placé.
\item \textbf{Version attendue} : La traduction française sous chaque phrase sert de corrigé type pour la phase d'évaluation croisée.
\end{itemize}
\end{exoresolu}

\courssub{Grille d'évaluation (Total : 20 points)}

\begin{center}
\begin{tabularx}{\linewidth}{|X|c|X|}
\hline
\rowcolor{popblueL} Critère & Points & Indicateurs de réussite \\
\hline
Contenu \& Plan & 4 & 4 parties identifiables (Intro, Comparaison, Cause, Opinion). Titre présent. 6-8 phrases. \\
\hline
Grammaire \& Connecteurs & 6 & \zh{比} correct (Adj/Adv), \zh{因为……所以……} complet, \zh{觉得} + phrase, 3 connecteurs min. utilisés à bon escient. Pas de \zh{了} intempestif. \\
\hline
Vocabulaire du thème & 4 & 6 mots minimum du glossaire employés correctement (nature, place). Précision lexicale (\zh{出口} vs \zh{进口}). \\
\hline
Écriture des caractères & 4 & Caractères lisibles, traits dans l'ordre, proportions respectées, pas de confusion \zh{出/山}, \zh{石/右}, \zh{港/江}. \\
\hline
Traduction (Version) & 2 & Traduction française fidèle, syntaxe française correcte, termes techniques rendus (\zh{发展得快} $\rightarrow$ « se développe plus vite »). \\
\hline
\textbf{Total} & \textbf{20} & \\
\hline
\end{tabularx}
\end{center}

\courssub{Bilan et prolongements}

Cette leçon t'a permis de mobiliser l'ensemble des compétences du Module 4 : \textbf{comprendre} un texte authentique simplifié, \textbf{maîtriser} le lexique économique de base, \textbf{structurer} ta pensée en chinois grâce aux connecteurs logiques (\zh{比}, \zh{因为/所以}, \zh{觉得}, \zh{但是/还有}), \textbf{écrire} les caractères avec rigueur graphique, et \textbf{traduire} dans les deux sens.

\begin{savaistu}[Regard croisé Cameroun - Chine : Les ports]
Le port de \zh{杜阿拉} (Douala) est le poumon économique du Cameroun (95\% du trafic). En Chine, le port de \zh{上海} (Shànghǎi) est le premier port à conteneurs du monde depuis 2010. Les deux pays coopèrent : des entreprises chinoises (ex. \zh{中国港湾} China Harbour Engineering) participent à l'extension du port de \zh{克里比} (Kribi), nouveau port en eau profonde au sud du Cameroun. Cette infrastructure illustre le lien \zh{一带一路} (Yīdài Yīlù, « Nouvelles Routes de la Soie ») entre développement portuaire et croissance économique.
\end{savaistu}

\begin{experience}[Prolongement oral : « Le Journal télévisé »]
En groupe de 4 (2 binômes). Chaque binôme présente son article à l'oral (1 min 30) comme un journal télévisé : un présentateur, un reporter « sur le terrain » (à Douala, Kribi, ou à Yaoundé au ministère de l'Économie). Les autres élèves notent : vocabulaire entendu, justesse des tons, fluidité. Vote du « Meilleur Reportage Éco ».
\end{experience}

\newpage

\coursec{Méthodes et exercices résolus}

\courssub{Méthodes et exercices résolus}

\begin{methode}[Analyser un texte modèle : structure, connecteurs et vocabulaire économique]
\textbf{Objectif :} Comprendre l'organisation d'un texte expositif comparatif (Cameroun vs Chine/Thaïlande) et repérer les outils linguistiques (connecteurs, structures de comparaison, vocabulaire spécialisé).\par
\textbf{Démarche en 5 étapes :}
\begin{enumerate}
    \item \textbf{Lecture globale :} Lire le texte une première fois sans s'arrêter au vocabulaire inconnu. Identifier le \cle{sujet} (l'économie camerounaise), la \cle{thèse} (la comparaison du développement) et la \cle{structure} générale (Introduction $\rightarrow$ Arguments $\rightarrow$ Conclusion).
    \item \textbf{Découpage en paragraphes :} Repérer les marqueurs de progression : \zh{首先}、\zh{其次}、\zh{最后}、\zh{总之}. Associer à chaque paragraphe une \cle{idée directrice} (ex. : ressources naturelles, industrialisation, infrastructures, éducation).
    \item \textbf{Chasse aux connecteurs logiques :} Surligner les mots-outils qui relient les idées : \zh{因为……所以……} (cause-conséquence), \zh{虽然……但是……} (concession), \zh{比} / \zh{没有} (comparaison), \zh{而且} (addition), \zh{例如} (exemple).
    \item \textbf{Extraction du vocabulaire thématique :} Lister les noms (économie, agriculture, industrie, infrastructures), les verbes (développer, investir, exporter) et les adjectifs (rapide, lent, important, faible) liés au thème. Vérifier le pinyin et le ton de chaque mot.
    \item \textbf{Modélisation :} Recopier le squelette du texte (phrases-types à trous) pour s'en servir comme \cle{patron de rédaction} lors de la production écrite.
\end{enumerate}
\textbf{Réflexe :} Toujours annoter le texte : entourez les connecteurs, soulignez les structures \zh{比}/\zh{没有}, encadrez le vocabulaire nouveau.
\end{methode}

\begin{textebox}[Texte modèle : Pourquoi l'économie du Cameroun se développe-t-elle moins vite que celle de la Thaïlande ?]
喀麦隆和泰国都是发展中国家，都有丰富的自然资源。\\
Kāmàilóng hé Tàiguó dōu shì fāzhǎnzhōng guójiā, dōu yǒu fēngfù de zìrán zīyuán.\\
« Le Cameroun et la Thaïlande sont tous deux des pays en développement, et disposent tous deux de riches ressources naturelles. »\\
首先，泰国的工业基础比喀麦隆好。因为泰国早期重视制造业，所以电子产品和汽车出口很多。\\
Shǒuxiān, Tàiguó de gōngyè jīchǔ bǐ Kāmàilóng hǎo. Yīnwèi Tàiguó zǎoqī zhòngshì zhìzàoyè, suǒyǐ diànzǐ chǎnpǐn hé qìchē chūkǒu hěn duō.\\
« Premièrement, la base industrielle de la Thaïlande est meilleure que celle du Cameroun. Parce que la Thaïlande a misé tôt sur l'industrie manufacturière, elle exporte beaucoup de produits électroniques et de voitures. »\\
其次，泰国的基础设施比喀麦隆完善。虽然喀麦隆也在修路和修桥，但是速度不够快。\\
Qícì, Tàiguó de jīchǔ shèshī bǐ Kāmàilóng wánshàn. Suīrán Kāmàilóng yě zài xiū lù hé xiū qiáo, dànshì sùdù bù gòu kuài.\\
« Deuxièmement, les infrastructures de la Thaïlande sont plus abouties que celles du Cameroun. Bien que le Cameroun construise aussi des routes et des ponts, le rythme n'est pas assez rapide. »\\
最后，泰国的人力资源质量比较高。泰国政府在教育上投资很多，所以劳动力技能强。\\
Zuìhòu, Tàiguó de rénlì zīyuán zhìliàng bǐjiào gāo. Tàiguó zhèngfǔ zài jiàoyù shang tóuzī hěn duō, suǒyǐ láodònglì jìnéng qiáng.\\
« Enfin, la qualité des ressources humaines de la Thaïlande est relativement élevée. Le gouvernement thaïlandais investit beaucoup dans l'éducation, donc la main-d'œuvre est qualifiée. »\\
总之，如果喀麦隆想追上泰国，必须加强工业化、改善基础设施、增加教育投资。\\
Zǒngzhī, rúguǒ Kāmàilóng xiǎng zhuīshàng Tàiguó, bìxū jiāqiáng gōngyèhuà, gǎishàn jīchǔ shèshī, zēngjiā jiàoyù tóuzī.\\
« En résumé, si le Cameroun veut rattraper la Thaïlande, il doit renforcer l'industrialisation, améliorer les infrastructures et augmenter l'investissement dans l'éducation. »
\end{textebox}

\begin{exoresolu}[Analyse du texte modèle]
\textbf{Énoncé :} Lisez le texte modèle ci-dessus. Répondez aux questions en français.
\begin{enumerate}
    \item Quel est le connecteur utilisé pour introduire le premier argument ? Le deuxième ? Le troisième ?
    \item Identifiez deux structures de comparaison dans le texte. Traduisez-les.
    \item Trouvez la phrase exprimant la cause et la conséquence (structure \zh{因为……所以……}).
    \item Quel mot exprime la concession ? Quelle est la structure complète ?
    \item Listez 5 mots de vocabulaire économique du texte avec leur traduction.
\end{enumerate}
\tcblower
\textbf{Solution détaillée :}
\begin{enumerate}
    \item \textbf{Connecteurs d'enchaînement :} \zh{首先} (premièrement), \zh{其次} (deuxièmement), \zh{最后} (enfin). Ils structurent la progression argumentative. \zh{总之} (en résumé) introduit la conclusion.
    \item \textbf{Comparaisons :}
    \begin{itemize}
        \item \zh{泰国的工业基础比喀麦隆好。} (Tàiguó de gōngyè jīchǔ bǐ Kāmàilóng hǎo.) : « La base industrielle de la Thaïlande est meilleure que celle du Cameroun. » (Supériorité avec \zh{比}.)
        \item \zh{速度不够快} (sùdù bù gòu kuài) : « le rythme n'est pas assez rapide ». Structure \zh{不够} + adjectif (insuffisance).
        \item \zh{泰国的人力资源质量比较高} (Tàiguó de rénlì zīyuán zhìliàng bǐjiào gāo) : comparatif de degré (\zh{比较} = relativement, assez).
    \end{itemize}
    \item \textbf{Cause $\rightarrow$ Conséquence :} \zh{因为泰国早期重视制造业，所以电子产品和汽车出口很多。} (Yīnwèi Tàiguó zǎoqī zhòngshì zhìzàoyè, suǒyǐ diànzǐ chǎnpǐn hé qìchē chūkǒu hěn duō.) Structure : \zh{因为} [cause] \zh{，所以} [conséquence].
    \item \textbf{Concession :} \zh{虽然喀麦隆也在修路和修桥，但是速度不够快。} (Suīrán Kāmàilóng yě zài xiū lù hé xiū qiáo, dànshì sùdù bù gòu kuài.) Structure : \zh{虽然} [fait admis] \zh{，但是} [réalité contrastée]. \zh{虽然} se place en tête de la première proposition.
    \item \textbf{Vocabulaire économique :}
    \begin{center}\begin{tabularx}{\linewidth}{|X|X|X|X|}\hline \rowcolor{popblueL} Caractères & Pinyin & Nature & Traduction \\ \hline
    发展中国家 & fāzhǎnzhōng guójiā & nom & pays en développement \\ \hline
    自然资源 & zìrán zīyuán & nom & ressources naturelles \\ \hline
    工业基础 & gōngyè jīchǔ & nom & base industrielle \\ \hline
    基础设施 & jīchǔ shèshī & nom & infrastructures \\ \hline
    人力资源 & rénlì zīyuán & nom & ressources humaines \\ \hline
    制造业 & zhìzàoyè & nom & industrie manufacturière \\ \hline
    出口 & chūkǒu & verbe / nom & exporter / exportation \\ \hline
    投资 & tóuzī & verbe / nom & investir / investissement \\ \hline
    追上 & zhuīshàng & verbe & rattraper \\ \hline
    工业化 & gōngyèhuà & nom & industrialisation \\ \hline
    \end{tabularx}\end{center}
\end{enumerate}
\textbf{Conclusion :} Ce texte suit un plan « Constat $\rightarrow$ 3 causes structurelles $\rightarrow$ Recommandations ». Les connecteurs \zh{首先}/\zh{其次}/\zh{最后} et la structure \zh{因为……所以……} sont les piliers de la cohérence.
\end{exoresolu}

\begin{methode}[Écrire correctement les caractères clés du thème économique]
\textbf{Objectif :} Maîtriser l'ordre des traits, identifier le radical (la clé) et éviter les confusions graphiques pour les caractères à haute fréquence du Module 4.\par
\textbf{Démarche en 4 étapes :}
\begin{enumerate}
    \item \textbf{Analyse graphique (clé + partie phonétique) :} Pour chaque caractère, identifier le \cle{radical} (qui porte le sens) et l'éventuelle \cle{phonétique} (qui suggère la prononciation). Ex. : \zh{资} (zī) : radical \zh{贝} (coquillage $\rightarrow$ argent, richesse) + phonétique \zh{次} (cì). \zh{源} (yuán) : radical \zh{氵} (eau) + phonétique \zh{原} (yuán).
    \item \textbf{Ordre des traits (règles d'or) :} Appliquer : haut $\rightarrow$ bas ; gauche $\rightarrow$ droite ; horizontal avant vertical ; extérieur $\rightarrow$ intérieur $\rightarrow$ fermeture ; trait central avant les côtés. \emph{Exemple \zh{经} (jīng), 8 traits :} 1-3. \zh{纟} (radical « soie », à gauche : pli, pli, relevé) ; 4-5. \zh{又} (à droite en haut) ; 6-8. \zh{工} (à droite en bas : horizontal, vertical, horizontal).
    \item \textbf{Caractères proches (pièges visuels) :} Créer des paires de discrimination.
    \begin{itemize}
        \item \zh{资} (zī, ressource, radical \zh{贝}) vs \zh{咨} (zī, consulter, radical \zh{口}).
        \item \zh{产} (chǎn, produire) vs \zh{严} (yán, strict) : mêmes premiers traits, mais \zh{产} finit par un trait oblique vers la gauche.
        \item \zh{构} (gòu, structure, radical \zh{木}) vs \zh{沟} (gōu, fossé, radical \zh{氵}) : même phonétique \zh{勾}, radicaux différents.
        \item \zh{经} (jīng, économie) vs \zh{轻} (qīng, léger, radical \zh{车}) : même partie droite.
    \end{itemize}
    \item \textbf{Entraînement guidé :} Écrire 5 fois chaque caractère en dictant l'ordre des traits à voix haute. Puis écrire 3 mots composés contenant ce caractère (ex. : \zh{资源}、\zh{资本}、\zh{投资}).
\end{enumerate}
\textbf{Outil :} la grille de 9 cases (九宫格) est obligatoire pour respecter les proportions.
\end{methode}

\begin{exoresolu}[Écriture et discrimination de caractères : \zh{资}、\zh{源}、\zh{构}、\zh{产}、\zh{基}]
\textbf{Énoncé :}
\begin{enumerate}
    \item Pour chacun des 5 caractères, donnez : le radical (nom + sens), le nombre total de traits, et un mot composé du vocabulaire de la leçon.
    \item Complétez les phrases avec le bon mot (attention aux paires proches) :
    \begin{itemize}
        \item[a.] 喀麦隆有丰富的自然\_\_\_\_。 (ressources)
        \item[b.] 我们要调整经济\_\_\_\_。 (structure)
        \item[c.] 这家工厂\_\_\_\_很多汽车。 (produit)
        \item[d.] 公司需要更多的\_\_\_\_。 (capitaux)
    \end{itemize}
    \item Identifiez l'erreur d'ordre des traits dans l'écriture proposée pour \zh{经} (jīng) : l'élève écrit d'abord la partie droite (\zh{又} puis \zh{工}), puis le radical \zh{纟}.
\end{enumerate}
\tcblower
\textbf{Solution détaillée :}
\begin{enumerate}
    \item \textbf{Fiches caractères :}
    \begin{center}\begin{tabularx}{\linewidth}{|X|X|X|X|}\hline \rowcolor{popblueL} Caractère & Radical (clé) & Nb de traits & Mot composé \\ \hline
    \zh{资} & \zh{贝} (bèi, coquillage $\rightarrow$ argent) & 10 & \zh{资源} (zīyuán), \zh{资本} (zīběn) \\ \hline
    \zh{源} & \zh{氵} (shuǐ, eau) & 13 & \zh{资源} (zīyuán), \zh{能源} (néngyuán) \\ \hline
    \zh{构} & \zh{木} (mù, bois) & 8 & \zh{结构} (jiégòu), \zh{构成} (gòuchéng) \\ \hline
    \zh{产} & \zh{亠} (tóu, couvercle) & 6 & \zh{产品} (chǎnpǐn), \zh{生产} (shēngchǎn) \\ \hline
    \zh{基} & \zh{土} (tǔ, terre, base) & 11 & \zh{基础} (jīchǔ), \zh{基本} (jīběn) \\ \hline
    \end{tabularx}\end{center}
    \item \textbf{Complétion (discrimination) :}
    \begin{itemize}
        \item[a.] \zh{资源} (zīyuán). \textbf{Piège :} ne pas écrire \zh{咨} (consulter, radical \zh{口}).
        \item[b.] \zh{结构} (jiégòu). \textbf{Piège :} ne pas écrire \zh{沟} (gōu, fossé, radical \zh{氵}).
        \item[c.] \zh{生产} (shēngchǎn). \zh{产} (6 traits) : point, horizontal, point, oblique, horizontal, oblique descendante à gauche.
        \item[d.] \zh{资本} (zīběn). \zh{资} (richesse) + \zh{本} (base, capital).
    \end{itemize}
    \item \textbf{Erreur d'ordre (\zh{经}) :} L'élève viole la règle \textbf{« gauche avant droite »} : le radical \zh{纟} (à gauche) s'écrit toujours en premier. Ordre correct : 1-3. \zh{纟} ; 4-5. \zh{又} ; 6-8. \zh{工}.
\end{enumerate}
\textbf{Conclusion :} La connaissance du radical guide le sens (argent $\rightarrow$ \zh{贝} ; bois $\rightarrow$ \zh{木} ; terre $\rightarrow$ \zh{土} ; eau $\rightarrow$ \zh{氵}) et l'ordre des traits assure la vitesse et la lisibilité (l'écriture est notée à l'examen).
\end{exoresolu}

\begin{methode}[Utiliser les structures de comparaison, de cause et de concession pour argumenter]
\textbf{Objectif :} Construire des phrases complexes correctes pour comparer deux économies, expliquer les causes et nuancer l'argumentation.\par
\textbf{3 structures indispensables (schéma + emploi + exceptions) :}
\begin{enumerate}
    \item \textbf{Comparaison de supériorité : \zh{A 比 B + Adj.}} (A est plus [Adj] que B).
        \begin{itemize}
            \item \emph{Degré :} \zh{一点儿} / \zh{一些} (un peu), \zh{得多} / \zh{多了} (beaucoup) se placent \emph{après} l'adjectif.
            \item \emph{Infériorité :} \zh{A 没有 B + Adj.} (A n'est pas aussi [Adj] que B).
            \item \textbf{Interdits :} pas de \zh{很} ni de \zh{更} avec \zh{比} ; pas de \zh{不} devant l'adjectif avec \zh{没有}.
            \item \emph{Exemples :} \zh{中国经济比喀麦隆经济大得多。} (Zhōngguó jīngjì bǐ Kāmàilóng jīngjì dà de duō.) « L'économie chinoise est bien plus grande que celle du Cameroun. » / \zh{喀麦隆经济没有中国经济大。} (Kāmàilóng jīngjì méiyǒu Zhōngguó jīngjì dà.) « L'économie camerounaise n'est pas aussi grande que celle de la Chine. »
        \end{itemize}
    \item \textbf{Cause $\rightarrow$ Conséquence : \zh{因为……，所以……}} (formel : \zh{由于……，因此……}).
        \begin{itemize}
            \item \zh{因为} en tête de la proposition 1, \zh{所以} en tête de la proposition 2. Virgule obligatoire entre les deux.
            \item \emph{Exemple :} \zh{因为基础设施差，所以投资少。} (Yīnwèi jīchǔ shèshī chà, suǒyǐ tóuzī shǎo.) « Parce que les infrastructures sont mauvaises, les investissements sont faibles. »
        \end{itemize}
    \item \textbf{Concession : \zh{虽然……，但是……}} (bien que..., mais...).
        \begin{itemize}
            \item \zh{虽然} en tête de la proposition 1 ; \zh{但是} (ou \zh{却} après le sujet) en tête de la proposition 2.
            \item \textbf{Interdit :} ne pas combiner \zh{虽然……但是……} avec \zh{因为……所以……} dans la même phrase (redondance logique).
            \item \emph{Exemple :} \zh{虽然喀麦隆有石油，但是炼油厂不够。} (Suīrán Kāmàilóng yǒu shíyóu, dànshì liànyóuchǎng bù gòu.) « Bien que le Cameroun ait du pétrole, les raffineries ne suffisent pas. »
        \end{itemize}
\end{enumerate}
\textbf{Connecteurs d'enchaînement pour le paragraphe :} \zh{首先} (premièrement), \zh{其次} (deuxièmement), \zh{此外} (de plus), \zh{最后} (enfin), \zh{因此} (c'est pourquoi), \zh{总之} (en résumé).
\end{methode}

\begin{exoresolu}[Transformation et production de phrases complexes (comparaison / cause / concession)]
\textbf{Énoncé :}
\begin{enumerate}
    \item Transformez les phrases en utilisant \zh{比} (supériorité) puis \zh{没有} (infériorité).
    \begin{itemize}
        \item[a.] 中国经济大，喀麦隆经济小。
        \item[b.] 泰国工业发达，喀麦隆工业不发达。
    \end{itemize}
    \item Reliez les deux propositions par \zh{因为……所以……}.
    \begin{itemize}
        \item[a.] 政府投资教育少。劳动力技能不强。
        \item[b.] 基础设施不好。运输成本高。
    \end{itemize}
    \item Écrivez une phrase avec \zh{虽然……但是……} pour exprimer : « Le Cameroun a du pétrole, mais il importe de l'essence. »
    \item \textbf{Production guidée :} En utilisant \zh{首先}、\zh{其次}、\zh{最后}、\zh{总之}, rédigez un court paragraphe (4-5 phrases) comparant l'agriculture du Cameroun et de la Chine. Vocabulaire imposé : \zh{农业} (nóngyè), \zh{机械化} (jīxièhuà), \zh{产量} (chǎnliàng), \zh{人口} (rénkǒu).
\end{enumerate}
\tcblower
\textbf{Solution détaillée :}
\begin{enumerate}
    \item \textbf{Comparaison :}
    \begin{itemize}
        \item[a.] \zh{中国经济比喀麦隆经济大得多。} (Zhōngguó jīngjì bǐ Kāmàilóng jīngjì dà de duō.) — \zh{喀麦隆经济没有中国经济大。} (Kāmàilóng jīngjì méiyǒu Zhōngguó jīngjì dà.)
        \item[b.] \zh{泰国工业比喀麦隆工业发达。} (Tàiguó gōngyè bǐ Kāmàilóng gōngyè fādá.) — \zh{喀麦隆工业没有泰国工业发达。} (Kāmàilóng gōngyè méiyǒu Tàiguó gōngyè fādá.)
    \end{itemize}
    \item \textbf{Cause-conséquence :}
    \begin{itemize}
        \item[a.] \zh{因为政府在教育上投资少，所以劳动力技能不强。} (Yīnwèi zhèngfǔ zài jiàoyù shang tóuzī shǎo, suǒyǐ láodònglì jìnéng bù qiáng.)
        \item[b.] \zh{因为基础设施不好，所以运输成本很高。} (Yīnwèi jīchǔ shèshī bù hǎo, suǒyǐ yùnshū chéngběn hěn gāo.)
    \end{itemize}
    \item \textbf{Concession :}
    \zh{虽然喀麦隆有石油，但是要进口汽油。} (Suīrán Kāmàilóng yǒu shíyóu, dànshì yào jìnkǒu qìyóu.) « Bien que le Cameroun ait du pétrole, il doit importer de l'essence. »
    \item \textbf{Paragraphe modèle :}\\
    \zh{首先，中国农业的机械化程度比喀麦隆高得多。其次，因为中国在农业技术上投资很多，所以粮食产量很大。最后，虽然喀麦隆土地肥沃，但是缺少机械和资金。总之，喀麦隆应该学习中国的农业现代化经验。}\\
    \emph{Shǒuxiān, Zhōngguó nóngyè de jīxièhuà chéngdù bǐ Kāmàilóng gāo de duō. Qícì, yīnwèi Zhōngguó zài nóngyè jìshù shang tóuzī hěn duō, suǒyǐ liángshi chǎnliàng hěn dà. Zuìhòu, suīrán Kāmàilóng tǔdì féiwò, dànshì quēshǎo jīxiè hé zījīn. Zǒngzhī, Kāmàilóng yīnggāi xuéxí Zhōngguó de nóngyè xiàndàihuà jīngyàn.}\\
    « Premièrement, le degré de mécanisation de l'agriculture chinoise est bien plus élevé que celui du Cameroun. Deuxièmement, parce que la Chine investit beaucoup dans les technologies agricoles, sa production céréalière est très importante. Enfin, bien que le Cameroun ait des terres fertiles, il manque de machines et de capitaux. En résumé, le Cameroun devrait s'inspirer de l'expérience chinoise de modernisation agricole. »\\
    \textbf{Analyse grammaticale :}
    \begin{itemize}
        \item \zh{程度……高} (niveau... élevé) : collocation correcte.
        \item \zh{粮食产量} (production céréalière) : vocabulaire précis.
        \item \zh{土地肥沃} (terres fertiles) : argument de concession réaliste.
        \item \zh{学习……的经验} (s'inspirer de l'expérience de...) : structure finale proactive.
    \end{itemize}
\end{enumerate}
\end{exoresolu}

\begin{methode}[Traduire du français vers le chinois (thème) : économie et développement]
\textbf{Objectif :} Passer du français (articles, genres, place libre des compléments) au chinois (SVO strict, pas d'articles, particules d'aspect, classificateurs, ordre fixe des compléments).\par
\textbf{Démarche en 5 étapes :}
\begin{enumerate}
    \item \textbf{Analyse de la phrase source :} Identifier sujet, verbe (temps), compléments (objet, lieu, temps, manière), négations, comparatifs.
    \item \textbf{Transposition lexicale :} Choisir le bon mot chinois (ex. : « développement » économique $\rightarrow$ \zh{发展}, pas \zh{生长} qui concerne les êtres vivants). Vérifier le \cle{classificateur} si nombre + nom (\zh{一个项目}, \zh{三条路}).
    \item \textbf{Réordonnancement syntaxique (piège n°1) :} Placer les compléments dans l'ordre chinois : \cle{Sujet + Temps + Lieu + Verbe + Objet}.
        \begin{itemize}
            \item Fr. : « Il a investi \textbf{beaucoup d'argent} \textbf{au Cameroun} \textbf{l'année dernière}. »
            \item Zh. : \zh{他去年在喀麦隆投资了很多钱。} (Tā qùnián zài Kāmàilóng tóuzīle hěn duō qián.) — Temps \zh{去年} avant lieu \zh{在喀麦隆} ; \zh{了} après le verbe (action accomplie) ; \zh{很多钱} après le verbe.
        \end{itemize}
    \item \textbf{Gestion des mots-outils (piège n°2) :}
        \begin{itemize}
            \item \zh{了} (action accomplie / changement d'état) vs \zh{过} (expérience vécue).
            \item \zh{的} (épithète), \zh{地} (complément circonstanciel de manière), \zh{得} (complément de degré).
            \item Comparatif : \zh{比} / \zh{没有} (et non \zh{不比} pour « pas aussi... que »).
        \end{itemize}
    \item \textbf{Relecture « tons et caractères » :} Vérifier les tons (notamment \zh{不} bù $\rightarrow$ bú devant un 4\up{e} ton), les homophones, l'orthographe des caractères.
\end{enumerate}
\textbf{Astuce examen :} Traduire par « blocs de sens » : découper $\rightarrow$ transposer $\rightarrow$ réordonner $\rightarrow$ relire.
\end{methode}

\begin{exoresolu}[Thème : 5 phrases sur l'économie camerounaise (français $\rightarrow$ chinois)]
\textbf{Énoncé :} Traduisez en chinois (caractères + pinyin).
\begin{enumerate}
    \item Le Cameroun a de riches ressources naturelles. (\zh{有} / \zh{丰富} / \zh{自然资源})
    \item L'année dernière, le gouvernement a investi beaucoup d'argent dans les infrastructures. (\zh{去年} / \zh{政府} / \zh{投资} / \zh{基础设施})
    \item L'industrialisation du Cameroun n'est pas aussi rapide que celle de la Chine. (\zh{工业化} / \zh{快} / \zh{没有})
    \item Bien que le marché soit grand, le pouvoir d'achat est faible. (\zh{虽然} / \zh{市场} / \zh{但是} / \zh{购买力} / \zh{低})
    \item Si l'on veut développer l'économie, il faut améliorer l'éducation. (\zh{如果} / \zh{想} / \zh{发展} / \zh{经济} / \zh{必须} / \zh{改善} / \zh{教育})
\end{enumerate}
\tcblower
\textbf{Solution détaillée (avec justifications) :}
\begin{enumerate}
    \item \zh{喀麦隆有丰富的自然资源。} (Kāmàilóng yǒu fēngfù de zìrán zīyuán.)\\
    \emph{Justification :} \zh{有} (avoir) est un verbe transitif direct. \zh{丰富的} : adjectif à deux syllabes + \zh{的} obligatoire devant le nom.
    \item \zh{去年，政府在基础设施建设上投资了很多钱。} (Qùnián, zhèngfǔ zài jīchǔ shèshī jiànshè shang tóuzīle hěn duō qián.)\\
    \emph{Justification :} \zh{去年} (temps) en tête. Le complément de lieu/domaine \zh{在……上} se place \emph{avant} le verbe. \zh{投资} + \zh{了} (action accomplie et datée). \zh{很多钱} après le verbe. \textbf{Piège évité :} ne pas dire \zh{投资了很多钱在基础设施} (ordre français calqué).
    \item \zh{喀麦隆的工业化没有中国的快。} (Kāmàilóng de gōngyèhuà méiyǒu Zhōngguó de kuài.)\\
    \emph{Justification :} Comparaison entre éléments de même nature, nominalisés par \zh{的} (« celle de la Chine » = \zh{中国的}). Structure d'infériorité \zh{没有……快} (« pas aussi rapide que »).
    \item \zh{虽然市场很大，但是购买力很低。} (Suīrán shìchǎng hěn dà, dànshì gòumǎilì hěn dī.)\\
    \emph{Justification :} \zh{虽然} en tête de la 1\up{re} proposition, \zh{但是} en tête de la 2\up{e}. \zh{购买力} (pouvoir d'achat) ; \zh{低} (faible, bas) est l'adjectif correct pour un niveau. \zh{很} est requis devant l'adjectif monosyllabique prédicat (valeur neutre, non comparative).
    \item \zh{如果想发展经济，就必须改善教育。} (Rúguǒ xiǎng fāzhǎn jīngjì, jiù bìxū gǎishàn jiàoyù.)\\
    \emph{Justification :} \zh{如果……就……} (si... alors...) : structure conditionnelle standard. Le sujet générique « on » est omis en chinois. \zh{想} + verbe (avoir l'intention de) ; \zh{必须} + verbe (il faut).
\end{enumerate}
\end{exoresolu}

\begin{methode}[Traduire du chinois vers le français (version) : précision terminologique et syntaxe française]
\textbf{Objectif :} Rendre le sens exact du chinois en français correct, en gérant l'absence de conjugaison (les particules d'aspect), la topicalisation et les structures spécifiques.\par
\textbf{Démarche en 4 étapes :}
\begin{enumerate}
    \item \textbf{Segmentation (découper en mots) :} Le chinois s'écrit sans espaces. Repérer les mots : \zh{喀麦隆 / 的 / 经济 / 发展 / 很快} (Cameroun / de / économie / développer / très vite).
    \item \textbf{Analyse grammaticale (qui fait quoi ?) :} Identifier le \cle{sujet} (souvent en tête), le \cle{verbe}, les \cle{particules} (\zh{了、过、着、的、地、得}). Déterminer le \cle{temps/aspect} : \zh{了} (accompli), \zh{过} (expérience), \zh{正在}/\zh{在} (progressif), \zh{将}/\zh{要}/\zh{会} (futur).
    \item \textbf{Transposition lexicale contextuelle :} Choisir le mot français précis selon le contexte économique.
        \begin{itemize}
            \item \zh{发展} : développer (v.t.), se développer (v.i.), développement (n.).
            \item \zh{投资} : investir (v.), investissement (n.).
            \item \zh{基础设施} : infrastructures (au pluriel en français).
            \item \zh{人力资源} : ressources humaines / main-d'œuvre.
        \end{itemize}
    \item \textbf{Rédaction française :} Reconstruire une phrase française naturelle (SVO, accords, articles). Rendre la \cle{topicalisation} chinoise par « Quant à... », « En ce qui concerne... » si nécessaire. \textbf{Attention :} \zh{被} (bèi) = voix passive ; \zh{把} (bǎ) = antéposition du complément d'objet (insistance sur le résultat de l'action).
\end{enumerate}
\textbf{Règle d'or :} Ne jamais traduire mot à mot. Traduire la \cle{phrase}, pas les \cle{mots}.
\end{methode}

\begin{textebox}[Texte source pour la version]
近年来，喀麦隆政府非常重视基础设施建设，已经修建了很多高速公路和大桥。\\
Jìnniánlái, Kāmàilóng zhèngfǔ fēicháng zhòngshì jīchǔ shèshī jiànshè, yǐjing xiūjiànle hěn duō gāosù gōnglù hé dà qiáo.\\
但是，由于资金不足，电力供应还是一个大问题。\\
Dànshì, yóuyú zījīn bùzú, diànlì gōngyìng háishì yí ge dà wèntí.\\
如果电力问题解决了，工厂的生产效率就会提高，外资也会增加。\\
Rúguǒ diànlì wèntí jiějuéle, gōngchǎng de shēngchǎn xiàolǜ jiù huì tígāo, wàizī yě huì zēngjiā.\\
希望将来喀麦隆能成为中非的经济枢纽。\\
Xīwàng jiānglái Kāmàilóng néng chéngwéi Zhōng-Fēi de jīngjì shūniǔ.
\end{textebox}

\begin{exoresolu}[Version : texte économique chinois $\rightarrow$ français]
\textbf{Énoncé :} Traduisez en français correct et naturel le texte source ci-dessus.
\tcblower
\textbf{Solution détaillée (segmentation $\rightarrow$ analyse $\rightarrow$ traduction) :}
\begin{enumerate}
    \item \zh{近年来} (jìnniánlái) : « ces dernières années » (marqueur temporel $\rightarrow$ présent ou passé composé en français).
    \item \zh{喀麦隆政府非常重视基础设施建设} : sujet \zh{政府} ; \zh{重视} = « accorder une grande importance à » (et non « regarder ») ; objet : \zh{基础设施建设} « la construction d'infrastructures ».
    \item \zh{已经修建了} : \zh{已经} + verbe + \zh{了} = action achevée $\rightarrow$ passé composé (« a déjà construit »). \zh{高速公路} (autoroutes), \zh{大桥} (grands ponts).
    \item \zh{由于资金不足} : « en raison d'un manque de financement » (\zh{由于} = « en raison de », registre formel). \zh{电力供应还是一个大问题} : \zh{还是} = « reste, demeure » ; \zh{电力供应} = « l'approvisionnement en électricité ».
    \item \zh{如果……解决了，……就会提高，……也会增加} : conditionnelle \zh{如果……就……} ; \zh{了} marque ici l'accomplissement hypothétique (« si le problème est résolu ») ; \zh{就会} = « alors... (conséquence certaine) » ; \zh{生产效率} = « l'efficacité de la production » ; \zh{外资} = « les investissements étrangers ».
    \item \zh{希望将来喀麦隆能成为中非的经济枢纽} : \zh{希望} = « on espère que » ; \zh{将来} = « à l'avenir » ; \zh{成为} = « devenir » ; \zh{中非} = « sino-africain / entre la Chine et l'Afrique » (\zh{中} = Chine, \zh{非} = Afrique) ; \zh{枢纽} = « plaque tournante, hub ».
\end{enumerate}
\textbf{Traduction finale proposée :}\par
\emph{« Ces dernières années, le gouvernement camerounais accorde une grande importance à la construction d'infrastructures : il a déjà construit de nombreuses autoroutes et de grands ponts. Cependant, en raison d'un manque de financement, l'approvisionnement en électricité reste un problème majeur. Si le problème de l'électricité est résolu, l'efficacité de la production des usines s'améliorera et les investissements étrangers augmenteront également. On espère qu'à l'avenir le Cameroun pourra devenir une plaque tournante économique entre la Chine et l'Afrique. »}\par
\textbf{Points de vigilance pour l'examen :}
\begin{itemize}
    \item \zh{重视} $\neq$ « regarder » mais « accorder de l'importance à ».
    \item \zh{已经……了} $\rightarrow$ passé composé (« a déjà construit »), pas d'imparfait.
    \item \zh{还是} dans un contexte de problème persistant $\rightarrow$ « reste / demeure / est encore ».
    \item \zh{中非} = Chine-Afrique, et non « Afrique centrale » (qui serait \zh{中非} uniquement comme nom propre du pays ; ici le contexte économique Chine-Cameroun impose « sino-africain »).
\end{itemize}
\end{exoresolu}

\begin{methode}[Rédiger un texte argumenté court (production écrite guidée : 80-100 caractères)]
\textbf{Objectif :} Produire un paragraphe cohérent, structuré, lexicalement riche et grammaticalement correct sur un sujet économique camerounais, en réinvestissant le modèle, le vocabulaire, les structures et les caractères de la leçon.\par
\textbf{Démarche en 6 étapes (processus d'écriture) :}
\begin{enumerate}
    \item \textbf{Analyse du sujet :} Identifier le \cle{type de texte} (exposé comparatif, opinion, rapport simple), la \cle{contrainte de longueur} (ex. : « environ 80 caractères ») et le \cle{thème imposé} (ex. : « Les atouts et les défis de l'économie camerounaise »).
    \item \textbf{Remue-méninges et plan (5 min) :} Lister idées et mots-clés (\zh{资源、农业、工业、基础设施、教育、投资、挑战、优势}). Structurer en 3 parties : \textbf{Introduction} (thèse) $\rightarrow$ \textbf{Développement} (2 arguments : un atout, un défi, avec \zh{首先}、\zh{但是}) $\rightarrow$ \textbf{Conclusion} (recommandation avec \zh{应该}、\zh{必须}).
    \item \textbf{Rédaction du brouillon (en chinois directement) :} Suivre le plan. Utiliser les \cle{patrons de phrases} vus en classe (\zh{因为……所以……}, \zh{虽然……但是……}, \zh{比}/\zh{没有}). Ne pas traduire depuis le français !
    \item \textbf{Vérification grammaticale (auto-correction) :} Liste de contrôle :
        \begin{itemize}
            \item Ordre des mots : Sujet + Temps + Lieu + Verbe ?
            \item \zh{的}/\zh{地}/\zh{得} correctement employés ?
            \item \zh{了}/\zh{过} selon le sens (accompli vs expérience) ?
            \item Classificateurs si nombres ?
            \item Comparatifs : \zh{比} + Adj. (sans \zh{很}) / \zh{没有} + Adj. (sans \zh{不}) ?
        \end{itemize}
    \item \textbf{Calligraphie et ponctuation (copie finale) :} Écrire au propre en style régulier (楷书), caractères bien proportionnés. Ponctuation chinoise pleine chasse (\zh{，。；！？}). Pas de pinyin sur la copie finale (sauf consigne contraire).
    \item \textbf{Relecture finale :} Le texte répond-il au sujet ? Le nombre de caractères est-il respecté ($\pm$ 10 \%) ?
\end{enumerate}
\textbf{Critères d'évaluation (grille simplifiée) :} Contenu et idées (4 pts) + Vocabulaire et structures (4 pts) + Grammaire et précision (4 pts) + Cohérence et connecteurs (4 pts) + Présentation et caractères (4 pts) = /20.
\end{methode}

\begin{exoresolu}[Production écrite : « Les atouts et les défis de l'économie du Cameroun »]
\textbf{Énoncé :} Rédigez un paragraphe d'environ 80-100 caractères chinois sur le thème : \zh{喀麦隆经济的优势与挑战} (Les atouts et les défis de l'économie camerounaise). Utilisez au moins : \zh{首先}、\zh{但是}、\zh{因为……所以……}、\zh{应该}.
\tcblower
\textbf{Solution détaillée (modèle de copie « excellente ») :}\par
\zh{喀麦隆经济有优势，也有挑战。首先，喀麦隆自然资源丰富，农业基础好，土地肥沃。其次，因为人口年轻，所以人力资源潜力很大。但是，基础设施不够完善，工业化程度低，电力供应不稳定。因此，政府应该加大投资，改善基础设施，发展制造业。总之，只要努力，喀麦隆经济一定会发展得更好。}\par
\emph{Kāmàilóng jīngjì yǒu yōushì, yě yǒu tiǎozhàn. Shǒuxiān, Kāmàilóng zìrán zīyuán fēngfù, nóngyè jīchǔ hǎo, tǔdì féiwò. Qícì, yīnwèi rénkǒu niánqīng, suǒyǐ rénlì zīyuán qiánlì hěn dà. Dànshì, jīchǔ shèshī bù gòu wánshàn, gōngyèhuà chéngdù dī, diànlì gōngyìng bù wěndìng. Yīncǐ, zhèngfǔ yīnggāi jiādà tóuzī, gǎishàn jīchǔ shèshī, fāzhǎn zhìzàoyè. Zǒngzhī, zhǐyào nǔlì, Kāmàilóng jīngjì yídìng huì fāzhǎn de gèng hǎo.}\par
« L'économie camerounaise a des atouts, mais aussi des défis. Premièrement, le Cameroun possède de riches ressources naturelles, une bonne base agricole et des terres fertiles. Deuxièmement, parce que sa population est jeune, ses ressources humaines ont un grand potentiel. Cependant, les infrastructures ne sont pas assez abouties, le degré d'industrialisation est faible et l'approvisionnement en électricité est instable. C'est pourquoi le gouvernement devrait augmenter les investissements, améliorer les infrastructures et développer l'industrie manufacturière. En résumé, pourvu qu'il fasse des efforts, l'économie du Cameroun se développera certainement mieux. »\par
\textbf{Analyse du modèle (pourquoi cela fonctionne) :}
\begin{itemize}
    \item \textbf{Structure :} Thèse (\zh{有优势，也有挑战}) $\rightarrow$ Atouts (\zh{首先}... ressources ; \zh{其次}... démographie) $\rightarrow$ Défis (\zh{但是}... infrastructures, industrie, électricité) $\rightarrow$ Solutions (\zh{因此}... \zh{应该}...) $\rightarrow$ Conclusion optimiste (\zh{总之}...).
    \item \textbf{Grammaire validée :}
        \begin{itemize}
            \item \zh{自然资源丰富} (sujet + adjectif prédicat, sans \zh{很} car énoncé factuel).
            \item \zh{因为人口年轻，所以人力资源潜力很大} (cause-conséquence parfaite).
            \item \zh{基础设施不够完善} (\zh{不够} + Adj. = « pas assez »).
            \item \zh{政府应该加大投资} (\zh{应该} + verbe = conseil, obligation morale).
            \item \zh{发展得更好} (verbe + \zh{得} + complément de degré : « se développer mieux »).
        \end{itemize}
    \item \textbf{Vocabulaire :} \zh{优势} (atout), \zh{挑战} (défi), \zh{肥沃} (fertile), \zh{潜力} (potentiel), \zh{完善} (abouti), \zh{程度} (degré), \zh{稳定} (stable), \zh{加大} (augmenter), \zh{制造业} (industrie manufacturière).
    \item \textbf{Longueur :} environ 100 caractères (dans la fourchette).
    \item \textbf{Connecteurs :} \zh{首先}、\zh{其次}、\zh{但是}、\zh{因此}、\zh{总之} (5 connecteurs logiques bien employés).
\end{itemize}
\textbf{Variante « niveau passable » (60-70 caractères) :} moins de détails, structures plus simples (\zh{有……也有……}, \zh{因为……所以……}, \zh{应该……}).
\end{exoresolu}

\begin{attention}[Les 5 erreurs qui coûtent des points à l'examen (spécial économie / expression écrite)]
\begin{enumerate}
    \item \textbf{Tons et homophones (perte de sens) :} Confondre \zh{投资} (tóuzī, investir) et \zh{头子} (tóuzi, chef de bande) ; \zh{基础} (jīchǔ, base) et \zh{机器} (jīqì, machine) ; \zh{产量} (chǎnliàng, production) et \zh{车辆} (chēliàng, véhicules). \textbf{Réflexe :} répéter le vocabulaire \emph{à voix haute} avec les tons ; écrire le pinyin \emph{avec les tons} sur le brouillon.
    \item \textbf{Caractères « jumeaux » (perte de points en écriture) :}
        \begin{itemize}
            \item \zh{资} (zī, richesse, radical \zh{贝}) vs \zh{咨} (zī, consulter, radical \zh{口}).
            \item \zh{构} (gòu, structure, radical \zh{木}) vs \zh{沟} (gōu, fossé, radical \zh{氵}).
            \item \zh{经} (jīng, économie, radical \zh{纟}) vs \zh{轻} (qīng, léger, radical \zh{车}).
            \item \zh{产} (chǎn, produire) vs \zh{严} (yán, strict).
        \end{itemize}
        \textbf{Règle :} le radical donne le sens. Argent, économie $\rightarrow$ \zh{贝} (\zh{资、贸、财}) ; bois, construction $\rightarrow$ \zh{木} (\zh{构、桥、楼}) ; eau $\rightarrow$ \zh{氵} (\zh{源、沟、油}).
    \item \textbf{Ordre des mots : compléments de temps et de lieu :}
        \textbf{Erreur :} \zh{他投资了很多钱去年在喀麦隆。} (ordre français calqué) $\rightarrow$ \textbf{Correct :} \zh{他去年在喀麦隆投资了很多钱。} (Tā qùnián zài Kāmàilóng tóuzīle hěn duō qián.) \textbf{Règle d'or :} \cle{Sujet + Temps + Lieu + Verbe + Objet}.
    \item \textbf{Mots-outils \zh{的}、\zh{地}、\zh{得} et \zh{了}、\zh{过} :}
        \begin{itemize}
            \item \zh{快速的发展} (Adj. + \zh{的} + nom : « un développement rapide ») vs \zh{快速地发展} (Adv. + \zh{地} + verbe : « se développer rapidement ») vs \zh{发展得很快} (verbe + \zh{得} + degré : « se développe très vite »).
            \item \zh{了} vs \zh{过} : \zh{我去过喀麦隆。} (Wǒ qùguo Kāmàilóng. « Je suis déjà allé au Cameroun », expérience) vs \zh{我去年去了喀麦隆。} (Wǒ qùnián qùle Kāmàilóng. « Je suis allé au Cameroun l'année dernière », action datée et accomplie).
        \end{itemize}
    \item \textbf{Calques du français et comparatifs :}
        \begin{itemize}
            \item \textbf{Erreur :} \zh{中国经济比喀麦隆经济更大。} (\zh{比} + \zh{更} = redondance). \textbf{Correct :} \zh{中国经济比喀麦隆经济大。} ou \zh{中国经济比喀麦隆经济大得多。}
            \item \textbf{Erreur :} \zh{喀麦隆没有泰国经济发达。} (comparaison asymétrique : un pays comparé à une économie). \textbf{Correct :} \zh{喀麦隆的经济没有泰国的经济发达。} ou \zh{喀麦隆没有泰国发达。} (comparer un pays à un pays, une économie à une économie).
            \item \textbf{Erreur :} combiner \zh{因为……所以……} et \zh{虽然……但是……} dans la même phrase. \textbf{Règle :} choisir l'une ou l'autre selon la logique (cause ou concession).
        \end{itemize}
\end{enumerate}
\textbf{Conseil final :} Relisez votre copie en vous demandant : « Chaque caractère a-t-il son sens précis ? Chaque particule a-t-elle sa fonction grammaticale ? L'ordre des mots suit-il la logique chinoise (Sujet + Temps + Lieu + Verbe) ? »
\end{attention}

\newpage

\coursec{Exercices}

\courssub{Je vérifie mes connaissances}

\begin{exercice}[QCM : Vocabulaire et structures clés]
Chaque question ne comporte qu'une seule réponse correcte. Choisis la bonne proposition.
\begin{enumerate}
    \item Comment dit-on « économie » en chinois ?
    \begin{itemize}
        \item A. 经济 (jīngjì)
        \item B. 经纪 (jīngjì)
        \item C. 京基 (jīngjī)
    \end{itemize}
    \item Quelle structure exprime correctement la cause et la conséquence : « Parce que le Cameroun a du pétrole, son économie se développe » ?
    \begin{itemize}
        \item A. 喀麦隆有石油，\textbf{但是} 经济发展。
        \item B. \textbf{因为} 喀麦隆有石油，\textbf{所以} 经济发展。
        \item C. 喀麦隆有石油，\textbf{而且} 经济发展。
    \end{itemize}
    \item Pour comparer le développement du Cameroun et de la Thaïlande (« Le Cameroun se développe plus vite que la Thaïlande »), on utilise :
    \begin{itemize}
        \item A. 喀麦隆 \textbf{比} 泰国 发展 快。
        \item B. 喀麦隆 \textbf{跟} 泰国 发展 一样 快。
        \item C. 喀麦隆 \textbf{没有} 泰国 发展 快。
    \end{itemize}
    \item Quel mot signifie « produit agricole » ?
    \begin{itemize}
        \item A. 石油 (shíyóu)
        \item B. 农产品 (nóngchǎnpǐn)
        \item C. 基础设施 (jīchǔ shèshī)
    \end{itemize}
\end{enumerate}
\end{exercice}

\begin{exercice}[Vrai ou Faux : Justifiez en français]
Lis le court texte ci-dessous et détermine si les affirmations sont vraies (V) ou fausses (F). Justifie ta réponse en citant le texte.
\begin{textebox}[Texte : L'économie du Cameroun]
喀麦隆 有 很多 石油 和 可可。\\
因为 自然 资源 丰富，所以 经济 发展 得 比较 快。\\
但是 基础设施 还 不够 好。\\
我觉得 将来 会 越来越 好。
\end{textebox}
\begin{enumerate}
    \item Le Cameroun possède du pétrole et du cacao.
    \item L'économie se développe vite parce que la population est nombreuse.
    \item Les infrastructures sont déjà parfaites.
    \item L'auteur est optimiste pour l'avenir.
\end{enumerate}
\end{exercice}

\begin{exercice}[Vocabulaire : Complétez le tableau]
Complète les cases vides du tableau ci-dessous (caractères, pinyin, nature ou traduction).
\begin{center}
\begin{tabularx}{\linewidth}{|X|X|X|X|}
\hline
\rowcolor{popblueL} \textbf{Caractères} & \textbf{Pinyin} & \textbf{Nature} & \textbf{Traduction} \\
\hline
经济 & jīngjì & \cle{nom} &  \\
\hline
 & fāzhǎn & verbe &  \\
\hline
石油 &  & \cle{nom} & pétrole \\
\hline
 & kěkě & \cle{nom} & cacao \\
\hline
农产品 &  & \cle{nom} &  \\
\hline
丰富 & fēngfù &  & riche, abondant \\
\hline
基础设施 &  & \cle{nom} & infrastructures \\
\hline
投资 & tóuzī & verbe / nom &  \\
\hline
\end{tabularx}
\end{center}
\end{exercice}

\begin{exercice}[Association Caractères / Pinyin]
Relie chaque caractère (ou mot) à son pinyin correct (tons inclus).
\begin{center}
\begin{tabular}{llll}
\textbf{Caractères} & & \textbf{Pinyin} & \\
\hline
1. 经 & \hspace{3cm} & a. zhǎn \\
2. 济 & & b. jì \\
3. 展 & & c. yóu \\
4. 油 & & d. jīng \\
5. 产 & & e. chǎn \\
6. 品 & & f. pǐn \\
\hline
\end{tabular}
\end{center}
\emph{Écris les paires correspondantes (ex: 1-d).}
\end{exercice}

\courssub{Je m'entraîne}

\begin{exercice}[Thème : Traduisez en chinois (caractères + pinyin)]
Utilise le vocabulaire de la leçon (比, 因为...所以, 越来越, 觉得, 资源, 丰富).
\begin{enumerate}
    \item Le Cameroun a beaucoup de ressources naturelles.
    \item Parce que les produits agricoles sont abondants, l'économie se développe de plus en plus vite.
    \item La Thaïlande a plus d'usines, mais le Cameroun a plus de pétrole.
    \item Je pense que l'avenir du Cameroun sera de mieux en mieux.
    \item Les infrastructures ne sont pas encore très bonnes.
\end{enumerate}
\end{exercice}

\begin{exercice}[Transformation de phrases]
Réécris chaque phrase en utilisant la structure indiquée entre parenthèses. Écris en caractères et en pinyin.
\begin{enumerate}
    \item 喀麦隆有石油。喀麦隆有可可。 (用 \textbf{而且})
    \item 喀麦隆经济发展快。泰国经济发展慢。 (用 \textbf{比})
    \item 因为有农产品，所以经济好。 (用 \textbf{越来越} pour exprimer l'évolution)
    \item 基础设施不好。 (用 \textbf{还} pour nuancer : « pas encore bon »)
    \item 我觉得喀麦隆会发展。 (Remplace \textbf{觉得} par \textbf{认为}  et ajoute \textbf{将来})
\end{enumerate}
\end{exercice}

\begin{exercice}[Texte à trous : Connecteurs et mots-outils]
Complète le texte ci-dessous avec les mots manquants choisis dans la liste : \textbf{因为、所以、但是、而且、比、越来越}。
\begin{textebox}[Pourquoi l'économie camerounaise se développe-t-elle ?]
喀麦隆 自然 资源 很 丰富，\_\_\_\_\_\_ 经济 发展 得 快。\\
\_\_\_\_\_\_ 有 石油，\_\_\_\_\_\_ 有 很多 农产品，\_\_\_\_\_\_ 有 木材。\\
\_\_\_\_\_\_ 基础设施 还 不够 完善。\\
如果 投资 多，将来 会 \_\_\_\_\_\_ 好。
\end{textebox}
\end{exercice}

\begin{exercice}[Remise en ordre : Mots → Phrases]
Remets les mots dans l'ordre pour former des phrases correctes. Écris le résultat en caractères et en pinyin.
\begin{enumerate}
    \item 经济 / 喀麦隆 / 发展 / 比 / 泰国 / 快
    \item 所以 / 石油 / 有 / 因为 / 经济 / 好 / 喀麦隆
    \item 觉得 / 我 / 将来 / 会 / 越来越 / 好 / 国家 / 这个
    \item 产品 / 农 / 丰富 / 很 / 喀麦隆 / 的
    \item 不 / 基础设施 / 够 / 还 / 好 / 是
\end{enumerate}
\end{exercice}

\courssub{Je me prépare au Bac}

\begin{exercice}[Compréhension de l'écrit et langue (Type Bac)]
Lis le texte original suivant, puis réponds aux questions en français.
\begin{textebox}[Article : 喀麦隆与泰国的经济对比]
喀麦隆 和 泰国 都 是 发展 中国家。\\
喀麦隆 的 优势 是 自然 资源 非常 丰富。有 石油、天然气、可可、咖啡 和 木材。农业 占 经济 的 大 部分。\\
泰国 的 优势 是 工业 和 旅游 业 比较 发达。汽车 零件、电子 产品 和 大米 出口 很多。\\
\cle{因为} 喀麦隆 的 基础设施 \cle{不如} 泰国 好，\cle{所以} 外资 少 一些。\\
\cle{但是} 现在 政府 \cle{投资} 修路、建厂、办学，经济 增长 速度 \cle{越来越} 快。\\
专家 说：``喀麦隆 有 潜力，将来 一定 能 赶上 来。''
\end{textebox}
\textbf{Questions de compréhension (en français) :}
\begin{enumerate}
    \item Cite trois ressources naturelles du Cameroun mentionnées dans le texte.
    \item Quels sont les deux piliers de l'économie thaïlandaise selon le texte ?
    \item Pourquoi les investissements étrangers (外资) sont-ils moins nombreux au Cameroun qu'en Thaïlande ? Cite la structure de cause à effet utilisée.
    \item Quel connecteur marque l'opposition au 4\textsuperscript{e} paragraphe ? Quelle mesure le gouvernement prend-il pour améliorer la situation ?
    \item Traduire en français la dernière phrase (citation de l'expert).
\end{enumerate}
\textbf{Questions de langue :}
\begin{enumerate}
    \setcounter{enumi}{5}
    \item Explique la différence de sens entre \zh{不如} (bùrú) et \zh{比} (bǐ) dans une phrase de comparaison. Donne un exemple pour chacun.
    \item Identifie la structure grammaticale de \zh{越来越快} et donne sa traduction.
    \item Trouve dans le texte le mot qui correspond à « investir / investissement » et écris son pinyin.
    \item Réécris la phrase « 喀麦隆的基础设施不如泰国好 » en utilisant \zh{比} (le Cameroun est le sujet).
\end{enumerate}
\end{exercice}

\begin{exercice}[Thème et Version (Type Bac)]
\textbf{Partie A -- Thème (Français $\rightarrow$ Chinois) :}
Traduis la phrase suivante en chinois (caractères simplifiés + pinyin avec tons).
\begin{quote}
« Le Cameroun a beaucoup de produits agricoles, donc son économie se développe de plus en plus vite. »
\end{quote}
\textbf{Partie B -- Version (Chinois $\rightarrow$ Français) :}
Traduis la phrase suivante en français naturel.
\begin{quote}
因为喀麦隆有很多石油和可可，所以经济越来越好。
\end{quote}
\end{exercice}

\begin{exercice}[Expression écrite guidée : Mon pays et son économie (Type Bac)]
\textbf{Consigne :} Rédige un court texte en chinois (caractères simplifiés) d'environ \textbf{80 à 100 caractères} sur le thème : \zh{我的国家和经济} (Mon pays et son économie).
\textbf{Contraintes obligatoires :}
\begin{itemize}
    \item Utilise au moins \textbf{6 phrases} complètes.
    \item Inclus \textbf{au moins une comparaison avec \zh{比}} (ex: comparer le Cameroun à un autre pays ou comparer passé/présent).
    \item Inclus \textbf{ton opinion personnelle avec \zh{我觉得} (ou \zh{我认为})}.
    \item Utilise \textbf{au moins deux connecteurs logiques} parmi : \zh{因为...所以}、\zh{但是}、\zh{而且}、\zh{如果...就}。
    \item Utilise du vocabulaire de la leçon : \zh{资源}、\zh{农产品}、\zh{发展}、\zh{基础设施}、\zh{投资}、\zh{将来}.
\end{itemize}
\textbf{Indication :} Tu peux t'inspirer du modèle de la leçon (texte commenté) pour structurer ton paragraphe (Situation $\rightarrow$ Atouts $\rightarrow$ Défis $\rightarrow$ Actions $\rightarrow$ Avis).
\end{exercice}

\newpage

\coursec{Corrigés des exercices}

\begin{corrige}[Exercice 1 — QCM : Vocabulaire et structures clés]
\begin{enumerate}
    \item \textbf{Réponse A.} \zh{经济} (jīngjì) signifie « économie ». \zh{经纪} (jīngjì) signifie « agent / courtier » : c'est un \cle{homophone} (même pinyin, sens différent). \zh{京基} n'existe pas comme mot courant.
    \item \textbf{Réponse B.} La structure de cause et conséquence est \zh{因为……所以……} (yīnwèi... suǒyǐ...) : « parce que... donc... ». \zh{但是} marque l'opposition, \zh{而且} l'addition : ils ne conviennent pas ici.
    \item \textbf{Réponse A.} La comparaison de supériorité se construit : A + \zh{比} + B + verbe/adjectif. \zh{喀麦隆比泰国发展快。} (Kāmàilóng bǐ Tàiguó fāzhǎn kuài.) « Le Cameroun se développe plus vite que la Thaïlande. » La proposition B exprime l'égalité (\zh{跟……一样}), la C l'infériorité (\zh{没有}).
    \item \textbf{Réponse B.} \zh{农产品} (nóngchǎnpǐn) = « produit agricole ». \zh{石油} = pétrole ; \zh{基础设施} = infrastructures.
\end{enumerate}
\textbf{Points de vigilance :} attention aux homophones \zh{经济}/\zh{经纪} : seuls les caractères permettent de les distinguer à l'écrit. Ne jamais traduire « que » de la comparaison française par un mot séparé : c'est \zh{比} qui porte la comparaison.
\end{corrige}

\begin{corrige}[Exercice 2 — Vrai ou Faux : Justifiez en français]
\begin{enumerate}
    \item \textbf{Vrai.} Le texte dit : \zh{喀麦隆有很多石油和可可。} « Le Cameroun a beaucoup de pétrole et de cacao. »
    \item \textbf{Faux.} Le texte donne une autre cause : \zh{因为自然资源丰富，所以经济发展得比较快。} C'est la richesse des ressources naturelles, et non la population, qui explique le développement.
    \item \textbf{Faux.} Le texte dit : \zh{但是基础设施还不够好。} « Mais les infrastructures ne sont pas encore assez bonnes » (\zh{还不够好} = « pas encore suffisamment bonnes »).
    \item \textbf{Vrai.} Le texte conclut : \zh{我觉得将来会越来越好。} « Je pense qu'à l'avenir cela ira de mieux en mieux » : l'auteur est donc optimiste.
\end{enumerate}
\textbf{Point de vigilance :} \zh{得} dans \zh{发展得比较快} est le mot-outil qui introduit le complément de degré après un verbe : « se développer (de manière) assez vite ».
\end{corrige}

\begin{corrige}[Exercice 3 — Vocabulaire : Complétez le tableau]
\begin{center}
\begin{tabularx}{\linewidth}{|X|X|X|X|}
\hline
\rowcolor{popblueL} \textbf{Caractères} & \textbf{Pinyin} & \textbf{Nature} & \textbf{Traduction} \\
\hline
经济 & jīngjì & \cle{nom} & économie \\
\hline
发展 & fāzhǎn & verbe & (se) développer \\
\hline
石油 & shíyóu & \cle{nom} & pétrole \\
\hline
可可 & kěkě & \cle{nom} & cacao \\
\hline
农产品 & nóngchǎnpǐn & \cle{nom} & produit agricole \\
\hline
丰富 & fēngfù & adjectif & riche, abondant \\
\hline
基础设施 & jīchǔ shèshī & \cle{nom} & infrastructures \\
\hline
投资 & tóuzī & verbe / nom & investir / investissement \\
\hline
\end{tabularx}
\end{center}
\textbf{Points de vigilance :} \zh{丰富} est un adjectif (on peut dire \zh{很丰富}) ; \zh{石油} se prononce shíyóu (2\textsuperscript{e} ton sur les deux syllabes) ; dans \zh{可可}, les deux syllabes portent le 3\textsuperscript{e} ton (kěkě).
\end{corrige}

\begin{corrige}[Exercice 4 — Association Caractères / Pinyin]
Les paires correctes sont :
\begin{itemize}
    \item 1 -- d : \zh{经} = jīng (1\textsuperscript{er} ton)
    \item 2 -- b : \zh{济} = jì (4\textsuperscript{e} ton)
    \item 3 -- a : \zh{展} = zhǎn (3\textsuperscript{e} ton)
    \item 4 -- c : \zh{油} = yóu (2\textsuperscript{e} ton)
    \item 5 -- e : \zh{产} = chǎn (3\textsuperscript{e} ton)
    \item 6 -- f : \zh{品} = pǐn (3\textsuperscript{e} ton)
\end{itemize}
\textbf{Points de vigilance :} ne pas confondre \zh{经} (jīng, sans ton montant) et \zh{济} (jì, ton descendant) ; \zh{油} (yóu) commence par la voyelle « iou » contractée en « -iu » après une consonne (\zh{石油} shíyóu), mais s'écrit « you » en début de syllabe.
\end{corrige}

\begin{corrige}[Exercice 5 — Thème : Traduisez en chinois]
\begin{enumerate}
    \item \zh{喀麦隆有很多自然资源。}\\
    \emph{Kāmàilóng yǒu hěn duō zìrán zīyuán.}\\
    « Le Cameroun a beaucoup de ressources naturelles. »
    \item \zh{因为农产品很丰富，所以经济发展得越来越快。}\\
    \emph{Yīnwèi nóngchǎnpǐn hěn fēngfù, suǒyǐ jīngjì fāzhǎn de yuè lái yuè kuài.}\\
    « Parce que les produits agricoles sont abondants, l'économie se développe de plus en plus vite. »
    \item \zh{泰国的工厂比较多，但是喀麦隆的石油比较多。}\\
    \emph{Tàiguó de gōngchǎng bǐjiào duō, dànshì Kāmàilóng de shíyóu bǐjiào duō.}\\
    « La Thaïlande a plus d'usines, mais le Cameroun a plus de pétrole. »\\
    (On peut aussi écrire : \zh{泰国比喀麦隆工厂多，但是喀麦隆比泰国石油多。})
    \item \zh{我觉得喀麦隆的将来会越来越好。}\\
    \emph{Wǒ juéde Kāmàilóng de jiānglái huì yuè lái yuè hǎo.}\\
    « Je pense que l'avenir du Cameroun sera de mieux en mieux. »
    \item \zh{基础设施还不是很好。}\\
    \emph{Jīchǔ shèshī hái bú shì hěn hǎo.}\\
    « Les infrastructures ne sont pas encore très bonnes. »
\end{enumerate}
\textbf{Points de vigilance :} dans la phrase 2, \zh{得} (et non \zh{的}) introduit le complément de degré après le verbe \zh{发展} ; dans la phrase 4, \zh{会} exprime le futur probable ; « de plus en plus » se rend toujours par \zh{越来越} + adjectif, jamais par une répétition de « plus ».
\end{corrige}

\begin{corrige}[Exercice 6 — Transformation de phrases]
\begin{enumerate}
    \item \zh{喀麦隆有石油，而且有可可。}\\
    \emph{Kāmàilóng yǒu shíyóu, érqiě yǒu kěkě.}\\
    « Le Cameroun a du pétrole, et aussi du cacao. » (\zh{而且} relie deux informations qui s'ajoutent.)
    \item \zh{喀麦隆比泰国经济发展快。} ou mieux : \zh{喀麦隆的经济比泰国发展快。}\\
    \emph{Kāmàilóng de jīngjì bǐ Tàiguó fāzhǎn kuài.}\\
    « L'économie du Cameroun se développe plus vite que celle de la Thaïlande. » (Structure : A + \zh{比} + B + verbe + adjectif.)
    \item \zh{因为有农产品，所以经济越来越好。}\\
    \emph{Yīnwèi yǒu nóngchǎnpǐn, suǒyǐ jīngjì yuè lái yuè hǎo.}\\
    « Parce qu'il y a des produits agricoles, l'économie va de mieux en mieux. »
    \item \zh{基础设施还不好。} ou \zh{基础设施还不够好。}\\
    \emph{Jīchǔ shèshī hái bù hǎo. / Jīchǔ shèshī hái búgòu hǎo.}\\
    « Les infrastructures ne sont pas encore bonnes / pas encore assez bonnes. » (\zh{还} marque ici « pas encore ».)
    \item \zh{我认为喀麦隆将来会发展。}\\
    \emph{Wǒ rènwéi Kāmàilóng jiānglái huì fāzhǎn.}\\
    « Je pense (considère) que le Cameroun se développera à l'avenir. » (\zh{认为} est un peu plus formel que \zh{觉得} ; \zh{将来} se place avant le verbe.)
\end{enumerate}
\textbf{Point de vigilance :} avec \zh{比}, on ne dit jamais « plus » : l'adjectif seul suffit (\zh{快} = « plus vite » dans ce contexte). On peut renforcer avec \zh{得多} ou \zh{一点儿} après l'adjectif.
\end{corrige}

\begin{corrige}[Exercice 7 — Texte à trous : Connecteurs et mots-outils]
Texte complété :
\begin{quote}
喀麦隆自然资源很丰富，\textbf{所以}经济发展得快。\\
\textbf{因为}有石油，\textbf{而且}有很多农产品，\textbf{还}有木材。\\
\textbf{但是}基础设施还不够完善。\\
如果投资多，将来会\textbf{越来越}好。
\end{quote}
\emph{Kāmàilóng zìrán zīyuán hěn fēngfù, suǒyǐ jīngjì fāzhǎn de kuài. Yīnwèi yǒu shíyóu, érqiě yǒu hěn duō nóngchǎnpǐn, hái yǒu mùcái. Dànshì jīchǔ shèshī hái búgòu wánshàn. Rúguǒ tóuzī duō, jiānglái huì yuè lái yuè hǎo.}

Traduction : « Le Cameroun a des ressources naturelles très abondantes, donc son économie se développe vite. Parce qu'il a du pétrole, et aussi beaucoup de produits agricoles, et encore du bois. Mais les infrastructures ne sont pas encore assez abouties. Si les investissements sont nombreux, l'avenir ira de mieux en mieux. »

\textbf{Justifications :}
\begin{itemize}
    \item Trou 1 : \zh{所以} (« donc ») introduit la conséquence de la richesse des ressources.
    \item Trous 2 et 3 : \zh{因为} ouvre l'énumération des causes ; \zh{而且} ajoute un élément.
    \item Trou 4 : \zh{还} (« encore, de plus ») poursuit l'énumération (pétrole + produits agricoles + bois). Note : \zh{还} n'était pas dans la liste mais est la solution naturelle ; si l'on s'en tient strictement à la liste, on peut écrire \zh{而且有木材} au trou 3 et laisser la phrase en deux éléments.
    \item Trou 5 : \zh{但是} marque l'opposition (atouts / défis).
    \item Trou 6 : \zh{越来越} + adjectif exprime l'évolution progressive.
\end{itemize}
\end{corrige}

\begin{corrige}[Exercice 8 — Remise en ordre : Mots → Phrases]
\begin{enumerate}
    \item \zh{喀麦隆比泰国经济发展快。}\\
    \emph{Kāmàilóng bǐ Tàiguó jīngjì fāzhǎn kuài.}\\
    « Le Cameroun se développe économiquement plus vite que la Thaïlande. » (Ordre : sujet + \zh{比} + terme de comparaison + verbe.)
    \item \zh{因为喀麦隆有石油，所以经济好。}\\
    \emph{Yīnwèi Kāmàilóng yǒu shíyóu, suǒyǐ jīngjì hǎo.}\\
    « Parce que le Cameroun a du pétrole, son économie va bien. » (\zh{因为} en tête de la subordonnée, \zh{所以} en tête de la principale.)
    \item \zh{我觉得这个国家将来会越来越好。}\\
    \emph{Wǒ juéde zhège guójiā jiānglái huì yuè lái yuè hǎo.}\\
    « Je pense que ce pays ira de mieux en mieux à l'avenir. »
    \item \zh{喀麦隆的农产品很丰富。}\\
    \emph{Kāmàilóng de nóngchǎnpǐn hěn fēngfù.}\\
    « Les produits agricoles du Cameroun sont très abondants. » (Attention : \zh{农产品} est un seul mot, on ne peut pas séparer \zh{农} et \zh{产品} ; la phrase donnée exige de reconstituer le mot.)
    \item \zh{基础设施还是不够好。}\\
    \emph{Jīchǔ shèshī háishì búgòu hǎo.}\\
    « Les infrastructures ne sont toujours pas assez bonnes. » (\zh{还是} = « quand même, toujours » ; ordre : sujet + \zh{还是} + \zh{不够} + adjectif.)
\end{enumerate}
\textbf{Point de vigilance :} en chinois, l'ordre des mots est fixe : sujet -- adverbes (\zh{还}, \zh{都}) -- verbe -- compléments. Les adverbes ne se placent jamais après le verbe.
\end{corrige}

\begin{corrige}[Exercice 9 — Compréhension de l'écrit et langue (Type Bac)]
\textbf{Barème indicatif : 10 points} (compréhension : 5 pts ; langue : 5 pts).

\textbf{Questions de compréhension :}
\begin{enumerate}
    \item Trois ressources naturelles citées : le pétrole (\zh{石油}), le gaz naturel (\zh{天然气}), le cacao (\zh{可可}) — on accepte aussi le café (\zh{咖啡}) et le bois (\zh{木材}). (1 pt)
    \item Les deux piliers de l'économie thaïlandaise : l'industrie et le tourisme (\zh{工业和旅游业比较发达}). (1 pt)
    \item Les investissements étrangers sont moins nombreux parce que les infrastructures du Cameroun sont moins bonnes que celles de la Thaïlande. La structure de cause à effet utilisée est \zh{因为……所以……} : \zh{因为喀麦隆的基础设施不如泰国好，所以外资少一些。} (1 pt)
    \item Le connecteur d'opposition est \zh{但是} (« mais »). Le gouvernement investit pour construire des routes, des usines et des écoles : \zh{政府投资修路、建厂、办学}. (1 pt)
    \item Traduction de \zh{喀麦隆有潜力，将来一定能赶上来。} : « Le Cameroun a du potentiel, à l'avenir il pourra certainement rattraper son retard. » (1 pt)
\end{enumerate}
\textbf{Questions de langue :}
\begin{enumerate}
    \setcounter{enumi}{5}
    \item \zh{比} exprime la supériorité : A + \zh{比} + B + adjectif = « A est plus... que B ». Exemple : \zh{喀麦隆的石油比泰国多。} « Le Cameroun a plus de pétrole que la Thaïlande. » \zh{不如} exprime l'infériorité : A + \zh{不如} + B (+ adjectif) = « A n'est pas aussi... que B ». Exemple : \zh{喀麦隆的基础设施不如泰国。} « Les infrastructures du Cameroun ne valent pas celles de la Thaïlande. » (1,5 pt)
    \item \zh{越来越快} est la structure \zh{越来越} + adjectif, qui exprime un changement progressif : « de plus en plus vite ». (1 pt)
    \item Le mot « investir / investissement » est \zh{投资}, pinyin : tóuzī. (0,5 pt)
    \item Avec \zh{比} (le Cameroun comme sujet, il faut inverser la comparaison) : \zh{泰国的基础设施比喀麦隆好。} n'est pas accepté car le sujet serait la Thaïlande. Formulation correcte avec le Cameroun sujet : \zh{喀麦隆的基础设施比泰国差。} (Kāmàilóng de jīchǔ shèshī bǐ Tàiguó chà.) « Les infrastructures du Cameroun sont plus mauvaises que celles de la Thaïlande. » (2 pts)
\end{enumerate}
\textbf{Points de vigilance :} la question 9 est un piège classique : avec \zh{比}, pour garder le Cameroun comme sujet, il faut changer l'adjectif (\zh{好} $\rightarrow$ \zh{差}) ; \zh{不如} permet de garder l'adjectif positif. \zh{一定} exprime la certitude (« certainement »), \zh{赶上来} = « rattraper ».
\end{corrige}

\begin{corrige}[Exercice 10 — Thème et Version (Type Bac)]
\textbf{Barème indicatif : 5 points} (thème : 3 pts ; version : 2 pts).

\textbf{Partie A -- Thème (3 pts) :}
\zh{喀麦隆有很多农产品，所以经济发展得越来越快。}\\
\emph{Kāmàilóng yǒu hěn duō nóngchǎnpǐn, suǒyǐ jīngjì fāzhǎn de yuè lái yuè kuài.}

Détail du barème : \zh{农产品} correct (0,5) ; \zh{所以} pour « donc » (0,5) ; \zh{得} de complément de degré (0,5) ; \zh{越来越快} pour « de plus en plus vite » (1) ; ordre des mots et caractères corrects (0,5).

\textbf{Partie B -- Version (2 pts) :}
\zh{因为喀麦隆有很多石油和可可，所以经济越来越好。}\\
\emph{Yīnwèi Kāmàilóng yǒu hěn duō shíyóu hé kěkě, suǒyǐ jīngjì yuè lái yuè hǎo.}\\
Traduction : « Parce que le Cameroun possède beaucoup de pétrole et de cacao, son économie va de mieux en mieux. »

\textbf{Points de vigilance :} en français, on ne dit pas « parce que... donc... » dans la même phrase, mais en chinois \zh{因为} et \zh{所以} s'emploient ensemble : en version, il faut alléger (« Parce que..., son économie... » ou « Comme..., ... »). Ne pas confondre \zh{的} (possession) et \zh{得} (complément de degré) dans le thème.
\end{corrige}

\begin{corrige}[Exercice 11 — Expression écrite guidée : Mon pays et son économie (Type Bac)]
\textbf{Barème indicatif : 10 points} — respect des contraintes (6 phrases, comparaison avec \zh{比}, opinion avec \zh{觉得}/\zh{认为}, deux connecteurs, vocabulaire du thème) : 4 pts ; correction grammaticale (ordre des mots, mots-outils) : 3 pts ; exactitude des caractères : 2 pts ; cohérence et richesse du propos : 1 pt.

\textbf{Proposition de rédaction modèle (environ 95 caractères) :}
\begin{quote}
我的国家是喀麦隆。\\
喀麦隆的自然资源很丰富，有石油、可可和木材。\\
因为农产品很多，所以经济发展得比较快。\\
但是基础设施还不够好，外资比泰国少。\\
现在政府投资修路、建厂，经济越来越好。\\
我觉得喀麦隆的将来一定会更好。
\end{quote}
\emph{Wǒ de guójiā shì Kāmàilóng. Kāmàilóng de zìrán zīyuán hěn fēngfù, yǒu shíyóu, kěkě hé mùcái. Yīnwèi nóngchǎnpǐn hěn duō, suǒyǐ jīngjì fāzhǎn de bǐjiào kuài. Dànshì jīchǔ shèshī hái búgòu hǎo, wàizī bǐ Tàiguó shǎo. Xiànzài zhèngfǔ tóuzī xiūlù, jiànchǎng, jīngjì yuè lái yuè hǎo. Wǒ juéde Kāmàilóng de jiānglái yídìng huì gèng hǎo.}

Traduction : « Mon pays est le Cameroun. Les ressources naturelles du Cameroun sont très abondantes : il y a du pétrole, du cacao et du bois. Parce que les produits agricoles sont nombreux, l'économie se développe assez vite. Mais les infrastructures ne sont pas encore assez bonnes, et les capitaux étrangers sont moins nombreux qu'en Thaïlande. Maintenant, le gouvernement investit pour construire des routes et des usines, et l'économie va de mieux en mieux. Je pense que l'avenir du Cameroun sera certainement meilleur. »

\textbf{Vérification des contraintes :} 6 phrases complètes ; comparaison avec \zh{比} (\zh{外资比泰国少}) ; opinion avec \zh{我觉得} ; deux connecteurs (\zh{因为……所以……} et \zh{但是}) ; vocabulaire du thème (\zh{资源}、\zh{农产品}、\zh{发展}、\zh{基础设施}、\zh{投资}、\zh{将来}).

\textbf{Points de vigilance :}
\begin{itemize}
    \item Bien distinguer \zh{的} (\zh{我的国家}), \zh{得} (\zh{发展得比较快}) et \zh{地} (non requis ici).
    \item \zh{比} + adjectif sans « plus » : \zh{外资比泰国少} suffit.
    \item Ne pas traduire mot à mot « mon pays » par un possessif lourd : \zh{我的国家是喀麦隆} est correct et naturel.
    \item Compter ses caractères : rester entre 80 et 100, phrases courtes plutôt qu'une longue phrase risquée.
\end{itemize}
\end{corrige}

\newpage

\coursec{Fiche bilan}

\begin{popfigure}
\begin{tikzpicture}[mindmap, grow cyclic, every node/.style={concept, execute at begin node=\hspace{0pt}}, concept color=popblue, root concept/.append style={concept color=popblue, font=\small\bfseries}, level 1/.append style={level distance=3.4cm, sibling angle=51}, text=white]
\node[root concept]{Économie du Cameroun : expression écrite}
  child[concept color=poporange]{ node[align=center]{Lexique clé\\经济 农业 出口} }
  child[concept color=popgreen]{ node[align=center]{Comparatif\\A 比 B + adj.} }
  child[concept color=popgold]{ node[align=center]{Cause et but\\因为 所以 为了} }
  child[concept color=poppurple]{ node[align=center]{Connecteurs\\首先 然后 但是} }
  child[concept color=poppink]{ node[align=center]{Caractères\\clés et traits} }
  child[concept color=popgreen!70!black]{ node[align=center]{Thème et version\\traduire} }
  child[concept color=poporange!80!black]{ node[align=center]{Production\\texte simple} };
\end{tikzpicture}
\legende{Carte mentale de la leçon 34 : écrire sur l'économie du Cameroun.}
\end{popfigure}

\begin{bilan}
\textbf{1. Lexique clé (à connaître par cœur)}
\begin{itemize}
  \item \zh{经济} jīngjì : économie ; \zh{发展} fāzhǎn : développement / se développer ; \zh{发达} fādá : développé.
  \item \zh{农业} nóngyè : agriculture ; \zh{工业} gōngyè : industrie ; \zh{石油} shíyóu : pétrole ; \zh{咖啡} kāfēi : café ; \zh{可可} kěkě : cacao.
  \item \zh{出口} chūkǒu : exporter ; \zh{进口} jìnkǒu : importer ; \zh{市场} shìchǎng : marché ; \zh{工作} gōngzuò : travail.
  \item \zh{国家} guójiā : pays ; \zh{丰富} fēngfù : riche, abondant ; \zh{资源} zīyuán : ressources ; \zh{越来越} yuè lái yuè : de plus en plus.
\end{itemize}

\textbf{2. Grammaire à connaître par cœur}
\begin{center}
\begin{tabularx}{\linewidth}{|X|X|X|}
\hline
\rowcolor{popblueL} Structure & Exemple + pinyin & Traduction \\
\hline
A 比 B + adjectif & \zh{喀麦隆的农业比工业发达。} Kāmàilóng de nóngyè bǐ gōngyè fādá. & L'agriculture du Cameroun est plus développée que l'industrie. \\
\hline
因为 … 所以 … & \zh{因为资源很丰富，所以经济在发展。} Yīnwèi zīyuán hěn fēngfù, suǒyǐ jīngjì zài fāzhǎn. & Parce que les ressources sont riches, l'économie se développe. \\
\hline
越来越 + adjectif & \zh{喀麦隆的经济越来越发达。} Kāmàilóng de jīngjì yuè lái yuè fādá. & L'économie du Cameroun est de plus en plus développée. \\
\hline
为了 + but & \zh{为了发展经济，喀麦隆出口石油。} Wèile fāzhǎn jīngjì, Kāmàilóng chūkǒu shíyóu. & Pour développer son économie, le Cameroun exporte du pétrole. \\
\hline
\end{tabularx}
\end{center}

\textbf{3. Connecteurs pour écrire}
\begin{itemize}
  \item Enchaîner : \zh{首先} shǒuxiān (d'abord), \zh{然后} ránhòu (ensuite), \zh{最后} zuìhòu (enfin).
  \item Nuancer : \zh{但是} dànshì (mais), \zh{可是} kěshì (cependant).
  \item Conclure : \zh{所以} suǒyǐ (donc), \zh{总之} zǒngzhī (en résumé).
\end{itemize}

\textbf{4. Méthode express pour rédiger}
\begin{enumerate}
  \item Faire un plan en 3 parties : présentation, causes, conclusion.
  \item Utiliser une structure étudiée par phrase (比, 因为…所以…, 为了).
  \item Phrases courtes : Sujet + Verbe + Complément, jamais de phrase à rallonge.
  \item Relire : tons, ordre des mots, place du complément de temps avant le verbe.
\end{enumerate}

\textbf{5. Caractères du thème : pièges}
\begin{itemize}
  \item \zh{济} (jì) : clé de l'eau 氵 à gauche ; ne pas confondre avec \zh{挤} (jǐ, presser).
  \item \zh{发} (fā) : 5 traits, horizontal d'abord ; ne pas confondre avec \zh{友} (yǒu, ami).
  \item \zh{富} (fù) : clé du toit 宀 en haut ; ne pas oublier le trait horizontal sous le toit.
\end{itemize}
\end{bilan}

\begin{aretenir}[Checklist avant le Bac]
\begin{itemize}
  \item $\square$ Je connais le lexique de l'économie : 经济, 发展, 农业, 工业, 出口, 进口, 资源.
  \item $\square$ Je sais construire le comparatif : A 比 B + adjectif (jamais « plus » devant l'adjectif).
  \item $\square$ Je sais exprimer la cause avec 因为 … 所以 … (les deux ensemble, pas un seul).
  \item $\square$ Je sais exprimer le but avec 为了 + verbe.
  \item $\square$ J'utilise 越来越 + adjectif pour parler d'une évolution.
  \item $\square$ Je place le complément de temps avant le verbe : \zh{喀麦隆1960年独立了。}
  \item $\square$ Je sais écrire les caractères clés : 经, 济, 发, 展, 富, 源.
  \item $\square$ Je distingue 的 (nom), 地 (verbe), 得 (complément de degré).
  \item $\square$ Je sais traduire en version : je repère d'abord le sujet, puis le verbe, puis les compléments.
  \item $\square$ Je sais traduire en thème : je pense en chinois simple, pas mot à mot depuis le français.
  \item $\square$ Je structure mon texte avec 首先, 然后, 但是, 所以.
  \item $\square$ Je relis mon texte : ordre des mots, tons du pinyin, ponctuation chinoise 。，.
\end{itemize}
\end{aretenir}

\findecours

\end{document}
